% !TEX encoding = UTF-8 Unicode
\documentclass[UTF8,fontset=windows]{csoarticle}

% --- theorem environments (Chinese) ---
\newtheorem{theorem}{定理}
\newtheorem{lemma}{引理}
\renewcommand{\proofname}{证明}

% --- extra packages needed by the paper body ---
\usepackage{tabularx}
\usepackage{booktabs}
\usepackage{array}
\usepackage{ulem}
\usepackage{enumitem}

% --- framework annotation system ---
\newcommand{\mtag}[1]{\textbf{[#1]}}
\newcommand{\up}[1]{\ensuremath{^{#1}}}
\newcommand{\down}[1]{\ensuremath{_{#1}}}

% --- article header (China Science & Technology Online / paper.edu.cn) ---
\titleCHN{正反物质对风暴：一个统一四种基本力的涌现统计框架}
\titleENG{Matter--Antimatter Pair Storm: An Emergent Statistical Framework for the Unification of Fundamental Forces}
\authorCHN{杨韬\affil{1}}
\authorENG{taoyang\affil{1}}
\affiliationCHN{
    \affil{1} 作者单位，城市　邮编
}
\affiliationENG{
    \affil{1} Author Affiliation, City 000000
}
\abstractCHN{本文提出一个以\textbf{正反物质对风暴}为唯一底层机制的统一物理框架，旨在为四种基本力（引力、电磁、强核力、弱核力）的统一提供涌现式的统计力学论证。框架以四条公理为起点：（i）宇宙底层是正反物质对持续相遇---湮灭---再生的事件流，各向同性、50:50 平衡、普朗克尺度频率；（ii）湮灭存在极小概率的"泄漏"，其能标依赖形式统一解释重子不对称与暗能量；（iii）时间箭头分叉为两支互为 CPT 镜像的流形，时间 = 湮灭事件频率；（iv）空间可被卷曲成时间（黑洞时空互换的极限）。  在公理之上，我们建立以下论证链：  *时空生成\textbf{：事件流经因果集涌撒（sprinkling）与 Malament 恢复定理构造光滑洛伦兹流形，"流形 = 微变的统计平均"获得严格数学语言；离散网格的六方向结构张成 3+3 维时空（三空间 + 三时间），观测者生活在二维"时间表面"上（\S{}3.3）。  *引力的涌现}（\S{}4，核心）：质量遮蔽通量产生负压差，通量密度按引力势的\textbf{玻尔兹曼分布} \ensuremath{\rho}\ensuremath{\propto}e\up{-\Phi{}/c^{2}} 衰减（由遍历性 + 细致平衡 + 混沌传播导出，命题 2.1），给出 g\down{00} = \ensuremath{-}e\up{-2\Phi{}/c^{2}}，从而 \textbf{PPN 参数 \ensuremath{\gamma}=1、\ensuremath{\beta}=1 全部通过实验判决}；"微变权守恒"命题固定空间分量；经 Jacobson 热力学路径，Einstein 方程作为视界热力学恒等式涌现。黑洞熵 S = A/(4l\down{P}\up{2}) 获得完整论证链： \textbf{微态计数}（\S{}4.8）：通量管相干块（每 4 普朗克面积 1 比特，k = 2\up{1/4}），G、\ensuremath{\hbar} 作为普朗克面积的复合出现； \textbf{packing 命题包}（\S{}4.10）：1/4 系数的几何来源------\textbf{定理 1（k\ensuremath{\geq}4 下界，新证）} + Tammes 正四面体 + Voronoi 精确均分 + \textbf{命题 6（内空间编码优先性，框架内演绎）}------五处"4"同源交叉验证； \textbf{观测者模型}（\S{}4.9）：观测者 = 现在平面孤子（O1--O4），可访问信息 = 互信息 I(R;\ensuremath{\Sigma})，全息 = 观测者限制的几何必然；"现在平面选择"（N1--N4）由涡动力学（4\ensuremath{\pi} 相位锁定）+ 对称性自发破缺闭合；初始微扰 = \ensuremath{\varepsilon} 泄漏事件（P1--P4，"必然的偶然"）。 \textbf{v1.21 引力侧完整化}：(f) \textbf{各向同性指数度规识别}------玻尔兹曼 g\down{00}=\ensuremath{-}e\up{-2x} + 微变权守恒 g\down{rr}=e\up{2x} + 公理 1 各向同性 \ensuremath{\implies} g\down{\mu{}\nu{}}=diag(\ensuremath{-}e\up{-2x}, e\up{2x}, r\up{2}e\up{2x}, r\up{2}sin\up{2}\ensuremath{\theta}\ensuremath{\cdot}e\up{2x})（即文献 exponential metric），\textbf{完整 PPN 参数集 \ensuremath{\gamma}=\ensuremath{\beta}=1 其余全 0（1PN 与 GR 不可区分）}，2PN g\down{rr} 分叉 33\%、3PN g\down{00} 分叉 11\%（B3）；(g) \textbf{阴影与 EHT 相容}------各向同性约定下框架阴影 b\down{min}=2e\ensuremath{\cdot}GM/c\up{2}=5.436 GM/c\up{2} 与 EHT 5.2\ensuremath{\pm}0.3 相容（+4.55\%），修正 v1.14"强场与 EHT 冲突"（那是标准坐标约定的产物）------框架强场从"无度规承诺"升级为"有与观测相容的几何预言"；(h) \textbf{方向离散化}------普朗克管径\mtag{本框架设定：方向离散化}\ensuremath{\implies} 方向簇 k=4（王虹/Tammes）+ 正四面体凸覆盖 + 统计平均涌现连续方向 \ensuremath{\implies} \textbf{Bekenstein 上界 S\ensuremath{\leq}A/(4l\down{P}\up{2}) 从方向离散化恢复}，与 k=2\up{1/4} 微态计数（"="方向）双闭合（B4）；(i) \textbf{观测者互信息面积律}------体=表面投影 \ensuremath{\implies} H(R|\ensuremath{\Sigma})=0 \ensuremath{\implies} I(R;\ensuremath{\Sigma})=H(\ensuremath{\Sigma}) \ensuremath{\propto} A；观测者原始可读 4 bits/面元经 16 倍相干压缩（4面元\ensuremath{\times}4方向）\ensuremath{\implies} 熵密度 1/4 bit/面元 = 1/(4l\down{P}\up{2})（B5）。  *规范群结构\textbf{：标准模型规范群 SU(3)\ensuremath{\times}SU(2)\ensuremath{\times}U(1) 对应"配对阶数" n=3,2,1 的表示论结构（}v2.5：三准则结构筛选唯一化\textbf{------d\down{min} 匹配/同号配对阶数/三阶反对称，10 候选逐一排除，\S{}5.2），色禁闭获得三阶反对称张量的数学根源；三代费米子对应涡的三种卷绕模；多表示统一场给出四力 = 通量场涨落的分波分解。}v1.19 新增规范侧涌现论证\textbf{（\S{}5.6--5.8）：(a) }量子双代数命题包\textbf{------框架事件流三操作 \ensuremath{\leftrightarrow} Hopf 结构，非阿贝尔规范结构可从离散代数精确涌现（Kitaev 定理，Z\down{2}/Z\down{3}/S\down{3} 数值验证）；(b) }框架哈密顿量\textbf{ H\down{F} = \ensuremath{-}J\ensuremath{\Sigma}A\down{v} \ensuremath{-} K\ensuremath{\Sigma}B\down{p}------[H\down{F},A\down{v}]=[H\down{F},B\down{p}]=0 证明规范对称性是动力学对称性（涌现），2+1D 正能隙（拓扑序）为质量隙低维严格实现，环面简并 = |G|\up{2}；(c) }Wen 弦网 3+1D 路线\textbf{------量子直线 = 弦网边（本框架设定：宇宙弦=动态涌现网络），四力 = 弦网序涌现通道（电磁\ensuremath{\to}U(1) 弦网光子、引力\ensuremath{\to}弦网涌现自旋 2）；}v1.28--1.30 严格化推进\textbf{：\ensuremath{\lambda}=0 弦网凝聚严格证明（E0=\ensuremath{-}N\down{c}、简并 2\up{n-N\down{c}}、\ensuremath{\Delta}=2）+ 对偶映射严格 + }Z\down{2} \ensuremath{\subset} U(1) 群嵌入定理 + 麦克斯韦线性化命题 + 库仑相光子色散数值验证\textbf{（连续极限重整化群 = 世界级开放）+ }自旋 2 = 二阶张量旋转本征值 \ensuremath{\pm}2 严格验证\textbf{（本框架自旋 2 机制）；}v1.34/1.36 深化\textbf{：}弦网低能谱判据(4)数值化\textbf{（C4 螺旋度分解，自旋 2 模式有隙 \ensuremath{\Delta}\ensuremath{\pm}2\ensuremath{\approx}0.49、二阶张量扇区承载）+ }3+1D \ensuremath{\lambda}\down{c}\ensuremath{\approx}0.14\textbf{（四尺寸锁定，与 2D toric code 构成维数-临界耦合递减谱系）+ }Wilson RG 路径\textbf{（U(1) 高斯固定点 + 2D Ising 块 RG 流向连续极限------定理 A + 命题 B + 数值验证 C 四层闭合）------弦网路径从半严格升级为部分严格，与引力统计力学涌现（\S{}4）双路径印证。}v1.20 补充禁闭弦张力数值实现\textbf{（\S{}5.7）：(d) }X 场版电荷对\textbf{ \ensuremath{\sigma}\down{E}\ensuremath{\approx}2\ensuremath{\kappa}------三尺寸（3\ensuremath{\times}3/3\ensuremath{\times}4/4\ensuremath{\times}4）热力学极限收敛 <0.02\%，\ensuremath{\kappa}=0 严格 deconfined；(e) }Z 场版磁通对\textbf{ \ensuremath{\sigma}\down{m}\ensuremath{\approx}2\ensuremath{\lambda}（约束扇区方法：B\down{p} 扇区不变子空间）------0.21/1.09/3.04 @ \ensuremath{\lambda}=0.5/1.0/2.0（3\ensuremath{\times}4），\ensuremath{\lambda}=0 严格 deconfined------框架"禁闭 = 通量管线性势"\mtag{本框架机制：禁闭=通量管线性势}获得定量数值实现（早期数值经 P0-1 复核替换）；(f) }相变点定位（P0-3 开边界 FSS \ensuremath{\to} v1.23 环面精确化）\textbf{------序参量 m=\ensuremath{\langle}X\ensuremath{\rangle} 交叉点：开边界 FSS \ensuremath{\kappa}\down{c}\ensuremath{\approx}0.25\ensuremath{\pm}0.05（v1.22）；}环面（周期边界）精确化半极化点 \ensuremath{\kappa}\*\ensuremath{\approx}0.20\ensuremath{\pm}0.01\textbf{（2\ensuremath{\times}2/2\ensuremath{\times}3/2\ensuremath{\times}4/3\ensuremath{\times}3 尺寸无关，消除边界污染）；}判据修正\textbf{------m=0.5 交叉为半极化点非相变点（环面实测 m(h=0.32)\ensuremath{\approx}0.84），真相变 }\ensuremath{\kappa}\down{c}\ensuremath{\approx}0.25--0.32，与文献 X 场（\ensuremath{\sigma}\up{x} 场）版本 h\down{xc}\ensuremath{\approx}0.32 一致\textbf{（注：h\down{c}=0.328 为 \ensuremath{\sigma}\up{z} 场版本值，P0-3 曾混用，v1.23 已修正）------拓扑序\ensuremath{\to}禁闭转变确认，禁闭相图完整（\S{}5.7）。  *量子力学基础}：经典事件流无法产生贝尔违背，框架必须非交换化------事件流升级为算子代数，CHSH = 2\ensuremath{\sqrt{2}}、互补性 V\up{2}+I\up{2}\ensuremath{\leq}1 获得证明；定位是 't Hooft 式的：为量子力学的希尔伯特结构提供"事件流不交换"的本体论动机。  *宇宙学\textbf{：泄漏概率 \ensuremath{\varepsilon}(E) 的能标依赖（\ensuremath{\beta}\ensuremath{\approx}0.615）统一解释重子不对称（\ensuremath{\eta}\ensuremath{\approx}6\ensuremath{\times}10\up{-10}）与暗能量密度（\ensuremath{\rho}\down{\Lambda{}}\ensuremath{\sim}(meV)\up{4}）；统计衰减因子 \ensuremath{\varepsilon}\down{B}=(l\down{P}/R\down{U})\up{2} 给出暗物质密度；暗物质 = 本支普朗克涨落残余 + 镜像支引力投影的叠加模型；}初始微扰（\ensuremath{\varepsilon}\down{0}=7.7\ensuremath{\times}10\up{-20}）与重子不对称/暗能量三源同根\textbf{（\S{}4.9 P4）。  框架给出多项可检验预言，包括：}黑洞强场预言包（B1--B5）\textbf{（极端 Kerr 霍金辐射角分布偏离轴对称、原初黑洞熵体积修正 \ensuremath{\sim}10\up{-8}、霍金辐射-落入物质纠缠、ER 桥 packing 上限、非热修正 1/S）、}EMRI 进动差判决\textbf{（S-E1--S-E4：}v1.35 修正后 \ensuremath{\Delta}lnL 判据 27\ensuremath{\sigma}/28\ensuremath{\sigma}（a-\ensuremath{\varphi}0 全自由稳健，\ensuremath{\sigma}\down{\delta{}}=1.6\ensuremath{\times}10\up{-6}）\textbf{）、}FQH 引力子模 \ensuremath{\varepsilon} 项\textbf{（S-E5，本框架预言）、吸积盘火墙非热修正（S4，原初黑洞窗口）、空洞-镜像支引力投影关联（void-CPT-\ensuremath{\kappa} 异常）、暗能量 w=\ensuremath{-}1+4\ensuremath{\varepsilon}/3 的常数偏离、SU(3) \ensuremath{\beta} 函数阶梯跳跃、纠缠熵面积律的 ln 2 因子。  *诚实声明}：本框架是启发式统一框架，非已确立物理。文中所有\mtag{推导}为框架内自洽论证，所有\mtag{待验证}项不得作为推导依据。\textbf{v1.0 的强场度规外推预言（g\down{00}=\ensuremath{-}e\up{-2x} 中子星偏离）已正式撤回}（超出适用域，\S{}4.6）；框架强场内容为七环机制链 + 各向同性指数度规（阴影与 EHT 相容，\S{}4.5b/\S{}4.6）。剩余主要开放问题：非阿贝尔分量的完全涌现（4D）、J/K 耦合具体值、3+1D 弦网涌现 U(1) 连续极限与自旋 2（\S{}10.2）。  --}
\abstractENG{This paper proposes a unified physical framework based on a single underlying mechanism, the \textbf{matter--antimatter pair storm}, providing an emergent statistical-mechanical account of the unification of the four fundamental forces (gravity, electromagnetism, strong and weak nuclear forces). The framework begins with four axioms: (i) the universe at its deepest level is an event flow of continuously meeting--annihilating--regenerating matter--antimatter pairs, isotropic, 50:50 balanced, at Planck-scale frequency; (ii) annihilation has an extremely small probability of "leakage," whose energy-scale-dependent form unifies the baryon asymmetry and dark energy; (iii) the arrow of time bifurcates into two manifolds that are mutual CPT mirrors, with time = frequency of annihilation events; (iv) space can be curled into time (the limit of black-hole space--time interchange).  On top of the axioms, we establish the following chain of arguments:  *Emergence of space--time\textbf{: The event flow constructs a smooth Lorentzian manifold via causal-set sprinkling and the Malament recovery theorem, giving rigorous mathematical language to "manifold = statistical average of micro-variations"; the six-direction structure of the discrete grid spans 3+3 dimensional space--time (three spatial + three temporal), and observers live on a two-dimensional "time surface" (\S{}3.3).  *Emergence of gravity} (\S{}4, core): mass-shielded flux produces a negative pressure difference, and the flux density decays according to the \textbf{Boltzmann distribution} \ensuremath{\rho}\ensuremath{\propto}e\up{-\Phi{}/c^{2}} of the gravitational potential (derived from ergodicity + detailed balance + chaotic propagation, Proposition 2.1), giving g\down{00} = \ensuremath{-}e\up{-2\Phi{}/c^{2}} and hence \textbf{PPN parameters \ensuremath{\gamma}=1, \ensuremath{\beta}=1, all passing experimental verdicts}; the "micro-variation weight conservation" proposition fixes the spatial components; via the Jacobson thermodynamic path, the Einstein equation emerges as a horizon-thermodynamics identity. Black-hole entropy S = A/(4l\down{P}\up{2}) acquires a complete argument chain: \textbf{Microstate counting} (\S{}4.8): flux-tube coherent blocks (1 bit per 4 Planck areas, k = 2\up{1/4}), with G, \ensuremath{\hbar} appearing as composites of the Planck area; \textbf{Packing proposition package} (\S{}4.10): geometric origin of the 1/4 coefficient --- \textbf{Theorem 1 (k\ensuremath{\geq}4 lower bound, new proof)} + Tammes regular tetrahedron + exact Voronoi partition + \textbf{Proposition 6 (interior-space encoding priority, framework-internal deduction)} --- five occurrences of "4" cross-validated as one source; \textbf{Observer model} (\S{}4.9): observer = present-plane soliton (O1--O4), accessible information = mutual information I(R;\ensuremath{\Sigma}), holography = geometric necessity of observer limitation; "present-plane selection" (N1--N4) closed by vortex dynamics (4\ensuremath{\pi} phase locking) + spontaneous symmetry breaking; initial perturbation = \ensuremath{\varepsilon} leakage event (P1--P4, "necessary contingency"). \textbf{v1.21 completion of the gravity side}: (f) \textbf{isotropic exponential metric identification} --- Boltzmann g\down{00}=\ensuremath{-}e\up{-2x} + micro-variation weight conservation g\down{rr}=e\up{2x} + Axiom 1 isotropy \ensuremath{\implies} g\down{\mu{}\nu{}}=diag(\ensuremath{-}e\up{-2x}, e\up{2x}, r\up{2}e\up{2x}, r\up{2}sin\up{2}\ensuremath{\theta}\ensuremath{\cdot}e\up{2x}) (the literature exponential metric), \textbf{complete PPN parameter set \ensuremath{\gamma}=\ensuremath{\beta}=1, all others 0 (indistinguishable from GR at 1PN)}, 2PN g\down{rr} bifurcation 33\%, 3PN g\down{00} bifurcation 11\% (B3); (g) \textbf{shadow compatible with EHT} --- in the isotropic convention the framework shadow b\down{min}=2e\ensuremath{\cdot}GM/c\up{2}=5.436 GM/c\up{2} is compatible with EHT 5.2\ensuremath{\pm}0.3 (+4.55\%), correcting v1.14's "strong-field conflict with EHT" (an artifact of the standard coordinate convention) --- the framework's strong field is upgraded from "no metric commitment" to "geometric predictions compatible with observation"; (h) \textbf{direction discretization} --- Planck-tube diameter \mtag{Framework assumption: direction discretization}\ensuremath{\implies} direction cluster k=4 (Wang Hong/Tammes) + regular-tetrahedron convex covering + statistical average yielding continuous directions \ensuremath{\implies} \textbf{Bekenstein bound S\ensuremath{\leq}A/(4l\down{P}\up{2}) recovered from direction discretization}, doubly closed with the k=2\up{1/4} microstate counting ("=" directions) (B4); (i) \textbf{observer mutual-information area law} --- bulk = surface projection \ensuremath{\implies} H(R|\ensuremath{\Sigma})=0 \ensuremath{\implies} I(R;\ensuremath{\Sigma})=H(\ensuremath{\Sigma}) \ensuremath{\propto} A; the observer's raw 4 bits/plaquette compressed coherently by a factor 16 (4 plaquettes \ensuremath{\times} 4 directions) \ensuremath{\implies} entropy density 1/4 bit/plaquette = 1/(4l\down{P}\up{2}) (B5).  *Gauge-group structure\textbf{: The Standard Model gauge group SU(3)\ensuremath{\times}SU(2)\ensuremath{\times}U(1) corresponds to the representation-theoretic structure of "pairing orders" n=3,2,1 (}v2.5: uniqueness via three-criterion structural screening\textbf{ --- d\down{min} matching / same-sign pairing order / third-order antisymmetric tensor, 10 candidates eliminated one by one, \S{}5.2); color confinement acquires its mathematical root in the third-order antisymmetric tensor; the three fermion generations correspond to three winding modes of the vortex; the multi-representation unified field gives the four forces as a partial-wave decomposition of flux-field fluctuations. }v1.19 gauge-side emergence arguments\textbf{ (\S{}5.6--5.8): (a) }quantum-double algebra proposition package\textbf{ --- the framework's three event-flow operations \ensuremath{\leftrightarrow} Hopf structure, non-Abelian gauge structure emerging exactly from discrete algebra (Kitaev theorem, Z\down{2}/Z\down{3}/S\down{3} numerical verification); (b) }framework Hamiltonian\textbf{ H\down{F} = \ensuremath{-}J\ensuremath{\Sigma}A\down{v} \ensuremath{-} K\ensuremath{\Sigma}B\down{p} --- [H\down{F},A\down{v}]=[H\down{F},B\down{p}]=0 proves gauge symmetry is a dynamical symmetry (emergent), 2+1D positive energy gap (topological order) is a low-dimensional rigorous realization of the mass gap, torus degeneracy = |G|\up{2}; (c) }Wen string-net 3+1D route\textbf{ --- quantum lines = string-net edges (Framework assumption: cosmic strings = dynamically emergent network), four forces = emergent channels of string-net order (EM\ensuremath{\to}U(1) string-net photon, gravity\ensuremath{\to}string-net emergent spin-2); }v1.28--1.30 rigorization\textbf{: rigorous proof of \ensuremath{\lambda}=0 string-net condensation (E0=\ensuremath{-}N\down{c}, degeneracy 2\up{n-N\down{c}}, \ensuremath{\Delta}=2) + rigorous duality mapping + }Z\down{2} \ensuremath{\subset} U(1) group-embedding theorem + Maxwell linearization proposition + numerical verification of Coulomb-phase photon dispersion\textbf{ (continuum-limit renormalization group = world-class open problem) + }spin 2 = rigorous verification of second-order tensor rotation eigenvalues \ensuremath{\pm}2\textbf{ (framework spin-2 mechanism); }v1.34/1.36 deepening\textbf{: }criterion (4) of the string-net low-energy spectrum numericalized\textbf{ (C4 helicity decomposition, spin-2 mode gapped \ensuremath{\Delta}\ensuremath{\pm}2\ensuremath{\approx}0.49, carried by second-order tensor sector) + }3+1D \ensuremath{\lambda}\down{c}\ensuremath{\approx}0.14\textbf{ (locked by four system sizes, forming a dimension--critical-coupling decreasing sequence with 2D toric code) + }Wilson RG path\textbf{ (U(1) Gaussian fixed point + 2D Ising block RG flowing to the continuum limit --- Theorem A + Proposition B + Numerical verification C, four-layer closure) --- the string-net path upgraded from semi-rigorous to partially rigorous, cross-validated with the gravitational statistical-mechanics emergence (\S{}4) as two independent paths. }v1.20 confined string-tension numerics\textbf{ (\S{}5.7): (d) }X-field charge pair\textbf{ \ensuremath{\sigma}\down{E}\ensuremath{\approx}2\ensuremath{\kappa} --- three system sizes (3\ensuremath{\times}3/3\ensuremath{\times}4/4\ensuremath{\times}4) thermodynamic-limit convergence <0.02\%, \ensuremath{\kappa}=0 strictly deconfined; (e) }Z-field flux pair\textbf{ \ensuremath{\sigma}\down{m}\ensuremath{\approx}2\ensuremath{\lambda} (constrained-sector method: B\down{p}-sector invariant subspace) --- 0.21/1.09/3.04 at \ensuremath{\lambda}=0.5/1.0/2.0 (3\ensuremath{\times}4), \ensuremath{\lambda}=0 strictly deconfined --- the framework's "confinement = flux-tube linear potential" \mtag{Framework mechanism: confinement = flux-tube linear potential} acquires quantitative numerical realization (early numerics superseded after P0-1 review); (f) }transition-point localization (P0-3 open-boundary FSS \ensuremath{\to} v1.23 torus refinement)\textbf{ --- order parameter m=\ensuremath{\langle}X\ensuremath{\rangle} crossing: open-boundary FSS \ensuremath{\kappa}\down{c}\ensuremath{\approx}0.25\ensuremath{\pm}0.05 (v1.22); }torus (periodic boundary) refinement: half-polarization point \ensuremath{\kappa}*\ensuremath{\approx}0.20\ensuremath{\pm}0.01\textbf{ (2\ensuremath{\times}2/2\ensuremath{\times}3/2\ensuremath{\times}4/3\ensuremath{\times}3 size-independent, boundary pollution removed); }criterion correction\textbf{ --- the m=0.5 crossing is a half-polarization point, not the transition point (torus measurement m(h=0.32)\ensuremath{\approx}0.84), true transition }\ensuremath{\kappa}\down{c}\ensuremath{\approx}0.25--0.32, consistent with the literature X-field (\ensuremath{\sigma}\up{x} field) value h\down{xc}\ensuremath{\approx}0.32\textbf{ (note: h\down{c}=0.328 is the \ensuremath{\sigma}\up{z-field} value, mixed up in P0-3, corrected in v1.23) --- topological-order\ensuremath{\to}confinement transition confirmed, confinement phase diagram complete (\S{}5.7).  *Foundations of quantum mechanics}: The classical event flow cannot produce Bell violation; the framework must be non-commutativized --- the event flow is upgraded to an operator algebra, with CHSH = 2\ensuremath{\sqrt{2}} and complementarity V\up{2}+I\up{2}\ensuremath{\leq}1 proven; the positioning is 't Hooft-like: providing an ontological motivation ("event flow does not commute") for the Hilbert-space structure of quantum mechanics.  *Cosmology\textbf{: The energy-scale dependence of the leakage probability \ensuremath{\varepsilon}(E) (\ensuremath{\beta}\ensuremath{\approx}0.615) unifies the baryon asymmetry (\ensuremath{\eta}\ensuremath{\approx}6\ensuremath{\times}10\up{-10}) and the dark-energy density (\ensuremath{\rho}\down{\Lambda{}}\ensuremath{\sim}(meV)\up{4}); the statistical decay factor \ensuremath{\varepsilon}\down{B}=(l\down{P}/R\down{U})\up{2} yields the dark-matter density; dark matter = superposition model of this-branch Planck-fluctuation remnants + mirror-branch gravitational projection; }the initial perturbation (\ensuremath{\varepsilon}\down{0}=7.7\ensuremath{\times}10\up{-20}) shares one root with baryon asymmetry and dark energy --- three sources, one origin\textbf{ (\S{}4.9 P4).  The framework yields a number of testable predictions, including: }black-hole strong-field prediction package (B1--B5)\textbf{ (extreme-Kerr Hawking-radiation angular distribution deviating from axisymmetry, primordial black-hole entropy volume correction \ensuremath{\sim}10\up{-8}, Hawking-radiation--infalling-matter entanglement, ER-bridge packing bound, non-thermal correction 1/S), }EMRI precession-difference verdict\textbf{ (S-E1--S-E4: }after v1.35 correction, \ensuremath{\Delta}lnL criterion 27\ensuremath{\sigma}/28\ensuremath{\sigma} (robust with a-\ensuremath{\varphi}0 fully free, \ensuremath{\sigma}\down{\delta{}}=1.6\ensuremath{\times}10\up{-6})\textbf{), }FQH graviton-mode \ensuremath{\varepsilon} term\textbf{ (S-E5, framework prediction), accretion-disk firewall non-thermal correction (S4, primordial black-hole window), void--mirror-branch gravitational-projection correlation (void-CPT-\ensuremath{\kappa} anomaly), constant deviation of dark energy w=\ensuremath{-}1+4\ensuremath{\varepsilon}/3, SU(3) \ensuremath{\beta}-function staircase jump, ln 2 factor in the entanglement-entropy area law.  *Honesty statement}: This framework is a heuristic unification framework, not established physics. All \mtag{Derivation}s in this paper are internally consistent arguments within the framework, and no \mtag{To be verified} item may be used as a basis for derivation. \textbf{The v1.0 strong-field metric extrapolation prediction (g\down{00}=\ensuremath{-}e\up{-2x} neutron-star deviation) has been formally withdrawn} (beyond the domain of validity, \S{}4.6); the framework's strong-field content is the seven-loop mechanism chain + isotropic exponential metric (shadow compatible with EHT, \S{}4.5b/\S{}4.6). Remaining major open problems: complete emergence of the non-Abelian component (4D), concrete values of J/K coupling, string-net emergent U(1) continuum limit and spin 2 (3+1D) (\S{}10.2).  --}
\keywordCHN{正反物质对风暴；涌现统计框架；四种基本力统一；量子引力；引力涌现；弦网}
\keywordENG{matter--antimatter pair storm; emergent statistical framework; unification of fundamental forces; quantum gravity; emergent gravity; string-net}
\cateidCHN{O412}
\authorIntroduction{作者姓名（出生年-），性别，职称，主要研究方向：***。}
\fund{基金项目（编号）}

\begin{document}
\maketitle



\begin{quote}
\textbf{论文初稿 v2.8}（M7 因子推导补强）\ensuremath{\cdot} 2026-08-24 \ensuremath{\cdot} 启发式统一框架（参考架构，非已确立物理）
\textbf{标注约定}：\mtag{公理}= 基本假设；\mtag{假设}= 本框架引入的物理假设；\mtag{推导}= 本框架内完成的推导；\mtag{引用}= 已确立的物理/数学结果；\mtag{待验证}= 需进一步验证的猜想。\textbf{术语强度规范（v2.2）}：\textbf{定理} = 数学上严格、无框架依赖（含文献已确立）；\textbf{命题} = 在框架公理下成立的演绎；\textbf{推论} = 命题的直接推论；\textbf{推测} = 未证明的猜想；\textbf{论证} = 启发式/结构论证；\mtag{数值}= 数值验证结果。
\end{quote}


\section{引言}

\subsection{动机与背景}

现代物理学的两大支柱------广义相对论（GR）与量子场论（QFT）------分别精确描述了引力与其它三种基本力，但两者的统一（量子引力）是物理学最深层的未解问题。标准模型以 SU(3)\ensuremath{\times}SU(2)\ensuremath{\times}U(1) 规范群统一描述强、弱、电磁力，GR 以黎曼几何描述引力，但：

\begin{itemize}
\item 引力不可重整化，且缺乏量子完备描述；
\item 标准模型含 19--26 个自由参数（层级问题、宇宙学常数问题、强 CP 问题）；
\item 四种力的相对强度（1 : 10\up{-2} : 10\up{-6} : 10\up{-38}）缺乏统一解释；
\item 大统一（GUT）、弦论、圈量子引力等候选路径均未获得实验支持。
\end{itemize}

本框架提出一种不同的统一策略：\textbf{不假设更大的对称群或额外的维度，而是寻找一个底层统计机制（事件流），使四种力成为同一机制在不同尺度/构型下的涌现模态}。这一策略与统计力学之于热力学的关系类似：GR 是几何描述，本框架的统计流体描述是其有效理论。

\subsection{核心概念：正反物质对风暴}

\mtag{公理 1}宇宙底层是正反物质对持续相遇---湮灭---再生的事件流，各向同性、50:50 平衡、普朗克尺度频率（\ensuremath{\nu} \ensuremath{\sim} t\down{P}\up{-1} \ensuremath{\approx} 10\up{43} s\up{-1}）。每一次相遇是一次"微变"；时空流形 = 微变的统计平均（"流形无限微变"）。

四种基本力 = 同一风暴的四种模态：

\begin{table}[htbp]
\centering
\small
\begin{tabularx}{\textwidth}{@{}lXX@{}}
\toprule
力 & 模态 & 机制 \\
\midrule
强力 & 屏蔽港湾 & 近距离粒子共享低压腔（色禁闭的直观） \\
弱力 & 穿堂风 & 极短程的管道耦合（手性、W/Z 质量） \\
电磁 & 选择性耦合 & 长程通量压力梯度（空腔模型） \\
引力 & 通量负压 & 长程极弱的通量遮挡亏损 \\
\bottomrule
\end{tabularx}
\end{table}

\textbf{统一 = 模态分界距离的消失}：电弱统一（100 GeV）、GUT（10\up{15} GeV）、普朗克（10\up{19} GeV）三个能标是模态边界逐一消失的阶梯；\textbf{破缺 = 通量密度下降后模态边界浮现}。

\subsection{论文结构}

\S{}2 建立四条公理与基本对象；\S{}3 论证时空的涌现生成（因果集、时间表面、3+3 维）；\S{}4 是核心章节------引力的完整涌现论证（玻尔兹曼分布、PPN 判决、微态计数、热平衡命题、\textbf{七环强场机制、packing 命题包、观测者模型 O1--O4、现在平面选择 N1--N4、初始微扰来源 P1--P4}）；\S{}5 规范群结构；\S{}6 动力学（最小作用量）；\S{}7 量子力学基础；\S{}8 宇宙学；\S{}9 可检验预言；\S{}10 结论与开放问题。



\section{公理与基本设定}

\subsection{四条公理}

\textbf{\mtag{公理 1（正反物质对风暴）}}：宇宙底层是正反物质对持续相遇---湮灭---再生事件流，各向同性、50:50 平衡、普朗克尺度频率；流形 = 微变的统计平均。

\textbf{\mtag{公理 2（泄漏机制）}}：每次湮灭有极小概率 \ensuremath{\varepsilon} 绑定失败，正份漏入正时间箭头支、反份漏入镜像支。\ensuremath{\varepsilon} 能标依赖：\ensuremath{\varepsilon}(E) = \ensuremath{\varepsilon}\down{0}(M\down{Pl}/E)\up{\beta{}，\beta{}} \ensuremath{\approx} 0.615。

\textbf{\mtag{公理 3（时间箭头分叉）}}：两支互为 CPT 镜像；时间 = 相遇/湮灭频率，由交叠积分 \ensuremath{\tau}(x) = \ensuremath{\int}\ensuremath{\bar{\psi}}\ensuremath{\psi} d\up{4}x 的生成率定义。

\textbf{\mtag{公理 4（空间卷曲成时间）}}：空间可被卷曲成时间；黑洞内部时空互换 = 空间卷曲到极限（径向被卷曲成帧推进方向）。

\subsection{基本对象与定义}

\begin{itemize}
\item \textbf{微变权流 J\up{\mu{}}}\mtag{推导}：湮灭事件流的守恒流（\ensuremath{\nabla}\down{\mu{}} J\up{\mu{}} = 0，\ensuremath{\varepsilon} \ensuremath{\ll} 1 时）。微变权 W[\ensuremath{\Sigma}] = \ensuremath{\int}\down{\Sigma{}} J\up{\mu{}} d\ensuremath{\Sigma}\down{\mu{}}。
\item \textbf{通量密度 \ensuremath{\rho}\down{pair}(x)}：局部湮灭事件率密度，决定时间流速。
\item \textbf{量子直线}：同源泄漏事件在两支间的关联通道（ER=EPR 的底层本体）。
\item \textbf{诸元}：构成流形的基本单元（六方向连接，张成三平面）。
\end{itemize}

\subsection{力学的层级结构}

\begin{verbatim}
事件流（普朗克噪声）
  |-- 统计平均 => 流形/时空
  |-- 孤子化 => 粒子（稳定局域结构）
  |-- 压强抵抗 => 质量/惯性（等效原理 = 同耦合的两侧面）
  |-- 通量响应 => 四力模态 => 能标上升 => 统一/破缺
  `-- 泄漏 eps => 时间箭头 / 重子不对称 / 暗能量
\end{verbatim}



\section{时空生成}

\subsection{因果集公理化}

\mtag{引用 + 推导}把公理 1 的事件流形式化为\textbf{因果集}（Causal Set）：事件集 C 配偏序 \ensuremath{\prec}，满足不自反、传递、局部有限。每次相遇---湮灭 \ensuremath{\in} C，"彼此可影响" = \ensuremath{\prec}。各向同性对应\textbf{泊松涌撒}（sprinkling）：在光滑洛伦兹流形 (M,g) 上按泊松过程以密度 \ensuremath{\rho} = l\down{P}\up{-4} 随机落点。

\textbf{Malament 恢复定理}\mtag{引用}：若 (M,g)、(M',g') 均为过去/将来可区分流形且因果序同构，则 M \ensuremath{\cong} M' 且 g' = \ensuremath{\Omega}\up{2}g------因果结构唯一决定度规的共形类。共形因子由事件计数固定：

\[g_{\mu\nu}(x) = \Omega^2(x)g_0{}_{\mu\nu}(x),\qquad \Omega^4(x) = \frac{dN(x)}{\rho\,dV_0(x)}\]

\textbf{"流形微变"三段命题链}\mtag{推导}：
\begin{enumerate}
\item 泊松涌撒（大数定律）\ensuremath{\to} 因果偏序；
\item Malament 恢复定理 \ensuremath{\to} 度规共形类；
\item 计数密度 \ensuremath{\to} 共形因子 \ensuremath{\Omega}。
\end{enumerate}

\subsection{时间 = 事件频率}

\mtag{推导}静态度规 ds\up{2} = g\down{00} dt\up{2} + h\down{ij} dx\up{i} dx\up{j} 下，理想钟 = 事件计数器。事件在固有时上以普朗克频率 f\down{0} 均匀发生，坐标时事件率 f(x)：

\[g_{00}(x) = -\left(\frac{f(x)}{f_0}\right)^2,\qquad \frac{f_B}{f_A} = \sqrt{\frac{g_{00}(B)}{g_{00}(A)}} \approx 1 + \frac{\Phi_B - \Phi_A}{c^2}\]

弱场给出引力红移------\textbf{频率梯度与 g\down{00} 梯度一一对应}。

\subsection{3+3 维离散时空与时间表面}

\mtag{推导}诸元六方向（前后/左右/上下）张成三坐标轴、三平面（xy/xz/yz），分别对应\textbf{过去/现在/未来三个时间维度}。我们只感知一维时间，因另外两维被卷曲进空间平面法线方向。

\textbf{推论（全息 = 时间表面的几何必然）}\mtag{推导，完整命题链见 \S{}4.9}：
\begin{itemize}
\item 观测者 = 现在平面上的局域孤子（O1），可访问信息 = 时间表面上的投影 = 互信息 I(R;\ensuremath{\Sigma})（O2）；
\item 体信息冗余，表面信息基本；信息容量 \ensuremath{\propto} 表面积（O4）；
\item 观测者的均匀先验由观测者几何导出（O3）------外空间方向无信息通道。
\end{itemize}

\textbf{"现在平面"为何被选中}\mtag{推导，命题链见 \S{}4.9}：涡必须锁定旋转轴（N1）且恰好一个平面（N2）；三平面等价（N3，"现在"= 被选中平面的定义）；初始微扰（= \ensuremath{\varepsilon} 泄漏事件，P1--P4）破缺 SO(3) 简并（N4）。

\subsection{连续极限的开放问题}

\mtag{待验证}(1) 从湮灭动力学证明"事件集 = 泊松涌撒"；(2) 连续极限收敛性（标准因果集理论亦未完全解决）；(3) "每事件 = 1 固有时元"公设证明。



\section{引力的涌现（核心章节）}

\subsection{通量遮挡 \ensuremath{\to} 引力势}

\mtag{推导}质量 M 的物体"吞"通量，产生稳态亏损。无源区通量亏损满足 \ensuremath{\nabla}\up{2}(\ensuremath{\delta}\ensuremath{\rho}) = 0，球对称解 \ensuremath{\delta}\ensuremath{\rho} \ensuremath{\propto} 1/r；由时间分量匹配牛顿势：

\[\frac{\rho_{pair}(r)}{\rho_{pair}(\infty)} = 1 - \frac{GM}{c^2 r} \quad\text{（一阶）}\]

弱场数值验证：太阳表面 GM/c\up{2}R = 2.1\ensuremath{\times}10\up{-6} \ensuremath{\ll} 1 \checkmark。

\textbf{因子 2 说明（v1.12 补充）}：通量亏损系数取 \textbf{1}（GM/c\up{2}r，与牛顿势 \ensuremath{\Phi}/c\up{2} 同形），而非 2GM/c\up{2}r------\textbf{g\down{00} 中的"2"来自平方结构 g\down{00} = \ensuremath{-}(\ensuremath{\rho}/\ensuremath{\rho}\ensuremath{\infty})\up{2}}（\S{}4.3，公理 3：时间=事件频率 \ensuremath{\implies} \ensuremath{\sqrt{-g\down{00}}}=\ensuremath{\rho}/\ensuremath{\rho}\ensuremath{\infty}）。即：
\begin{itemize}
\item 通量赤字 \ensuremath{\rho}\ensuremath{\infty}\ensuremath{-}\ensuremath{\rho}\down{ss} = \ensuremath{\rho}\ensuremath{\infty}\ensuremath{\cdot}GM/(c\up{2}r)：\textbf{一阶系数 1}（= \ensuremath{\Phi}/c\up{2}）；
\item g\down{00} = \ensuremath{-}(1\ensuremath{-}GM/(c\up{2}r))\up{2} = \ensuremath{-}(1\ensuremath{-}2GM/(c\up{2}r)+\ldots{})：\textbf{一阶系数 2}（与 GR 一致，\ensuremath{\beta}\down{ppn}=1 自洽）。
\end{itemize}
这同时确定有限 R 稳态（v1.11）的通量抽零半径 \textbf{R = GM/c\up{2}}（非史瓦西 2GM/c\up{2}）。数值验证：R=GM/c\up{2} 与玻尔兹曼弱场一致（r\ensuremath{\geq}10 偏差 <0.6\%），R=2GM/c\up{2} 偏差 >11\%；g\down{00} 平方结构与 GR 弱场 r=10 差 1.0\%、r=100 差 0.01\%。

\textbf{强场分叉裁决（v1.13）}\mtag{命题}:两条候选路径的 PPN \ensuremath{\beta} 参数分叉------玻尔兹曼 g\down{00}=\ensuremath{-}e\up{-2\Phi{}/c^{2}} 展开 \ensuremath{\beta}=1（通过观测 1\ensuremath{\pm}3\ensuremath{\times}10\up{-4}），平方稳态 g\down{00}=\ensuremath{-}(1\ensuremath{-}R/r)\up{2} 展开 \ensuremath{\beta}=1/2（\textbf{被 1600\ensuremath{\sigma} 排除}）。\ensuremath{\implies} \textbf{玻尔兹曼路径是弱场-中场的正确描述}；平方稳态仅为一阶近似（作为通量剖面仍有效，不能作 g\down{00} 精确映射）。\textbf{v1.12 视界预言 r=GM/c\up{2} 撤回}（该预言基于平方稳态 g\down{00}=0，已被 \ensuremath{\beta} 判决排除）。框架强场：玻尔兹曼永不归零 \ensuremath{\implies} 无 g\down{00}=0 几何视界，视界 = \textbf{风暴过载/通量抽零等势面}（物理机制，本框架设定）------与 v1.1 强场未定义、七环机制一致。

\subsection{玻尔兹曼分布（通量场热平衡）}

\mtag{推导 + 热平衡证明}通量密度按引力势的\textbf{玻尔兹曼权重}衰减：

\[\frac{\rho(r)}{\rho_\infty} = e^{-\Phi(r)/c^2} = 1 - x + \tfrac12 x^2 - \cdots,\qquad x = \frac{GM}{c^2 r}\]

\textbf{热平衡命题（六命题链）}\mtag{推导}：
\begin{enumerate}
\item 事件流为平稳随机过程（SDE dX\down{t} = b dt + \ensuremath{\sigma} dW，745 页原文）；
\item 遍历性（Birkhoff--Khinchin 定理）：时间平均 = 系综平均；
\item 微观可逆性（湮灭-重生 CPT 对称）；
\item 混沌传播（Boltzmann-Grad 极限，连接希尔伯特第六问题/邓煜-Hani-马骁硬球链条）\ensuremath{\implies} 玻尔兹曼方程；
\item 细致平衡 \ensuremath{\implies} 玻尔兹曼分布 \ensuremath{\rho} \ensuremath{\propto} e\up{-E/kT}；
\item kT = mc\up{2}（相对论极限）\ensuremath{\implies} \ensuremath{\rho} \ensuremath{\propto} e\up{-\Phi{}/c^{2}}。
\end{enumerate}

\textbf{数值模拟验证}\mtag{推导}：细致平衡马尔可夫链收敛到玻尔兹曼分布（偏差 2.12\%）；遍历性（前后两半时间平均偏差 <0.2\%）。

\textbf{重要限制}\mtag{待验证}：纯 1/r 引力势的玻尔兹曼分布发散（r\ensuremath{\to}0）------引力系统无全局热平衡（与标准统计力学一致）；e\up{-\Phi{}/c^{2}} 仅局部/弱场/截断下有效，弱场展开与 4.1 一致。

\textbf{\mtag{v2.6 补充：玻尔兹曼涌现的数值证明 + 循环风险回应（响应审稿 M1）}}：
审稿质疑"玻尔兹曼'证明'循环（假设热平衡证热平衡）"。\textbf{数值决定性回应}（`\down{m1}\down{boltzmann}.py`，输入仅为守恒律+碰撞，初始极端非平衡单能壳）：(1) \textbf{事件流语言（DSMC，N=2000）}------温度涌现 0.00\% 偏差、速度分布收敛麦克斯韦-玻尔兹曼（加权偏差 7.0\%）、\textbf{Boltzmann H 函数净不可逆下降 \ensuremath{\Delta}H=\ensuremath{-}2.27}、遍历性 0.00\%------\textbf{热平衡是长时极限的涌现产物，非输入假设}；(2) \textbf{确定性硬球（无分子混沌假设）}------温度 0.00\%、H 净下降 \ensuremath{\Delta}H=\ensuremath{-}2.38（2D long-time tail 致分布偏差 \ensuremath{\sim}18\%，诚实标注）；(3) \textbf{推导链前提显式化}------平稳性（公理 1 各向同性）、遍历性/混沌传播（动力学性质，实测 0.00\%）、细致平衡（事件流时间反演可逆）------全部来自公理 1 层级，\textbf{非热平衡假设}。

\textbf{事件流哈密顿量候选（H\down{event}）}：事件流从隐喻升级为可计算模型------经典层 H=\ensuremath{\Sigma}p\up{2}/2m+\ensuremath{\Sigma}V\down{contact}（碰撞=相遇-湮灭-再生）；统计层 S=Tr f(D\up{2}/\ensuremath{\Lambda}\up{2})（谱作用，\S{}5.4）；量子层开放系统 Lindblad d\ensuremath{\rho}/dt=\ensuremath{-}i[H\down{0},\ensuremath{\rho}]+\ensuremath{\Sigma}(L\ensuremath{\rho}L\ensuremath{^\dagger}\ensuremath{-}\ensuremath{\tfrac12}{L\ensuremath{^\dagger}L,\ensuremath{\rho}})（湮灭=耗散/再生=泵浦，稳态=细致平衡\ensuremath{\implies}玻尔兹曼）。可证伪：各层给出碰撞率/驰豫时标/稳态分布预言。

\textbf{与 LHC 衔接}：事件流单体能标 E\ensuremath{\sim}E\down{P}\ensuremath{\approx}10\up{19} GeV \ensuremath{\gg} LHC \ensuremath{\sim}10\up{4} GeV（差 15 量级）\ensuremath{\implies} 事件流仅以集体模态显现，\textbf{预言 LHC 无新物理}（与空结果一致）；各向同性+泊松涌撒 \ensuremath{\implies} 不破坏洛伦兹/CPT；泄漏 \ensuremath{\varepsilon}\ensuremath{\sim}10\up{-50} 不产生 LHC 信号；框架判别预言集中在引力侧（S-E1/E5/S4）。\textbf{诚实保留}：事件流为公理 1（与标准模型拉氏量同为输入）；H\down{quantum} 微观结构、遍历性解析证明开放（\S{}10.2 \#4）。

\subsection{时间分量与 PPN \ensuremath{\gamma}=1}

\mtag{推导}由引理 g\down{00} = \ensuremath{-}(\ensuremath{\rho}/\ensuremath{\rho}\ensuremath{\infty})\up{2}：

\[g_{00} = -e^{-2\Phi/c^2} = -\big(1 - 2x + 2x^2 - \tfrac43 x^3 + \cdots\big)\]

\subsection{微变权守恒命题（空间分量）}

\textbf{\mtag{命题（微变权守恒 \ensuremath{\Rightarrow} \ensuremath{\gamma}=1）}}：设湮灭事件流满足守恒（公理 1'）与局域各向同性（公理 1），则静态球对称弱场中 g\down{rr} = \ensuremath{-}1/g\down{00} = (1\ensuremath{-}2GM/c\up{2}r)\up{-1}，\textbf{PPN \ensuremath{\gamma} = 1}。

\textbf{证明}（三行）：(1) 局域惯性系各向同性 \ensuremath{\Rightarrow} 时间权 \ensuremath{\sqrt{-g\down{00}}} 与径向权 \ensuremath{\sqrt{g\down{rr}}} 互补（乘积=1）；(2) 引理：\ensuremath{\sqrt{-g\down{00}}} = f(r)/f\down{0} = e\up{-x}；(3) 代入 \ensuremath{\Rightarrow} \ensuremath{\gamma}=1。\ensuremath{\blacksquare}

\subsection{PPN \ensuremath{\beta}=1（二阶项）}

\mtag{推导}玻尔兹曼分布的二阶系数 a = +1/2：

\[g_{00} = -e^{-2x} = -\big(1 - 2x + 2\beta x^2 + \cdots\big) \quad\Longrightarrow\quad \beta_{PPN} = 1\]

\textbf{数值验证}：PPN 展开对照（Will 约定 g\down{00} = \ensuremath{-}(1\ensuremath{-}2x+2\ensuremath{\beta}x\up{2})）\ensuremath{\implies} \ensuremath{\beta} = a + 1/2 = 1 \checkmark。

\subsection{完整度规 = 各向同性指数度规 + 完整 PPN 参数集（v1.21 新增）}

\textbf{\mtag{推导 + 数值}三条独立推导收敛为唯一的坐标形式}------玻尔兹曼 g\down{00}=\ensuremath{-}e\up{-2x}（\S{}4.2）+ 微变权守恒 g\down{rr}=e\up{2x}（\S{}4.4）+ 公理 1 各向同性（g\down{ij}=\ensuremath{\delta}\down{ij}\ensuremath{\cdot}e\up{2x}）：

\[g_{\mu\nu} = \mathrm{diag}\!\big(-e^{-2x},\ e^{2x},\ r^2 e^{2x},\ r^2\sin^2\theta\, e^{2x}\big),\qquad x = \frac{GM}{c^2 r}\]

互洽检验：微变权守恒 g\down{rr}\ensuremath{\cdot}(\ensuremath{-}g\down{00}) = e\up{2x}\ensuremath{\cdot}e\up{-2x} = \textbf{1 精确成立}------三条推导无坐标矛盾。\mtag{引用}该度规即文献中已知的 exponential metric（Edery 2005 等）------框架从公理独立导出，与主流文献直接对接。

\textbf{完整 PPN 参数集}：静态各向同性纯对角度规 \ensuremath{\implies}

\begin{table}[htbp]
\centering
\small
\begin{tabularx}{\textwidth}{@{}lXXX@{}}
\toprule
参数 & \ensuremath{\gamma} & \ensuremath{\beta} & \ensuremath{\xi}, \ensuremath{\alpha}\down{1}, \ensuremath{\alpha}\down{2}, \ensuremath{\alpha}\down{3}, \ensuremath{\zeta}\down{1}--\ensuremath{\zeta}\down{4}, \ensuremath{\zeta}\down{W} \\
\midrule
框架值 & \textbf{1} & \textbf{1} & \textbf{0（自动 GR 值）} \\
\bottomrule
\end{tabularx}
\end{table}

\textbf{1PN 与 GR 完全不可区分}；分叉点精确定位（泰勒展开对照）：

\begin{table}[htbp]
\centering
\footnotesize
\begin{tabularx}{\textwidth}{@{}lXXXXX@{}}
\toprule
阶 & 框架 g\down{00} & GR g\down{00} & 框架 g\down{rr} & GR g\down{rr} & 分叉 \\
\midrule
x\up{1} & \ensuremath{-}2 & \ensuremath{-}2 & +2 & +2 & 无（1PN） \\
x\up{2} & +2 & +2 & \textbf{+2} & \textbf{+1.5} & g\down{rr} 33\%（2PN） \\
x\up{3} & \textbf{\ensuremath{-}4/3} & \textbf{\ensuremath{-}3/2} & +4/3 & +0.5 & g\down{00} 11\%（3PN） \\
\bottomrule
\end{tabularx}
\end{table}

\textbf{数值验证（测地线 RK4）}：近日点进动框架比 GR 大 0.57\%/1.02\%/1.55\%（相对差 \ensuremath{\delta}=0.30\ensuremath{\cdot}x\down{p} 线性，2PN 分叉量化）；光子偏折框架 > GR（b=10 差 5\%）。2PN 可观测效应 \ensuremath{\propto}U\up{2} 当前不可分辨（双脉冲星），3PN 需 LISA EMRI/脉冲星计时阵列未来检验。\textbf{\mtag{v1.26 补充}LISA EMRI 检验提案（候选 E，定量可检验性论证}\mtag{数值}）：EMRI（q\ensuremath{\sim}10\up{4}--10\up{7}）轨道圈数 N\ensuremath{\sim}10\up{4}--10\up{5} 把进动差放大为可测相位差------\textbf{\ensuremath{\Delta}\ensuremath{\Phi}\down{prec} = N\ensuremath{\cdot}\ensuremath{\Delta}\ensuremath{\omega}\down{GR}\ensuremath{\cdot}\ensuremath{\delta}}（\ensuremath{\delta}=0.30\ensuremath{\cdot}x\down{p}，x\down{p}=GM/c\up{2}p）：M=10\up{6} M\ensuremath{\odot}、\ensuremath{\mu}=10 M\ensuremath{\odot}、e=0.3、p=15M 源 \ensuremath{\Delta}\ensuremath{\Phi}=776 rad（2yr）/1552 rad（4yr），裕度（=\ensuremath{\Delta}\ensuremath{\Phi}\ensuremath{\cdot}SNR，SNR\ensuremath{\sim}100 文献值）\textbf{77610/232831}；6 源案例 + 轨道半径扫描 p\ensuremath{\in}[10,50]M \textbf{全部可检验}（裕度 1128--513240）；3PN g\down{00} 11\% 分叉 \ensuremath{\implies} 啁啾相位差 \ensuremath{\Delta}\ensuremath{\Phi}\down{chirp}\ensuremath{\sim}206--411 rad \checkmark------\textbf{EMRI 是框架 2PN g\down{rr} 33\% 分叉的第一道观测判决窗口}（进动差-半径标度律 \ensuremath{\delta}\ensuremath{\propto}x\down{p} 为框架独有签名，S-E1--S-E4 见 \S{}9；与 PTA 互补）。\mtag{诚实标注}提案基于解析相位累积估计；Kerr 版本未建立。\textbf{\mtag{v1.27 补充}EMRI 完整波形模板 + Fisher 精度（后续工作(4)，提案升级）}\mtag{数值}：Peters--Mathews 椭圆绝热演化（quadrupole 辐射收缩）+ 主导谐波波形 h(t)=h\down{0}\ensuremath{\cdot}cos[2\ensuremath{\lambda}+2(1+\ensuremath{\delta})\ensuremath{\omega}\down{peri}+\ensuremath{\varphi}\down{0}]（进动差 \ensuremath{\delta} 显式注入）+ LISA 噪声（Barack--Cutler 2004 拟合）+ \textbf{7 参数 Fisher 矩阵} {\ensuremath{\delta},lnM,ln\ensuremath{\mu},a\down{0},e\down{0},\ensuremath{\varphi}\down{0},lnD}（含参数简并）：代表源（M=10\up{6} M\ensuremath{\odot}、\ensuremath{\mu}=10 M\ensuremath{\odot}、e\down{0}=0.3、p\down{0}=15M、D=3Gpc）\ensuremath{\delta}\down{true}=2.00\% \ensuremath{\implies} \textbf{\ensuremath{\sigma}\down{\delta{}}=8.3\ensuremath{\times}10\up{-4}（24\ensuremath{\sigma}，2yr）/1.2\ensuremath{\times}10\up{-4}（167\ensuremath{\sigma}，4yr）}，\ensuremath{\sigma}\down{lnM}=0.03\%/0.004\%，\textbf{GR 模型被排除 \ensuremath{\Delta}lnL=288（2yr）/13930（4yr）------单源判决（模板内）}；e\down{0}=0.7 源 13.6\ensuremath{\sigma}、\ensuremath{\mu}=1 源 3.5\ensuremath{\sigma}（边缘）。\mtag{诚实标注}模板为 quadrupole 主导谐波（无高次谐波/Kerr/自旋-轨道），Fisher 线性化，\ensuremath{\sigma}\down{\delta{}} 为模板内下界------真实波形参数简并或更强。\textbf{\mtag{v1.32 补充}Kerr EMRI 波形与贝叶斯分析（后续工作(4)，推广到旋转中心体）}\mtag{推导 + 数值}：静态模板推广到 \textbf{Kerr（a=0.9）}------Kerr 圆轨道频率 \ensuremath{\Omega}=\ensuremath{\sqrt{GM}}/[r\up{3/2}+a\ensuremath{\sqrt{GM/c\up{2}}}] + quadrupole 绝热收缩 + 框架进动差注入 \ensuremath{\omega}\down{total}=(1+\ensuremath{\delta})\ensuremath{\omega}\down{GR}（\ensuremath{\delta}=0.30\ensuremath{\cdot}x\down{p}）；方法升级为\textbf{贝叶斯（剖面似然 + 1D MCMC）}：代表源（M=10\up{6} M\ensuremath{\odot}、\ensuremath{\mu}=10 M\ensuremath{\odot}、a=0.9、r\down{0}=15 GM/c\up{2}、D=3 Gpc、2yr、SNR=20）------\textbf{\ensuremath{\delta}\down{hat}=0.02000 精确恢复真值（0\ensuremath{\sigma} 偏离）、\ensuremath{\sigma}\down{\delta{}}=9.7\ensuremath{\times}10\up{-4}、\ensuremath{\Delta}lnL=338.4 \ensuremath{\implies} GR 模型被排除 \ensuremath{\sim}26\ensuremath{\sigma}（Kerr 下稳健）}；GR 数据对照 \ensuremath{\delta}\down{hat}=\ensuremath{-}0.00000\ensuremath{\pm}1.9\ensuremath{\times}10\up{-3}（\textbf{无假阳性}）；\ensuremath{\sigma}\down{\delta{}} 与静态 Fisher（8.3\ensuremath{\times}10\up{-4}）同量级------\textbf{框架 2PN 进动差对旋转黑洞单源判决成立}（v1.27 结论稳健）。\mtag{诚实标注}自旋 a 固定（a-\ensuremath{\delta} 联合后验未做，真实分析 \ensuremath{\sigma}\down{\delta{}} 或增大）、圆轨道简化（无偏心率/LT 节点进动波形调制）、非 Teukolsky 数值波形。\textbf{\mtag{v1.35 补充}EMRI a-\ensuremath{\delta} 联合 MCMC（行动清单 P1-2，a 自由落地；发现并修正 v1.32 严重数值缺陷）}\mtag{数值}：三处重大修复------(1) \textbf{欠采样}：v1.32 采样 dt/载波周期=34.6（严重欠采样）\ensuremath{\to} \ensuremath{\delta} 似然混叠成梳状、粗网格取随机齿；改为\textbf{解析旋近 r(t)=(r\down{0}\up{4}\ensuremath{-}kt)\up{1/4} + dt=400s 奈奎斯特采样}；(2) \textbf{轨道相位梳}：a 进入快速轨道相位 2\ensuremath{\lambda}（2yr 累积 4.3\ensuremath{\times}10\up{5} rad）使 a-似然呈周期 \ensuremath{\sim}7\ensuremath{\times}10\up{-4} 梳状、粗网格边际化被毁；改为 \textbf{\ensuremath{\varphi}0 剖面似然 |\ensuremath{\Sigma}d\ensuremath{\cdot}e\up{-i\Phi{}}|} 消除相位梳；(3) \textbf{Lense--Thirring 进动加入}：\ensuremath{\omega}\down{peri}=(1+\ensuremath{\delta})\ensuremath{\omega}\down{GR}+\ensuremath{\omega}\down{LT}(a)（LT\ensuremath{\propto}x\up{3/2} vs 1PN\ensuremath{\propto}x，r-标度差破除 \ensuremath{\delta}-a 简并）。\textbf{核心结果}：(i) \textbf{\ensuremath{\sigma}\down{\delta{}} 修正 600\ensuremath{\times}}：\ensuremath{\sigma}\down{\delta{}}(1D, \ensuremath{\varphi}0 剖面, a 固定) = \textbf{1.63\ensuremath{\times}10\up{-6}（2yr）/7.37\ensuremath{\times}10\up{-7}（4yr）}（v1.32 的 9.7\ensuremath{\times}10\up{-4} 因欠采样+网格虚高；物理 \ensuremath{\sigma}\down{\delta{}}\ensuremath{\approx}1/(2\ensuremath{\cdot}wGR\down{total}\ensuremath{\cdot}SNR)）；(ii) \textbf{a 自由下简并很小}：\ensuremath{\sigma}\down{a}=1.5\ensuremath{\times}10\up{-5}（LT 相位 2aW=1.4\ensuremath{\times}10\up{4} rad 钉死 a），\ensuremath{\sigma}\down{\delta{}}(3D 边际)\ensuremath{\approx}3.8\ensuremath{\times}10\up{-6}（膨胀仅 2.4\ensuremath{\times}）；(iii) \textbf{\ensuremath{\delta}=2\% 检测稳健}：\textbf{\ensuremath{\Delta}lnL(GR)=357（2yr）/396（4yr）\ensuremath{\implies} 27\ensuremath{\sigma}/28\ensuremath{\sigma}（a-\ensuremath{\varphi}0 全自由）}------v1.32 的 4yr"167\ensuremath{\sigma}"（\ensuremath{\Delta}lnL=13930）是混叠似然虚高，真实值 28\ensuremath{\sigma}；(iv) \textbf{主简并识别}：\ensuremath{\delta}-\ensuremath{\varphi}0（\ensuremath{\rho}\ensuremath{\approx}\ensuremath{-}0.85，轨道相位）\textbf{而非 \ensuremath{\delta}-a}（LT/1PN r-标度差+啁啾破除）------解释 v1.32 固定 \ensuremath{\varphi}0 获得"过高精度"的机制；(v) GR 对照 \ensuremath{\delta}\down{hat}\ensuremath{\approx}0、\ensuremath{\Delta}lnL\ensuremath{\approx}0（无假阳性）。\mtag{诚实标注}3D MCMC 混合困难（\ensuremath{\sigma}\down{a} 极窄）、\ensuremath{\sigma}\down{\delta{}} 为局部曲率（稳健判据是 \ensuremath{\Delta}lnL）、\ensuremath{\varphi}0 剖面 SNR=20 噪声底板、圆轨道/quadrupole/非 Teukolsky------\textbf{S-E1 预言在自旋与初始相位全自由下仍以 27-28\ensuremath{\sigma} 可测（观测判决不受简并威胁）}。

\textbf{\mtag{v2.4 补充}与 exponential metric 文献的对照（响应第三方审稿 M4）}\mtag{文献对照}：本框架指数度规 g\down{\mu{}\nu{}}=diag(\ensuremath{-}e\up{-2x}, e\up{2x}, r\up{2}e\up{2x}, r\up{2}sin\up{2}\ensuremath{\theta}e\up{2x}) \textbf{即 Yilmaz--Rosen exponential metric}（Papapetrou 1954 首次提出、Yilmaz 1958/1971、Rosen 1962；Vaidya 1954/Coleman 1971/Harvey 1972；近期 Makukov \& Mychelkin 2018、Boonserm et al. 2021）。\textbf{文献已知性质与本框架对照}：(i) \textbf{PPN \ensuremath{\alpha}=\ensuremath{\beta}=\ensuremath{\gamma}=1}（1PN 与 GR 不可区分）------文献已确立（Yilmaz 传统），本框架为其涌现机制提供统计力学来源（\S{}4.2 通量玻尔兹曼）而非假设；(ii) \textbf{2PN g\down{rr} 分叉 33\% / 3PN g\down{00} 分叉 11\% = 文献已知高阶 PPN 性质}（Yilmaz \ensuremath{\delta}=1 vs GR \ensuremath{\delta}=3/4；本框架解析展开验证一致：g\down{rr}=e\up{2x}=1+2x+2x\up{2}+\ensuremath{\cdots} vs GR (1+x/2)\up{4}=1+2x+1.5x\up{2}+\ensuremath{\cdots} \ensuremath{\implies} 33\%；g\down{00}=\ensuremath{-}e\up{-2x} vs GR \ensuremath{-}[(1\ensuremath{-}x/2)/(1+x/2)]\up{2} \ensuremath{\implies} 3PN 11\%）------\textbf{数值一致，但本框架的贡献是机制的涌现推导与判别方案}（EMRI \ensuremath{\delta}=0.30\ensuremath{\cdot}x\down{p} 标度律，S-E1--E4）；(iii) \textbf{无事件视界}（g\down{00} 永不为零）------文献已知，与本框架 v1.13 撤回几何视界、视界=风暴过载等势面\textbf{自洽}；(iv) \textbf{exponential metric 非 EFE 真空解}（R\down{\mu{}\nu{}}=\ensuremath{-}\ensuremath{\tfrac12}\ensuremath{\nabla}\ensuremath{\varphi}\ensuremath{\nabla}\ensuremath{\varphi}，需负动能 ghost 标量场支撑）------本框架以"质量遮蔽通量产生负压差"（\S{}4.1）作为该源项的物理对应，需在后续工作中严格化；(v) \textbf{可穿越虫洞}（Boonserm et al. 2021：YEM=traversable wormhole）------与本框架"跨支 ER 通道/量子直线"（\S{}4.6）对应。\textbf{涌现引力文献定位}：本框架 \S{}4.7 视界热力学路径 = Jacobson 1995 方法（Clausius \ensuremath{\implies} Einstein）；\S{}4.2 通量玻尔兹曼机制区别于 Verlinde 2011 熵引力与 Padmanabhan 2010 均衡（机制不同、结论相容）；\S{}3 因果集与 \S{}5.8 弦网均为标准工具（Bombelli 1987 / Malament 1977 / Wen 2003 / Kitaev 2003）。\mtag{创新性重估}度规与分叉为文献已知；\textbf{本框架新增}：玻尔兹曼机制涌现度规、通量遮挡统一语言、EMRI 判别标度律、钳制物态、packing 定理包、弦网严格化系列（详见《M4补强报告》）。

\subsection{强场机制与黑洞（修订版）}

\textbf{\mtag{v1.1 重要修订}}：早期版本（v1.0）将弱场玻尔兹曼度规 g\down{00}=\ensuremath{-}e\up{-2GM/c^{2}r} 外推到强场，得出"光子球 1GM vs 3GM"的 EHT 冲突预言。经复核，\textbf{该外推超出框架适用域}（热平衡证明已确认纯引力势无全局热平衡，e\up{-\Phi{}/c^{2}} 仅弱场/局部/截断内有效）------\textbf{正式撤回}。

\textbf{\mtag{v1.14 补充}强场度规的数学排除}：尝试在玻尔兹曼弱场形式（\ensuremath{\beta}=1 存活）上叠加通量饱和截断（软核/线性/对数），以 GR 强场事实（视界 2GM、光子球 3GM、阴影 5.196GM）为基准检验------\textbf{全部失败}：(1) 纯玻尔兹曼光子球解析证明在 1GM（V=e\up{-2x}x\up{2}，dV/dx=0\ensuremath{\implies}x=1）；(2) 饱和修正光子球 2.16--4.56GM（偏差 9--52\%）；(3) 饱和因子在弱场（x=0.01）泄漏破坏 \ensuremath{\beta}=1。\textbf{\ensuremath{\implies} 所有基于玻尔兹曼的饱和强场度规被数学排除，框架没有强场度规公式。} 强场 = 风暴过载物理机制（视界 = 压强临界等势面，非 g\down{00}=0 几何奇点）；与 EHT 无冲突（无强场度规承诺）。

\textbf{\mtag{v1.15 补充}风暴过载观测效应（单参数族 \ensuremath{\alpha}）}：框架强场观测内容编码为最小承诺单参数族------通量抽零半径 r\down{h} = \ensuremath{\alpha}\ensuremath{\cdot}2GM/c\up{2}（\ensuremath{\alpha}=1 为 GR 极限），预言：阴影 b=5.196\ensuremath{\alpha}\ensuremath{\cdot}GM/c\up{2}、QNM f=f\down{GR}/\ensuremath{\alpha}、火墙弥散角 \ensuremath{\delta}/R=\ensuremath{\sqrt{2l\down{P}/R}}。\textbf{EHT 约束 \ensuremath{\alpha} = 1.00\ensuremath{\pm}0.06（M87*/Sgr A* 阴影 5.2 GM）------框架与黑洞阴影观测完全兼容}；未来 ET/CE QNM 精度 \ensuremath{\sim}1\% \ensuremath{\implies} \ensuremath{\alpha} 约束 \ensuremath{\pm}0.01（第一道精确判决）；\textbf{框架独有签名：火墙弥散 \ensuremath{\delta}/R（原初黑洞 \ensuremath{\sim}10\up{-9} 可测，GR 阴影锐利）}；过载区无奇点（风暴压制，假设 H-7）vs GR 奇点。预言清单 S1--S5 见 \S{}9。

\textbf{\mtag{v1.16 补充}\ensuremath{\alpha} 微观来源 = 熵一致性（理论必然，非拟合）}：\ensuremath{\alpha} 由框架热力学自洽性强制为 1------(1) 框架挂谷微态计数\mtag{211/222：每端点独占 4\ensuremath{\pi}/4 方向锥}S=A/(4l\down{P}\up{2}) 与 Bekenstein-Hawking 系数同为 1/4 \ensuremath{\implies} 熵一致性 \ensuremath{\implies} 面积一致性 A\down{过载}=A\down{BH} \ensuremath{\implies} \ensuremath{\alpha}\up{2}=1 \ensuremath{\implies} \textbf{\ensuremath{\alpha}=1（唯一解）}；(2) Jacobson 路径交叉验证：\ensuremath{\alpha}\ensuremath{\neq}1 \ensuremath{\implies} G\down{eff}=G/\ensuremath{\alpha}\up{2}\ensuremath{\neq}G，与 PPN/引力波观测（1\ensuremath{\pm}10\up{-5}）冲突；(3) 数值：\ensuremath{\alpha}=1 时太阳质量黑洞 S=1.05\ensuremath{\times}10\up{77} 纳特精确等于标准 BH。\textbf{\ensuremath{\implies} EHT 观测 \ensuremath{\alpha}=1.00\ensuremath{\pm}0.06 是框架自洽性的直接验证}------过载区半径被熵一致性锁定为史瓦西半径。

\textbf{\mtag{v1.17 补充}压强临界实现面积锁定的微观机制（M\up{2}/A 普朗克临界）}：原文\mtag{88}"黑洞只与 M/A 有关"在 GR 中 M/A=c\up{4}/(16\ensuremath{\pi}GM)（依赖 M）；\textbf{平方不变量 M\up{2}/A = c\up{4}/(16\ensuremath{\pi}G\up{2}) \ensuremath{\approx} 3.6\ensuremath{\times}10\up{52} kg\up{2}/m\up{2}（普朗克常数，普适临界）}\ensuremath{\implies} A=16\ensuremath{\pi}G\up{2}M\up{2}/c\up{4} \ensuremath{\implies} r\down{h}=2GM/c\up{2}（精确，太阳/M87*/Sgr A* 全验证）------风暴压制强度 \ensuremath{\propto} M\up{2}（引力自束缚）、表面抵抗 \ensuremath{\propto} A，压制/抵抗比达普朗克常数 = 史瓦西面。候选排除：引力自压强 GM\up{2}/(4\ensuremath{\pi}r\up{4})（r\down{h} 偏差 98\%）、(M/A)c\up{2}/r（r\down{h}\ensuremath{\propto}M\up{1/3}）均不匹配。\textbf{\ensuremath{\alpha}=1 获得双重来源：熵一致性（v1.16 热力学路径）+ M\up{2}/A 普朗克临界（本批力学路径）------面积锁定从热力学要求升级为风暴压强的物理平衡面。}

\textbf{\mtag{v1.18 补充 + v1.24 微观推导}过载区流体模型（正物质浆糊物态）}：\mtag{假设 H-7}"最大熵状态的正物质浆糊\ldots{}\ldots{}风暴压强有限（普朗克尺度限制）"\ensuremath{\implies} 物态方程 \textbf{P(\ensuremath{\rho}) = P\down{P}\ensuremath{\cdot}tanh(\ensuremath{\rho}/\ensuremath{\rho}\down{P})}（普朗克钳制：线性区恢复 P\ensuremath{\approx}c\up{2}\ensuremath{\rho}，饱和区压强钳制在 P\down{P}\ensuremath{\approx}4.6\ensuremath{\times}10\up{113} Pa------\textbf{避免奇点的定量表达}）。\textbf{v1.24 微观推导（候选 D，从唯象假设升级为涌现）}------三条机制全部来自框架既有前提、无新假设：
\begin{itemize}
\item \textbf{机制(1) 普朗克填充率}：\ensuremath{\rho}/\ensuremath{\rho}\down{P} 中 \ensuremath{\rho}\down{P} = m\down{P}/l\down{P}\up{3}（每普朗克体积一个碎片质量）------几何标度，非拟合参数；
\item \textbf{机制(2) tanh 涌现}：过载区压强 = 碎片-风暴双向净动量流（反射 +J\down{max} / 放行 \ensuremath{-}J\down{max}），净流满足罗吉斯蒂饱和 dJ/d\ensuremath{\rho}=(1/\ensuremath{\rho}\down{P} J\down{max})(J\down{max}\up{2}\ensuremath{-}J\up{2}) \ensuremath{\implies} \textbf{J=J\down{max}\ensuremath{\cdot}tanh(\ensuremath{\rho}/\ensuremath{\rho}\down{P})}（数值积分 vs 解析偏差 8\ensuremath{\times}10\up{-4}）------两通道竞争 \ensuremath{\implies} 双曲正切；
\item \textbf{机制(3) 普朗克压强}：P\down{P} = E\down{P}/l\down{P}\up{3} = 单次碰撞动量 p\down{P}=\ensuremath{\hbar}/l\down{P} \ensuremath{\times} 最大碰撞频率 \ensuremath{\nu}\down{P}=c/l\down{P}（公理 1 事件流上限；\mtag{公理 1 推论}"吞灭有最大速率"）/ 单元面积 l\down{P}\up{2}；
\item \textbf{附产物}：合并质量 m\down{eff} = m\down{P}\ensuremath{\cdot}(\ensuremath{\rho}/\ensuremath{\rho}\down{P})/tanh(\ensuremath{\rho}/\ensuremath{\rho}\down{P})（碎片独立 \ensuremath{\to} 风暴压成团块，浆糊微观图像）；\textbf{新预言：声速剖面 c\down{s}\up{2} = \ensuremath{\partial}P/\ensuremath{\partial}\ensuremath{\rho} = c\up{2}\ensuremath{\cdot}sech\up{2}(\ensuremath{\rho}/\ensuremath{\rho}\down{P})}------\ensuremath{\rho}\ensuremath{\ll}\ensuremath{\rho}\down{P} 时 c\down{s}=c（相对论性最大因果）、\ensuremath{\rho}\ensuremath{\gg}\ensuremath{\rho}\down{P} 时 c\down{s}\ensuremath{\to}0（\textbf{钳制区无弹性波 = 浆糊无结构的动力学表达}，衔接挂谷集信息\ensuremath{\approx}0 与"过载区无特征辐射"）；\textbf{能量条件}：P \ensuremath{\leq} \ensuremath{\rho}c\up{2} 处处满足（\ensuremath{\rho}\ensuremath{\gg}\ensuremath{\rho}\down{P} 时深度亚因果 P/\ensuremath{\rho}c\up{2}\ensuremath{\to}\ensuremath{\rho}\down{P}/\ensuremath{\rho}\ensuremath{\to}0）、\ensuremath{\partial}P/\ensuremath{\partial}\ensuremath{\rho}\ensuremath{\geq}0 \ensuremath{\implies} 无热力学不稳定性 \ensuremath{\implies} 无坍缩（对比 GR 奇点）；
\end{itemize}
数值：宏观黑洞过载区平均密度远低于普朗克密度（太阳 \ensuremath{\rho}\down{avg}=1.8\ensuremath{\times}10\up{19} vs \ensuremath{\rho}\down{P}=5.2\ensuremath{\times}10\up{96} kg/m\up{3}，差 78 个数量级，P\ensuremath{\approx}10\up{-6} P\down{P}、c\down{s}\ensuremath{\approx}c），仅普朗克质量黑洞接近 \ensuremath{\rho}\down{P}（\ensuremath{\rho}\down{avg}=0.03\ensuremath{\rho}\down{P}、P=0.030P\down{P}、c\down{s}=0.9996c）------\textbf{奇点天然消除}；吸积行为：物质被风暴压碎成正物质碎片（非 GR 自由下落\ensuremath{\to}奇点）；辐射：过载区=最大熵无结构 \ensuremath{\implies} 无特征辐射（辐射来自视界表层饱和模态+霍金辐射），\mtag{223 页预言}霍金辐射子集与落入物质纠缠（信息翻转签名，S4 观测窗口）。

\textbf{\mtag{v1.25 补充}过载区径向结构与辐射传输（遗留(2)(3)关闭）}：\textbf{等温球排除（推导 + 数值）}------假设过载区内部有牛顿引力梯度（线性区 P\ensuremath{\approx}c\up{2}\ensuremath{\rho} 静水平衡）\ensuremath{\implies} 奇异等温球 \ensuremath{\rho}(r)=c\up{2}/(2\ensuremath{\pi}Gr\up{2})、M(r)=2c\up{2}r/G \ensuremath{\implies} \textbf{M(r\down{h})=4M 恒定超重}（太阳/Sgr A*/M87*/原初 10\up{11}/10\up{13} kg/普朗克质量 6 个质量全验证 4.0\ensuremath{\times}M，奇异等温球无标度自相似）\ensuremath{\implies} 与总质量守恒矛盾；\textbf{框架机制印证}：视界=通量=0 等势面 + 风暴压制区（七环环 1）\ensuremath{\implies} 内部无净通量梯度 \ensuremath{\implies} 无引力（引力=通量遮挡亏损消失）\ensuremath{\implies} dP/dr=0 \ensuremath{\implies} \textbf{均匀浆糊} \ensuremath{\bar{\rho}} = 3c\up{6}/(32\ensuremath{\pi}G\up{3}M\up{2})（总质量+体积锁定，与 v1.18 \ensuremath{\rho}\down{avg} 精确一致：太阳 1.842\ensuremath{\times}10\up{19} kg/m\up{3} \checkmark、Sgr A* 3.94\ensuremath{\times}10\up{6} \checkmark、M87* 1.73 \checkmark）；P\ensuremath{\approx}c\up{2}\ensuremath{\bar{\rho}}（线性区）、c\down{s}\ensuremath{\approx}c------\textbf{"均匀=最大熵无结构"的动力学表达}（无内部梯度 \ensuremath{\implies} 无内部结构 \ensuremath{\implies} 无特征辐射）。\textbf{辐射传输（定量）}：碎片截面 \ensuremath{\sigma}=l\down{P}\up{2}（普朗克单元最小截面假设）\ensuremath{\implies} 光学深度 \ensuremath{\tau}=(\ensuremath{\bar{\rho}}/m\down{P})l\down{P}\up{2}r\down{h}=3c\up{4}l\down{P}\up{2}/(16\ensuremath{\pi}G\up{2}m\down{P} M)------\textbf{宏观黑洞 \ensuremath{\tau}\ensuremath{\sim}10\up{-40}\ensuremath{\sim}10\up{-49}（光学透明，过载区不吸收/不发射辐射的定量原因，落入辐射直接穿过）}、原初 10\up{11} kg \ensuremath{\tau}=1.3\ensuremath{\times}10\up{-20}、\textbf{普朗克质量黑洞 \ensuremath{\tau}=0.0597（T=0.942，过载区首次"可见"的临界窗口）}。

\textbf{\mtag{v1.31 补充}吸积盘辐射传输（后续工作(3)，遗留(3)关闭）}\mtag{引用 + 数值}：标准薄盘（Shakura--Sunyaev 多温度黑体 T(r)=[3GM\ensuremath{\dot{M}}/(8\ensuremath{\pi}\ensuremath{\sigma}r\up{3})(1\ensuremath{-}\ensuremath{\sqrt{r\down{in}/r}})]\up{1/4}，r\down{in}=6GM/c\up{2}）为框架吸积盘\textbf{主体------与 GR 完全一致}（外区几何 1PN 不可区分，\S{}4.5b；与 EHT 相容）：M=10\up{6} M\ensuremath{\odot}、\ensuremath{\dot{M}}=0.01\ensuremath{\dot{M}}\down{Edd} 时 T(2r\down{in})=\textbf{1.2\ensuremath{\times}10\up{5} K}、谱峰 \ensuremath{\nu}\ensuremath{\approx}6.7\ensuremath{\times}10\up{15} Hz（X 射线）、L\down{acc}=1.05\ensuremath{\times}10\up{35} W\ensuremath{\approx}0.01L\down{Edd}；\textbf{过载区辐射传输}：\ensuremath{\tau}\ensuremath{\sim}10\up{-40}（太阳/Sgr A*/M87*）、原初 10\up{11} kg \ensuremath{\tau}=1.3\ensuremath{\times}10\up{-20}------\textbf{落入吸积盘辐射穿过过载区无次级吸收}（v1.25 定量延续，过载区不改变盘辐射）；\textbf{框架独有差异在视界表层}：火墙发射层（F4）非热修正 \ensuremath{\sim}\ensuremath{\sqrt{l\down{P}/R}}（太阳 7.4\ensuremath{\times}10\up{-20}、\textbf{原初 10\up{11} kg 3.3\ensuremath{\times}10\up{-10}}、普朗克 0.71）+ \textbf{霍金辐射-落入物质纠缠}\mtag{本框架预言 S4}------原初黑洞（10\up{11}--10\up{13} kg，T\down{H}=10\up{10}--10\up{13} K）蒸发末期 \ensuremath{\gamma} 射线窗口含框架独有签名。\mtag{诚实标注}薄盘为牛顿近似（无 GR 红移/光线）、火墙层内传输未建模、Kerr 盘未处理。

\textbf{\mtag{v1.21 补充}各向同性指数度规的阴影与 EHT 相容（修正 v1.14 结论）}：v1.14 的"玻尔兹曼光子球 1GM 与 EHT 冲突"源于\textbf{标准坐标约定}（g\down{\theta{}\theta{}}=r\up{2}、角向不弯曲，其阴影 b\down{min}=e=2.718 GM/c\up{2} 确实冲突）；但 \S{}4.5b 识别的\textbf{各向同性指数度规}（公理 1 各向同性的自然选择）给出坐标不变冲击参数 b\down{min}=\ensuremath{\sqrt{g\down{\varphi{}\varphi{}}/|g\down{00}|}} 的极小值：

\begin{table}[htbp]
\centering
\small
\begin{tabularx}{\textwidth}{@{}lXXX@{}}
\toprule
度规 & 光子球 r\down{ph} & 阴影 b\down{min} & 与 EHT（5.2\ensuremath{\pm}0.3 GM/c\up{2}） \\
\midrule
\textbf{框架各向同性指数度规} & \textbf{2.00 GM/c\up{2}} & \textbf{5.436 GM/c\up{2}（=2e）} & \textbf{\checkmark 相容（+4.55\% < 5.8\% 误差）} \\
GR Schwarzschild & 3.00 GM/c\up{2} & 5.196 GM/c\up{2}（=3\ensuremath{\sqrt{3}}） & 基准 \\
{[对比] 标准坐标 g\down{\theta{}\theta{}}=r\up{2}} & 1.00 GM/c\up{2} & 2.718 GM/c\up{2}（=e） & \ensuremath{\times} 冲突（v1.14） \\
\bottomrule
\end{tabularx}
\end{table}

\textbf{框架强场从"无度规承诺（v1.14）"升级为"有与观测相容的几何预言（v1.21）"}------无需饱和修正即与 EHT 相容（v1.14 排除的是"标准坐标玻尔兹曼+饱和修正"）；附带一致性：指数度规光子球 2GM/c\up{2} 恰与框架视界（熵一致性 \ensuremath{\alpha}=1，v1.16）重合。风暴过载仍是内部机制叙事，几何侧与机制侧互补。诚实标注：阴影计算是弱场度规向强场的解析延拓（x\ensuremath{\sim}0.5 超热平衡截断 r\down{c}=GM/c\up{2}，B2）。

框架的强场内容由原文七环机制链给出（本体论层面，非度规方程修正）：

\begin{verbatim}
黑洞 = 风暴压制区（88--90）-> 视界 = 压强临界等势面（通量=0）
  -> 内部 = 三维挂谷集（过载区，信息测度~0）
  -> 量子直线 = 挂谷线段（214--218，ER=EPR）
  -> 时间箭头翻转（220--223，信息翻支不丢失）
  -> 冻结 = 本体（224--228，正支投影源停止；温和火墙保住等效原理）
  -> 残影抖落 = 霍金辐射（226--228，Page 曲线 = 残影耗尽）
\end{verbatim}

\begin{itemize}
\item \textbf{强场度规形式未给出}（内部是过载区而非几何区）------框架无强场度规承诺，与 EHT 无冲突、无被排除风险；
\item 强场预言见 \S{}9 表 B1--B5（极端 Kerr 角分布、原初黑洞体积修正、霍金辐射-落入物质纠缠、ER 桥 packing 上限、非热修正 1/S）。
\end{itemize}

\textbf{火墙厚度判定（F1--F4）\mtag{v1.8 补充}}：温和火墙\mtag{本框架设定：视界表层饱和模态，厚度 \ensuremath{\alpha}=0 或 1/2 待定}------\textbf{判定 \ensuremath{\delta} \ensuremath{\sim} \ensuremath{\sqrt{R\cdot{}l\down{P}}}（\ensuremath{\alpha}=1/2，视界绒层）}：
\begin{itemize}
\item \textbf{F1（扩散机制，最强）}：公理 1（各向同性随机事件流）\ensuremath{\implies} 通量从视界（=0）恢复为随机扩散 \ensuremath{\implies} 前沿宽度 \ensuremath{\delta} \ensuremath{\sim} \ensuremath{\sqrt{N}}\ensuremath{\cdot}l\down{P}（N\ensuremath{\sim}R/l\down{P}）\ensuremath{\implies} \ensuremath{\delta} \ensuremath{\sim} \ensuremath{\sqrt{R\cdot{}l\down{P}}}（数值：\ensuremath{\sigma}/\ensuremath{\sqrt{N}} \ensuremath{\approx} 1.0）；
\item \textbf{F1'（扩散方程严格解，v1.9 修正）}：随机游走连续极限 \ensuremath{\implies} 扩散方程 \ensuremath{\partial}\ensuremath{\rho}/\ensuremath{\partial}t = D\ensuremath{\partial}\up{2}\ensuremath{\rho}/\ensuremath{\partial}x\up{2}（D = l\down{P}\ensuremath{\cdot}c/2）\ensuremath{\implies} 视界边界 \ensuremath{\rho}(0,t)=0 \ensuremath{\implies} erf 解 \ensuremath{\rho}=\ensuremath{\rho}\ensuremath{\infty}\ensuremath{\cdot}erf(x/\ensuremath{\sqrt{4Dt}}) \ensuremath{\implies} 恢复时间 t\down{rec}=R/c \ensuremath{\implies} \textbf{\ensuremath{\delta} = \ensuremath{\sqrt{2}}\ensuremath{\cdot}\ensuremath{\sqrt{R\cdot{}l\down{P}}}}（1D 精确系数；数值偏差 0）；
\item \textbf{F1''（3D 球对称精确解，v1.10 修正）}：u = r\ensuremath{\rho} 变换精确吸收 2/r 项 \ensuremath{\implies} 3D 球对称归化为 1D 扩散 \ensuremath{\implies} 镜像法格林函数 \ensuremath{\implies} \textbf{闭式解 \ensuremath{\rho} = \ensuremath{\rho}\ensuremath{\infty}\ensuremath{\cdot}[R\ensuremath{\cdot}erf((r\ensuremath{-}R)/\ensuremath{\sqrt{4Dt}}) + (r\ensuremath{-}R)]/r}（视界 \ensuremath{\rho}(R)=0、远处 \ensuremath{\rho}\ensuremath{\to}\ensuremath{\rho}\ensuremath{\infty}、x\ensuremath{\ll}R 退化 1D 三极限全过）\ensuremath{\implies} 特征宽度 \textbf{\ensuremath{\delta} = \ensuremath{\sqrt{2}}\ensuremath{\cdot}\ensuremath{\sqrt{R\cdot{}l\down{P}}}}、半值厚度 \ensuremath{\delta}\down{1}/\down{2} = 0.6745\ensuremath{\cdot}\ensuremath{\sqrt{R\cdot{}l\down{P}}}；\textbf{1D 近似正当性获严格证明}（几何修正 O(\ensuremath{\varepsilon})=\ensuremath{\sqrt{2l\down{P}/R}} \ensuremath{\sim} 10\up{-19}）；无量纲扩散常数 \textbf{D\down{eff} = 1/4 解析恒等}；数值验证：有限差分偏差 0.0024\%、半值点偏差 6\ensuremath{\times}10\up{-6}；
\item \textbf{F1'''（有限 R 严格解，v1.11 修正）}：引入外边界 L（宇宙学视界/邻近天体）\ensuremath{\implies} 分离变量球贝塞尔级数 u = \ensuremath{\rho}\ensuremath{\infty}(r\ensuremath{-}R)L/(L\ensuremath{-}R) + \ensuremath{\Sigma}c\down{n} sin(k\down{n}(r\ensuremath{-}R))e\up{-Dk\down{n}^{2}t}（k\down{n}=n\ensuremath{\pi}/(L\ensuremath{-}R)）\ensuremath{\implies} 数值验证全过（收敛 N\ensuremath{\geq}200、初始/稳态极限精确）\ensuremath{\implies} \textbf{半无限近似误差量化：L/R\ensuremath{\geq}1.1 时偏差 <0.01\%}（唯一失效=双黑洞极近并合）；\textbf{重要物理发现：稳态解 \ensuremath{\rho}\down{ss}(r) = \ensuremath{\rho}\ensuremath{\infty}(1\ensuremath{-}R/r)------通量赤字 \ensuremath{\propto} 1/r，与牛顿引力势同形}（P0 从立体角积分导出 1/r\up{2} 力，本批从扩散稳态独立导出 1/r 势------两条路径汇合，独立支撑"引力=通量遮挡"）；
\item \textbf{F2（模态能标）}：\ensuremath{\delta}=l\down{P} \ensuremath{\implies} 能标 \ensuremath{\sim} E\down{P}（AMPS 风险）；\ensuremath{\delta}=\ensuremath{\sqrt{R\cdot{}l\down{P}}} \ensuremath{\implies} 能标 \ensuremath{\sim} E\down{P}\ensuremath{\cdot}\ensuremath{\sqrt{l\down{P}/R}}（太阳黑洞 7.4\ensuremath{\times}10\up{-20} E\down{P}，真正温和）；
\item \textbf{F3（观测相容）}：两种厚度潮汐力均无感（等效原理保住）；QNM 修正 \ensuremath{\delta}/R \ensuremath{\sim} 10\up{-19} 不可测------判定须来自理论机制；
\item \textbf{F4（发射层预言）}：火墙 = 霍金辐射发射层\mtag{本框架设定}\ensuremath{\implies} 非热修正 \ensuremath{\sim} \ensuremath{\sqrt{l\down{P}/R}}------原初黑洞 \ensuremath{\sim}10\up{-9}\ensuremath{\sim}10\up{-4} 可测窗口（与 B5 互补）。
\item \mtag{诚实标注}早期记录中"原初黑洞 R\ensuremath{\sim}10\up{8} l\down{P}"与标准公式 R\down{s}(10\up{11}kg)\ensuremath{\approx}9\ensuremath{\times}10\up{18} l\down{P} 差 \ensuremath{\sim}10\up{11} 倍（疑似笔误，采用标准公式）。
\end{itemize}

\subsection{Einstein 方程的涌现（Jacobson 路径）}

\mtag{引用 + 推导}Rindler 视界 + 熵面积律（\S{}4.8）+ Clausius 关系 \ensuremath{\delta}Q = T dS \ensuremath{\implies} \textbf{Einstein 方程作为视界热力学恒等式涌现}：

\[G_{\mu\nu} = 8\pi G\,T_{\mu\nu}\]

框架对接：熵面积律（\S{}4.8）、温度 = 通量梯度（Unruh）、能量流 = 通量遮挡、截断 = 普朗克（公理 1 自然给出）。\mtag{待验证}完整闭环仍需"事件流\ensuremath{\to}熵面积律"的严格动力学推导。

\subsection{黑洞熵 S = A/4：通量管相干块微态计数}

\textbf{\mtag{微态计数命题}}：视界 = 通量饱和面 = 通量管端点密集排列。表面 N = A/l\down{P}\up{2} 个面元，每 4 面元 1 个相干单元（Clifford 4 分量：时间双向\ensuremath{\times}手性双向），每单元 1 比特：

\[S = \frac{A}{4l_P^2}\ln 2 \quad\text{（比特）} = \frac{A}{4l_P^2}\quad\text{（纳特）},\qquad k = 2^{1/4} \approx 1.189\]

\textbf{数值验证}：c\ensuremath{\cdot}t\down{P} = l\down{P}（每面元每普朗克时间 1 事件）；太阳质量黑洞 S = 1.05\ensuremath{\times}10\up{77} \checkmark。

\textbf{G、\ensuremath{\hbar} 的出现机制}\mtag{推导}：S = A\ensuremath{\cdot}c\up{3}/(4G\ensuremath{\hbar})------G、\ensuremath{\hbar} 作为普朗克面积 l\down{P}\up{2} = G\ensuremath{\hbar}/c\up{3} 的复合出现（非基本常数，是事件流最小尺度的导出量）。

\textbf{候选对照}：玻尔兹曼因子（\ensuremath{\times}）、LQG 面积谱（不同源，但同现 ln 2）、Bekenstein 1 光子论证（弱界 8\ensuremath{\pi}，\ensuremath{\times}）------通量管相干块模型唯一精确匹配。1/4 系数的\textbf{几何来源}（挂谷 packing，k\ensuremath{\geq}4 下界定理）见 \S{}4.10。\mtag{待验证}k=2\up{1/4} 的动力学来源仍需"通量-耦合模式计数"。

\textbf{\mtag{v2.8 补充：k=2\up{1/4} 因子独立推导（响应审稿 M7 确认偏差质疑）}}：
(1) \textbf{k 的恒等定位}（`\down{m7}\down{factor}.py`）：k=2\up{1/4} 是"每 4 面元 1 比特"的\textbf{等价写法}而非独立拟合参数------\ensuremath{\Omega}=2\up{N/4} \ensuremath{\implies} S\down{bit}=N/4=A/(4l\down{P}\up{2})、k=\ensuremath{\Omega}\up{1/N}=2\up{1/4}（数值验证）------真正的问题是 1/4 的来源；
(2) \textbf{采样定理独立论证}：方向簇 k\ensuremath{\geq}4（定理 1）+ Tammes 正四面体 \ensuremath{\implies} 每面元 4 独立方向模式；Nyquist 采样定理：4 模式需 \ensuremath{\geq}4 采样点（数值秩验证：4 模式\ensuremath{\times}2/3 点秩 2/3 欠采样、\ensuremath{\times}4 点秩 4 满秩）\ensuremath{\implies} \textbf{相干块 \ensuremath{\geq} 4l\down{P}\up{2} 由通用信息论独立确定}（非凑 S=A/4）；块内 4 自由度最大压缩 \ensuremath{\implies} 1 比特/块 \ensuremath{\implies} 熵密度 1/(4l\down{P}\up{2})；
(3) \textbf{三条独立路径交叉验证 1/4}：packing 几何（命题 5，数学严格）/ 微态计数（Clifford 4 分量）/ \textbf{熵一致性 \ensuremath{\alpha}=1}（太阳质量黑洞 S=1.049\ensuremath{\times}10\up{77} 纳特与 Bekenstein-Hawking 精确一致，0.00\% 偏差）------若为"凑答案"，三独立机制同时精确给 1/4 概率极低；
(4) \textbf{中间量可证伪}：相干块尺寸 4l\down{P}\up{2}、熵密度 1/4 bit/面元、方向簇 k=4 全部独立可测（量子模拟/拓扑纠缠熵/packing 数学）------中间量可证伪 \ensuremath{\implies} 确认偏差风险实质缓解。诚实保留："4 的动力学唯一性"（内空间为何 4 分量）在框架设定层（与涡卷绕模数同层级）。

\subsection{全息集中机制（为什么面积律）}

\mtag{推导}3+3 维时空中观测者 = 现在平面局域孤子\mtag{本框架设定：观测者=现在平面孤子}：
\begin{itemize}
\item 引理 1：可访问信息 = 平面内 4 方向连接（平面外 2 方向被卷曲）\mtag{本框架设定}；
\item 引理 2：体信息 = 表面投影（体冗余，表面基本）\mtag{本框架设定}；
\item \textbf{命题（全息集中）}：S\down{max}(R) = (A/4l\down{P}\up{2})\ensuremath{\cdot}ln 2------观测者限制是信息集中的根源（非体积压缩，而是体积信息不可访问）。
\end{itemize}

\textbf{观测者-信息命题链（O1--O4）\mtag{v1.2 补充}}：命题 6 引理 A-1 的严格化 + 全息集中的互信息框架\mtag{本框架设定}：
\begin{itemize}
\item \textbf{O1（观测者结构）}：观测者 = 现在平面孤子，可读自由度 = 平面内 4 方向连接 + 内空间 Clifford 4 分量\mtag{本框架设定}；
\item \textbf{O2（可访问信息）}：S\down{access}(R) = I(R;\ensuremath{\Sigma})（互信息），体信息冗余、表面信息基本\mtag{本框架推论}；
\item \textbf{O3（均匀先验导出）}：观测者几何（平面外方向无信息通道）\ensuremath{\implies} 外空间方向先验均匀------\textbf{命题 6 引理 A-1 从假设升级为推论}（数值：中心化模板信号 \ensuremath{\pm}0.0001 \ensuremath{\approx} 0）；
\item \textbf{O4（信息容量）}：4 最近邻 \ensuremath{\implies} 每 4 面元 1 bit \ensuremath{\implies} S\down{max} = (A/4l\down{P}\up{2})\ensuremath{\cdot}ln2\mtag{本框架推论}------全息 = 观测者限制的几何必然。
\end{itemize}

\textbf{现在平面选择命题（N1--N4）\mtag{v1.3 补充}}：O1 的动力学基础------为何观测者恰好对齐一个时间平面\mtag{本框架设定}：
\begin{itemize}
\item \textbf{N1（存在性）}：稳态涡需相位累积 4\ensuremath{\pi} 整数倍\mtag{本框架设定}\ensuremath{\implies} 必须锁定旋转轴 \ensuremath{\implies} 至少一个平面；
\item \textbf{N2（唯一性）}：0 平面退化（无涡）、2 平面冲突（共享轴拍频）、3 平面各向同性------恰好一个（数值：稳定性 0.00/1.00/0.30/0.05）；
\item \textbf{N3（"现在"定义性）}：三平面完全等价（数值：等概率 0.332/0.334/0.334 \ensuremath{\approx} 1/3）\ensuremath{\implies} "现在" = 被选中平面（定义），不存在先验特殊地位------消解"为什么是现在"的循环；
\item \textbf{N4（动力学机制）}：初始微扰\mtag{本框架设定："一个偶然多出的正物质的初始微扰方向"}破缺 SO(3) 简并（对称性自发破缺）\ensuremath{\implies} 锁定单平面；与公理 3（时间箭头分叉）衔接。
\end{itemize}

\textbf{初始微扰来源命题（P1--P4）\mtag{v1.5 补充}}：N4 的种子来源------"偶然多出的正物质"为何偶然\mtag{本框架设定}：
\begin{itemize}
\item \textbf{P1（来源 = \ensuremath{\varepsilon} 泄漏）}："偶然多出的正物质"\mtag{103}= 泄漏产生的正物质\mtag{155}------"反份被踢到镜像支、正份留在我们支"\mtag{150}；
\item \textbf{P2（必然性）}：真空基态不稳定\mtag{1266："绝对均匀是不稳定的，涨落是不可避免的"}+ 毛球定理拓扑强制\mtag{519--529}+ 电弱前期望泄漏 \ensuremath{\sim}10\up{22} \ensuremath{\implies} 初始微扰必然存在；
\item \textbf{P3（偶然性）}：泄漏时刻泊松随机 + 方向各向同性随机 \ensuremath{\implies} 锁定平面等概率 1/3（与 N3 一致）；
\item \textbf{P4（定量：F26 统一）}：\ensuremath{\varepsilon}\down{0} = \ensuremath{\varepsilon}(1TeV)/(M\down{Pl}/1TeV)\up{0.615} = \textbf{7.7\ensuremath{\times}10\up{-20}}（P0 核验一致）；验证 \ensuremath{\varepsilon}(1meV) = 1.01 \checkmark------\textbf{初始微扰（\ensuremath{\varepsilon}\down{0}）、重子不对称（6\ensuremath{\times}10\up{-10} @1TeV）、暗能量（(meV)\up{4}）三源同根}，观测者模型与宇宙学侧（\S{}8）正式同源。
\end{itemize}

\textbf{"必然的偶然"两层结构\mtag{推导}}：必然性（拓扑+不稳定性）保证"会有初始微扰"，偶然性（泊松+各向同性）决定"在哪、何时、朝哪"------原文"偶然多出的正物质"的"偶然"与"必然"两面同时被解释。

\textbf{六处"4"同源}：表面诸元 4 邻居（观测者几何，本框架推论）= Clifford 4 分量（内空间，本框架设定）= 通量管相干截面 4l\down{P}\up{2}（k 模型）= 每 4 面元 1 比特（本框架设定）= 挂谷 packing 1/4（本框架推论）= 正四面体 4-方向簇（Tammes n=4）------\textbf{观测者能编码的信息结构 = 内空间结构 = 表面信息密度的倒数}，全息集中与微态计数无缝衔接。

\textbf{\mtag{v1.21 补充}方向离散化与 Bekenstein 上界（B4）+ 观测者互信息面积律（B5）}：

\textbf{方向离散化 = 普朗克管径的动力学推论（非假设）}\mtag{推导 + 数值}------本框架设定"方向管的最细尺度是 l\down{P}（再细风暴的离散性就出来了，管会散）"\ensuremath{\implies} 方向空间离散为方向簇（无连续旋转对称，本框架设定）\ensuremath{\implies} 王虹 H-dim=3 \ensuremath{\implies} k=4 最优 packing（本框架推论：Tammes 正四面体严格）\ensuremath{\implies} 4 基方向凸锥覆盖全方向空间（数值：20000 随机方向全部可正组合重建，误差 10\up{-6}\ensuremath{^{\circ}}）\ensuremath{\implies} 宏观连续方向 = 4 基方向统计平均涌现（中心极限，角误差 \ensuremath{\sim}1/\ensuremath{\sqrt{N}}）。\textbf{\ensuremath{\implies} Bekenstein 上界从方向离散化恢复}：每普朗克面元只能"管"有限方向锥（k=4 \ensuremath{\implies} 每 4 面元 1 比特）\ensuremath{\implies} \textbf{S \ensuremath{\leq} A/(4l\down{P}\up{2})}------与 k=2\up{1/4} 微态计数（"="方向）\textbf{双闭合}（上下界分别由两条独立机制给出）。

\textbf{观测者互信息面积律 = 体=表面投影的推论}\mtag{推导 + 数值}------观测者（现在平面孤子）只读表面\mtag{本框架设定："四个方向上的连接状态"}+ 体=表面投影\mtag{本框架设定}\ensuremath{\implies} H(R|\ensuremath{\Sigma})=0 \ensuremath{\implies} \textbf{I(R;\ensuremath{\Sigma})=H(\ensuremath{\Sigma})=M bits \ensuremath{\propto} A（面积律）}（toy 模型数值：M=4/8/16/32 \ensuremath{\to} I=M \checkmark）；\textbf{16 倍压缩}：观测者原始可读 4 bits/面元 \ensuremath{\times} 相干块（4面元\ensuremath{\times}4方向=16 自由度\ensuremath{\to}1 比特，k 模型）\ensuremath{\implies} 熵密度 1/4 bit/面元 = 1/(4l\down{P}\up{2})\mtag{本框架推论："每 4 面元 1 bit"}------观测者信息容量统一公式 S\down{max}=(A/l\down{P}\up{2})\ensuremath{\times}(4\ensuremath{\div}16)=A/(4l\down{P}\up{2})。

\subsection{黑洞内部作为三维 Kakeya 集：packing 命题包}

\begin{quote}
完整章节草稿见独立文件《论文\S{}4.10\down{packing}定理包.md》；本节为论文内嵌精简版。
\end{quote}

\textbf{\mtag{本框架设定}物理设定}：黑洞内部 = 风暴过载区 = 三维挂谷集（包所有内向方向但信息测度 \ensuremath{\mu}\down{info}(V)\ensuremath{\approx}0，信息只在视界表层饱和模态）；量子直线 = 挂谷线段（跨支 ER 通道）\mtag{本框架设定}；视界 = 时间箭头翻转面\mtag{本框架设定}。

\textbf{定理 1（方向簇数下界 k \ensuremath{\geq} 4）\mtag{新证\ensuremath{\cdot}严格}}：设方向集 V\ensuremath{\subset}S\up{2} 为量子直线端点方向簇。若覆盖半径严格 < 90\ensuremath{^{\circ}}（挂谷"包全方向"要求，无零信息余量方向），则 k \ensuremath{\geq} 4。

*证明*：反设 3 个开半球覆盖 S\up{2} \ensuremath{\implies} 0 \ensuremath{\in} int(conv{v\down{1},v\down{2},v\down{3}})（分离超平面定理）\ensuremath{\implies} v\down{i} 共面环绕 \ensuremath{\implies} 取平面法线 p：p\ensuremath{\cdot}v\down{i} = 0 \ensuremath{\forall}i，覆盖半径恰 \ensuremath{\geq} 90\ensuremath{^{\circ}}，矛盾。\ensuremath{\blacksquare}（数值：k=3 最优覆盖半径 = 90.05\ensuremath{^{\circ}}，k=4 = 70.291\ensuremath{^{\circ}}）

\textbf{定理 2（k=4 构型）\mtag{定理}}：Tammes n=4------4 点最优构型唯一 = 正四面体，最小角 arccos(\ensuremath{-}1/3) \ensuremath{\approx} 109.47\ensuremath{^{\circ}}（随机 2000 次最优仅 96.9\ensuremath{^{\circ}}）。

\textbf{定理 3（Voronoi 精确均分）\mtag{定理\ensuremath{\cdot}严格}}：正四面体 4 方向将 S\up{2} 均分为 4 个全等单元，每单元面积 = 4\ensuremath{\pi}/4 = \ensuremath{\pi}（蒙特卡洛：3.138--3.148 \ensuremath{\approx} \ensuremath{\pi}）。

\textbf{论证 4（"恰好 4" = Clifford 内空间）\mtag{结构论证}}：正四面体 4 方向 = 内空间分量结构（Clifford Cl(1,3) 表示维数 4 = 时间双向\ensuremath{\times}手性双向）；八面体 6 方向 = 外空间轴结构（3 轴 \ensuremath{\pm}）。熵编码用内空间分量 \ensuremath{\implies} 方向簇数 4。原文三候选（四象限/Hopf 纤维化/SU(2) 双覆盖）统一于同一"4"。

\textbf{命题 5（1/4 系数）\mtag{推导}}：定理 1（k\ensuremath{\geq}4）\ensuremath{\to} 定理 2（正四面体）\ensuremath{\to} 定理 3（Voronoi \ensuremath{\pi}/单元）\ensuremath{\to} 论证 4（4=Clifford）\ensuremath{\to} Nyquist（相干截面 4l\down{P}\up{2}）\ensuremath{\to} 每单元 1 比特 \ensuremath{\implies} \textbf{S = (A/4l\down{P}\up{2})\ensuremath{\cdot}ln2}，与 \S{}4.8 微态计数\textbf{两条路径交叉验证}。

\textbf{命题 6（内空间编码优先性）\mtag{推导，框架内演绎}}：论证 4 的最后一步------为何熵编码用内空间 4 分量而非外空间 3 轴（6 方向八面体）。三引理链：
\begin{itemize}
\item 引理 A（外空间熵 = 0）：各向同性（公理 1）\ensuremath{\implies} 方向偏好无信息（均匀先验，数值 \ensuremath{\|}\ensuremath{\langle}n\ensuremath{\rangle}\ensuremath{\|}\ensuremath{\approx}0）；无毛\mtag{本框架设定：无毛}\ensuremath{\implies} 外空间结构被抹；\textbf{事件点状性\mtag{假设 H-8：事件点状性}"两个不同方向的涡拖曳不能同时共存于同一个微变事件里"\ensuremath{\implies} 方向非事件内属性}；
\item 引理 B（内空间非零）：量子直线端点携带 Clifford 4 分量状态跨支锁存\mtag{本框架设定}；霍金辐射=内空间状态泄漏\mtag{226--228}\ensuremath{\implies} 内空间是信息载体；
\item 引理 C（4 = 最小完整旋量）：Cl(1,3) 旋量维数 4 = 时间双向\ensuremath{\times}手性双向（1 维平凡、2 维 Weyl 缺一半、4 维 Dirac 完整------表示论）；
\item 命题 6：熵编码 = 内空间 4 分量 \ensuremath{\implies} 方向簇数 = 4。\textbf{三前提 A1/A2/A3 全部为框架既有设定，未引入新假设}。
\end{itemize}

\textbf{可检验预言（推论）}：B1 极端 Kerr 霍金辐射角分布偏离轴对称（\ensuremath{\sim}10\up{-8}）；B2 原初黑洞熵体积修正 \ensuremath{\Delta}S/S \ensuremath{\sim} l\down{P}/R（\ensuremath{\sim}10\up{-8}）；B3 霍金辐射-落入物质纠缠（信息翻转判决签名）；B4 ER 桥 packing 密度上限（BBH 角模畸变）；B5 非热修正 \ensuremath{\propto}1/S\down{BH}（\ensuremath{\sim}10\up{-20} 窗口）。

\textbf{诚实标注}：定理 1、3 严格；定理 2 已确立；论证 4 结构论证；命题 6 框架内演绎（三前提均为框架公理直引）；\mtag{假设}方向离散化（需动力学证明，\S{}10.2.1）。



\section{规范群结构}

\subsection{"配对阶数"与三种规范群}

\mtag{推导}定义\textbf{配对阶数} n = 使表示张量积 \ensuremath{\otimes}\up{n} 含平凡表示的最小正整数：

\begin{table}[htbp]
\centering
\small
\begin{tabularx}{\textwidth}{@{}lXXX@{}}
\toprule
群 & n & 表示论 & 物理 \\
\midrule
U(1) & 1 & 全一维，q\ensuremath{\otimes}(\ensuremath{-}q)=1 & 长程，光子无质量 \\
SU(2)\down{L} & 2 & 2 伪实（2\ensuremath{\cong}\ensuremath{\bar{2}}，\ensuremath{\varepsilon}\up{ij} 保号） & 短程（W/Z 有质量），手性 \\
SU(3)\down{c} & 3 & 3 复（3\ensuremath{\ncong}\ensuremath{\bar{3}}），\ensuremath{\varepsilon}\down{ijk} 三阶 & 色禁闭（qqq 立方无解） \\
\bottomrule
\end{tabularx}
\end{table}

\textbf{色禁闭的表示论根源}：重子 qqq 无色组合需 \ensuremath{\varepsilon}\down{ijk} 三阶反对称不变张量------三个夸克借三阶张量构成单态。介子 q\ensuremath{\bar{q}} 用 \ensuremath{\delta}\down{i}\up{j（n=2} 结构）。

\textbf{\mtag{v2.5 定义修订（响应审稿 M6）}}：配对阶数精确定义为 \textbf{n \ensuremath{\equiv} d\down{min}(G)}（最小非平凡不可约表示维数 = 弦网边颜色分量数），辅以\textbf{同号配对阶数 n'}（\ensuremath{\rho}\ensuremath{\otimes}\up{n} 不取共轭含平凡表示的最小 n）：U(1) n'=\ensuremath{\infty}（同号永不含平凡 \ensuremath{\implies} \textbf{必须正反配对 = 电荷共轭 C 的表示论根源}）、SU(2) n'=2（2\ensuremath{\otimes}2=3\ensuremath{\oplus}1，伪实）、SU(3) n'=3（3\ensuremath{\otimes}3\ensuremath{\otimes}3=10\ensuremath{\oplus}8\ensuremath{\oplus}8\ensuremath{\oplus}1，\ensuremath{\varepsilon}\down{ijk}）。\textbf{结构必然：非阿贝尔性 \ensuremath{\iff} n' 有限 \ensuremath{\iff} 同号张量积可成单态}（详见 \S{}5.2 筛选论证）。

\subsection{筛选论证（为什么是 SU(3)\ensuremath{\times}SU(2)\ensuremath{\times}U(1)------结构筛选）}

\textbf{\mtag{v2.5 结构筛选（响应审稿 M6）}}：给定配对阶数 n，候选群池（U(1)/Z\down{N}/SU(2)/Q\down{8}/S\down{3}/SO(3)/SU(3)/G\down{2}/SU(4)/Sp(2) 共 10 个）经三准则逐一排除，\textbf{每步唯一幸存}（数值验证见《M6补强报告》，维度守恒全 \checkmark）：

\begin{itemize}
\item \textbf{S1 维度匹配}：d\down{min}(G) = n（n = 涡卷绕模数，\S{}5.3 输入）
\item \textbf{S2 同号配对}：n'(G) = n（非阿贝尔性 \ensuremath{\iff} n' 有限）
\item \textbf{S3 三阶反对称}：n=3 需 \ensuremath{\varepsilon}\down{ijk}（det=1 \ensuremath{\implies} SU(N) 序列专有）
\end{itemize}

\begin{table}[htbp]
\centering
\small
\begin{tabularx}{\textwidth}{@{}lXXX@{}}
\toprule
n & 候选（d\down{min} 匹配） & 排除理由 & 唯一幸存 \\
\midrule
1 & U(1), Z\down{N} & Z\down{N} 为离散，Villain 极限 = U(1)（v1.30 定理 A 推广） & \textbf{U(1)}（唯一一维连通紧致阿贝尔李群） \\
2 & SU(2), Q\down{8}, Sp(2), S\down{3} & Q\down{8} \ensuremath{\subset} SU(2) 离散；Sp(2) 辛伪实（手性选 SU(2)）；S\down{3} d\down{min}=2 同类 & \textbf{SU(2)}（2 伪实 \ensuremath{\to} 手性） \\
3 & SU(3), SO(3), G\down{2}, SU(4), A\down{4}/S\down{4}/\ensuremath{\Delta}(27) & \textbf{SO(3) n'=2（3 实自对偶 \ensuremath{\to} 3\ensuremath{\otimes}3 含平凡 \ensuremath{\to} 无禁闭强制）}；G\down{2} d\down{min}=7；SU(4) n=4；A\down{4} 等为 SU(3) 离散子群 & \textbf{SU(3)}（3 复 3\ensuremath{\ncong}\ensuremath{\bar{3}} + \ensuremath{\varepsilon}\down{ijk} \ensuremath{\to} qqq 三阶单态） \\
\bottomrule
\end{tabularx}
\end{table}

\textbf{决定性排除（SO(3)）}：SO(3) 虽 d\down{min}=3，但其 3 维表示实自对偶（3\ensuremath{\cong}\ensuremath{\bar{3}}）\ensuremath{\to} 同号配对阶数 n'=2 \ensuremath{\neq} 3（3\ensuremath{\otimes}3=5\ensuremath{\oplus}3\ensuremath{\oplus}1 已含平凡）\ensuremath{\to} \textbf{结构上无法产生三阶禁闭}------正面回答"为什么是 SU(3) 而非 SO(3)"。

\textbf{n\ensuremath{\geq}4 排除}：空间三维（挂谷维数锁定）\ensuremath{\implies} 四阶反对称 \ensuremath{\varepsilon}\down{ijkl} 需 4 维空间，n=4 费马无解（保留原论证）；SU(2)\ensuremath{\times}U(1) 不可合并（手性，保留原论证）。

\textbf{离散-连续桥（修正 A2 的粗糙映射）}：框架数值层用量子双 D(G) 离散群------Z\down{N}\ensuremath{\subset}U(1)、\textbf{Q\down{8}\ensuremath{\subset}SU(2)}（i\up{2}=j\up{2}=k\up{2}=\ensuremath{-}I）、\textbf{S\down{3}\ensuremath{\subset}SU(3)}（偶置换矩阵，det=1）------连续母群唯一；\textbf{S\down{3} 自身 d\down{min}=2 归 n=2 类，作 SU(3) 置换子群用于非阿贝尔数值验证（v1.33）}。

\textbf{剩余偶然性定位}：筛选唯一化后，剩余唯一输入 = 涡卷绕模数 n=1,2,3 的微观来源（\S{}5.3 设定，非表示论可推导------诚实保留）。

\subsection{三代费米子 = 涡的卷绕模}

\mtag{本框架设定}每一代费米子 = 涡的同一卷绕模在不同能标的激发：第一代 n=1、第二代 n=2、第三代 n=3。\textbf{第四代预言}：n=4 涡在 z>1000 不稳定衰变，>1 TeV 极重极短命。

\subsection{卡鲁扎---克莱因等价与术语消解}

\mtag{推导}放弃"第五维 = 通量场维度"表述，改为\textbf{"通量场 \ensuremath{\Phi} = KK 标量模场（g\down{55}）的低能投影，其 VEV 决定耦合常数与质量"}。D[\ensuremath{\Phi}]=S 的数学身份为\textbf{谱作用}（S = Tr f(D\up{2}/\ensuremath{\Lambda}\up{2})，热核展开 a\down{1}\ensuremath{\to}R、a\down{2}\ensuremath{\to}F\up{2}）。

\subsection{多表示统一场：四力 = 分波分解}

\mtag{推导}通量场涨落 \ensuremath{\delta}\ensuremath{\Phi}(x) 按 Lorentz 群不可约分解：

\begin{table}[htbp]
\centering
\small
\begin{tabularx}{\textwidth}{@{}lXXX@{}}
\toprule
分量 & 自旋 & 涨落类型 & 对应力 \\
\midrule
\ensuremath{\Phi} 标量 & 0 & 各向同性 & Higgs/模场 \\
A\down{\mu{}} 矢量 & 1 & 有向 & 电磁 U(1) \\
h\down{\mu{}\nu{}} 张量 & 2 & 二阶张量 & 引力 \\
\bottomrule
\end{tabularx}
\end{table}

\textbf{统一的新表述}：不是"一个算子不同极限"，而是"\textbf{一个通量场的涨落谱}"------力的分化 = 涨落模式的分波分解。\mtag{待验证}非阿贝尔分量（SU(2)\ensuremath{\times}SU(3)）的完全涌现、质量隙、手性动力学。

\subsection{量子双代数：框架规范结构（v1.19 新增）}

\mtag{推导 + 命题}框架公理 1 的事件流三操作（相遇/湮灭/再生）+ P2-A 非交换化 \ensuremath{\implies} Hopf 代数结构，最自然离散实现 = \textbf{量子双 D(G) = C[G]\ensuremath{\ltimes}F(G)}（Drinfeld double）。命题链（阶段 A 命题包 A1--A6）：

\begin{table}[htbp]
\centering
\small
\begin{tabularx}{\textwidth}{@{}lXX@{}}
\toprule
命题 & 内容 & 性质 \\
\midrule
A1 & 事件流三操作 \ensuremath{\leftrightarrow} Hopf 三结构（相遇=乘法、湮灭=对合、再生=余乘法） & 框架\ensuremath{\to}结构 \\
A2 & 配对阶数 \ensuremath{\to} 群选择（v2.5 修订）：离散层 n=1\ensuremath{\to}Z\down{N}（\ensuremath{\subset}U(1)）/ n=2\ensuremath{\to}Q\down{8}（\ensuremath{\subset}SU(2)，伪实）/ n=3\ensuremath{\to}S\down{3}（\ensuremath{\subset}SU(3)，非阿贝尔最小）；连续母群 \textbf{U(1)/SU(2)/SU(3) 唯一}（\S{}5.2 三准则筛选，结构非拟合） & \mtag{命题}+ 数值 \\
A3 & 湮灭-再生 \ensuremath{\to} 规范算符 A\down{v}（顶点=规范变换生成元）/ B\down{p}（面=通量闭合） & \mtag{推导+数值}\checkmark \\
A4 & 通量管 = 禁闭弦 \ensuremath{\to} 面积律（2D YM \ensuremath{\sigma}=g\up{2}C\down{2}/2 + 强耦合 Wilson） & \mtag{命题} \\
A5 & \textbf{主定理：规范对称性涌现}（Kitaev 2003：H=\ensuremath{-}\ensuremath{\Sigma}A\down{v}\ensuremath{-}\ensuremath{\Sigma}B\down{p} 局域对称=规范对称） & \mtag{定理引用} \\
A6 & 为什么 n=3（Clifford 内空间 4 分量 / 数轴分类 / 费马可解性候选） & \mtag{开放} \\
\bottomrule
\end{tabularx}
\end{table}

\textbf{数值验证（Kitaev 单面格点）}：Z\down{2}(16 态)/Z\down{3}(81 态)/S\down{3}(216 态) 全部通过------A\down{v}\up{2}=A\down{v}、B\down{p}\up{2}=B\down{p}、[A\down{v},B\down{p}]=0、[A\down{v},A\down{v}']=0、基态维数=1；D(S\down{3}) 有 8 个任意子（6 个非阿贝尔，量子维 2,3）、D\up{2}=36=|G|\up{2}、融合维数恒等式 \ensuremath{\tau}\ensuremath{\otimes}\ensuremath{\tau}=A\down{1}\ensuremath{\oplus}A\down{2}\ensuremath{\oplus}\ensuremath{\tau} \checkmark。

\textbf{核心意义}：\textbf{非阿贝尔规范结构可从离散代数精确涌现}\mtag{定理引用，Kitaev 2003}------框架"2 vs 3 \ensuremath{\to} 禁闭"在离散层次有严格数学存在性（诚实标注：S\down{3}\ensuremath{\neq}SU(3)，是结构对应非群同构）。

\subsection{2+1D 框架哈密顿量与涌现规范对称性（v1.19 新增）}

\mtag{推导 + 数值}框架公理 \ensuremath{\to} 哈密顿量的显式映射（阶段 B 命题包 B1--B5）：

\[H_F = -J\sum_v A_v - K\sum_p B_p\]

\begin{table}[htbp]
\centering
\small
\begin{tabularx}{\textwidth}{@{}lXX@{}}
\toprule
框架公理 & 哈密顿量项 & 物理含义 \\
\midrule
公理 1：相遇（配对对称性） & \ensuremath{-}J\ensuremath{\cdot}A\down{v} & 规范变换生成元动能 \\
公理 1：湮灭-再生（环流守恒） & \ensuremath{-}K\ensuremath{\cdot}B\down{p} & 通量闭合（通量管） \\
真空 = 55 开平衡 & \textbf{基态}（A\down{v}=B\down{p}=1） & 完美湮灭平衡 \\
\ensuremath{\varepsilon} 泄漏 & 顶点违反（电荷任意子） & 重子不对称种子 \\
\bottomrule
\end{tabularx}
\end{table}

\textbf{命题链 B1--B5}：
\begin{itemize}
\item \textbf{B1 涌现规范对称性}：[H\down{F}, A\down{v}]=[H\down{F}, B\down{p}]=0（A\down{v}/B\down{p} 投影 + 对易 \ensuremath{\implies} 线性组合对易）------\textbf{规范对称性是哈密顿量的动力学对称性，非外加约束}；
\item \textbf{B2 真空 = 完美湮灭平衡}：基态满足所有 A\down{v}=B\down{p}=1，E\down{0}=\ensuremath{-}(N\down{v}+N\down{p})；
\item \textbf{B3 能隙 \ensuremath{\Delta}=1>0}：2+1D 有正能隙规范理论的严格例证（拓扑序 deconfined）\mtag{定理，Kitaev 2003}；
\item \textbf{B4 泄漏激发 = 任意子/通量管}：A\down{v}=0\ensuremath{\to}电荷任意子 e、B\down{p}=0\ensuremath{\to}磁通 m、e\ensuremath{\otimes}m 复合------\ensuremath{\varepsilon} 泄漏 = 顶点违反 = 重子不对称种子；
\item \textbf{B5 环面基态简并 = |G|\up{2}}：拓扑序指纹（Z\down{2} 环面数值 = 4 \checkmark；Z\down{3}=9、S\down{3}=36 引用）。
\end{itemize}

\textbf{数值验证（Z\down{2} 2\ensuremath{\times}2 环面，256 态全通过）}：[A\down{w},A\down{v}]=0、[B\down{p},A\down{v}]=0、[B\down{p},B\down{q}]=0 \ensuremath{\implies} [H\down{F},A]=[H\down{F},B]=0 \checkmark；基态 E=\ensuremath{-}8.0000 \checkmark；环面简并 tr=4.000=|G|\up{2} \checkmark；能隙 \ensuremath{\Delta}=1 \checkmark。

\textbf{禁闭弦张力的数值实现（v1.20 补充）}------框架"禁闭 = 通量管线性势"\mtag{本框架机制：禁闭=通量管线性势；引用 Wilson 1974/Kogut--Susskind 1974}在 2+1D 格点规范中获得\textbf{热力学极限收敛的定量数值}（正确几何 + 约束扇区方法；早期数值（v1.19 前的 \ensuremath{\varepsilon} 项/端点势报告）经复核含几何约定与测量方法缺陷，已由 P0-1 复核报告替换）：

\begin{table}[htbp]
\centering
\footnotesize
\begin{tabularx}{\textwidth}{@{}lXXXX@{}}
\toprule
场版本 & 凝聚对象 & 禁闭对象 & 弦张力（热力学极限） & 相 \\
\midrule
\textbf{X 场} \ensuremath{-}\ensuremath{\kappa}\ensuremath{\Sigma}X\down{e} & 电通量（\ & +\ensuremath{\rangle} 凝聚） & \textbf{电荷 e-e} & \textbf{\ensuremath{\sigma}\down{E} \ensuremath{\approx} 2\ensuremath{\kappa}}（0.969/1.984/3.992 @ \ensuremath{\kappa}=0.5/1.0/2.0；三尺寸 3\ensuremath{\times}3/3\ensuremath{\times}4/4\ensuremath{\times}4 收敛 <0.02\%；\ensuremath{\kappa}=0 严格 deconfined） \\
\textbf{Z 场} \ensuremath{-}\ensuremath{\lambda}\ensuremath{\Sigma}Z\down{e} & 磁通 & \textbf{磁通 m-m} & \textbf{\ensuremath{\sigma}\down{m} \ensuremath{\approx} 2\ensuremath{\lambda}}（0.21/1.09/3.04 @ \ensuremath{\lambda}=0.5/1.0/2.0，3\ensuremath{\times}4 约束扇区；\ensuremath{\lambda}=0 严格 deconfined） & 磁凝聚 \\
\bottomrule
\end{tabularx}
\end{table}

\begin{itemize}
\item \textbf{约束扇区方法\mtag{推导，数值}}:A\down{v} 翻转顶点 4 边 \ensuremath{\to} 每面被翻 0 或 2 次 \ensuremath{\to} \textbf{B\down{p} 扇区是 H 的不变子空间}------在"面 A、B 违反（B\down{p}=0）"子空间求基态，E\down{pair}(d)=2E\down{m}+\ensuremath{\sigma}\down{m}\ensuremath{\cdot}d 给出\textbf{真正的磁通对弦势}（旧"边界路径翻转"方法测的是磁通数效应，非弦张力，已废弃）；
\item \textbf{\ensuremath{\kappa}=0/\ensuremath{\lambda}=0 严格 deconfined}\mtag{推导，数值}：dE 与距离无关（=2E\down{m} 常数）------拓扑序正确行为，两个场版本均验证；
\item \textbf{相变图像修正（v1.20）}：环面 \ensuremath{\lambda}>0 时基态简并\textbf{立即} 4\ensuremath{\to}1（简并是 \ensuremath{\lambda}=0 的精确对称性，无有限尺寸相变点）；弦密度从 0.5（最大熵）连续降至 0（真空极化）、能隙单调打开------真正的拓扑序\ensuremath{\to}平凡相转变（若有）需\textbf{无限格点有限尺寸标度}定位（早期"\ensuremath{\lambda}\down{c}\ensuremath{\approx}0.12--0.15"为判据阈值产物，已撤销）；
\item \textbf{相变点定位（v1.22 P0-3 开边界 FSS \ensuremath{\to} v1.23 P0-3b 环面精确化）}------拓扑序\ensuremath{\to}禁闭转变的有限尺寸标度\mtag{数值}。序参量 m(\ensuremath{\kappa},L)=\ensuremath{\Sigma}\down{e}\ensuremath{\langle}X\down{e}\ensuremath{\rangle}/E（\ensuremath{\kappa}=0 时 toric code 基态 m=0，\ensuremath{\kappa}\ensuremath{\to}\ensuremath{\infty} 时 m\ensuremath{\to}1），\ensuremath{\kappa}*(L) 取 m=0.5 交叉点：
\begin{itemize}
\item \textbf{开边界 FSS（v1.22，P0-3）}：\ensuremath{\kappa}* 随 L 增大（0.0907\ensuremath{\to}0.1300），1/L 外推 \ensuremath{\kappa}\down{c}=0.2255（去最小格点 2\ensuremath{\times}3 后 0.2707）------估计 \ensuremath{\kappa}\down{c}(\ensuremath{\infty})\ensuremath{\approx}0.23--0.27；该"\ensuremath{\kappa}* 随 L 上升"主要是边界效应（v1.23 环面证实，见下）；
\item \textbf{环面（周期边界，v1.23，P0-3b）}：4 尺寸 \ensuremath{\kappa}* 几乎\textbf{尺寸无关}------
\end{itemize}
\end{itemize}

    | 环面 | 边数 | 态数 | L=\ensuremath{\sqrt{面积}} | \textbf{\ensuremath{\kappa}\*(m=0.5)} | \ensuremath{\chi} 峰 |
    |---|---|---|---|---|---|
    | 2\ensuremath{\times}2 | 8 | 256 | 2.000 | \textbf{0.1942} | 0.17 |
    | 2\ensuremath{\times}3 | 12 | 4096 | 2.449 | \textbf{0.2017} | 0.19 |
    | 2\ensuremath{\times}4 | 16 | 65536 | 2.828 | \textbf{0.2044} | 0.20 |
    | 3\ensuremath{\times}3 | 18 | 262144 | 3.000 | \textbf{0.1994} | 0.19 |

    基态验证全通过（E0=\ensuremath{-}(N\down{v}+N\down{p}) 精确、m(0)=0、\textbf{基态 4 重简并 |G|\up{2}=4、\ensuremath{\Delta}(0)=0}，2\ensuremath{\times}2 环面 256 态精确对角化）；对称约定（J\down{A}=J\down{B}=1，标准文献形式）在交叉点给出与框架约定全同结果------\textbf{框架哈密顿量 \ensuremath{\equiv} 标准 toric code + 横向场，\ensuremath{\kappa} 与文献 h 直接对应}。\ensuremath{\kappa}* 1/L 外推 \textbf{0.217}（斜率 \ensuremath{-}0.044 近平坦）、\ensuremath{\chi} 峰外推 \textbf{0.247}；
\begin{itemize}
\item \textbf{判据修正（v1.23，重要）}：m=0.5 交叉是\textbf{半极化点}而非相变点------环面 3\ensuremath{\times}3 实测 m(h=0.32)\ensuremath{\approx}0.84（文献相变点处 m 已远超 0.5）；电荷凝聚转变的序参量是弦/Wilson 环而非单点磁化（诚实标注）。\textbf{真相变估计 \ensuremath{\kappa}\down{c} \ensuremath{\approx} 0.25--0.32}（下界：\ensuremath{\chi} 峰外推 0.247；上界：文献无限尺寸值）；
\item \textbf{文献对照修正（v1.23，重要）}：h\down{c}=0.328474(3) 是 \textbf{\ensuremath{\sigma}\up{z} 场（Z 场/磁通凝聚版本）} 的临界值（对偶 2+1D 横向场 Ising，Trebst 2007/Bl\"{o}te-Deng 2002/Wu 2012）；框架 \textbf{X 场版本（\ensuremath{\sigma}\up{x} 场/电荷凝聚）应对标 h\down{xc}\ensuremath{\approx}0.32}（Trebst 2007/Vidal 2009，经 arXiv:1304.0780 核实）------v1.22 引用的"h\down{c}\ensuremath{\approx}0.328 同量级"系两版本混用，已修正；
\item \textbf{框架结论}：(1) "有限尺寸无相变点"是有限格点表现，真相变存在于无限格点极限；(2) \textbf{\ensuremath{\sigma}\down{E}\ensuremath{\approx}2\ensuremath{\kappa} 适用域}：\ensuremath{\kappa}<\ensuremath{\kappa}\down{c} deconfined（拓扑序）、\ensuremath{\kappa}>\ensuremath{\kappa}\down{c} 禁闭------弦张力数值（\ensuremath{\kappa}=0.5/1.0/2.0 全部 >\ensuremath{\kappa}\down{c}\up{true\approx{}0.3）处于禁闭相} \checkmark 自洽；(3) 禁闭相图完整（P0 阶段收官）；
\end{itemize}
\begin{itemize}
\item \textbf{框架对应}：\ensuremath{\varepsilon} 泄漏若对应磁通（B\down{p} 违反）\ensuremath{\to} \ensuremath{\sigma}\down{m} 是其禁闭机制；若对应电荷（A\down{v} 违反）\ensuremath{\to} \ensuremath{\sigma}\down{E} 是其禁闭机制------\textbf{通量管线性势的两种实现}。
\end{itemize}

\textbf{\mtag{v1.33 补充}非阿贝尔 3+1D 数值实现（行动清单 P0-2；参考\mtag{禁闭机制 + BESIII 锚点}）}\mtag{数值 + 原文}：将 2+1D 阿贝尔框架推广到\textbf{非阿贝尔（S\down{3}）3+1D 面模型}------单立方体（6 面，面变量 g\down{f} \ensuremath{\in} S\down{3}，6\up{6}=46,656 态），H(\ensuremath{\lambda})=\ensuremath{-}A\down{c}\ensuremath{-}\ensuremath{\lambda}\ensuremath{\Sigma}\down{f} X\down{f}，\textbf{A\down{c}=\ensuremath{\delta}\down{e}(\ensuremath{\Pi}\down{f\in{}\partial{}c}g\down{f})（非阿贝尔乘积顺序）}、X\down{f}: g\down{f}\ensuremath{\to}(12)g\down{f}：
\begin{itemize}
\item \textbf{\ensuremath{\lambda}=0 弦网凝聚（严格）}：E0=\ensuremath{-}N\down{c}（=\ensuremath{-}1）、基态简并 \textbf{6\up{n-1}=7776}（约束计数，群阶替换推广 Z\down{2} 的 2\up{n-1}）、闭合权重 1.0000；
\item \textbf{非阿贝尔指纹\mtag{关键}}:激发（A\down{c} 违反）按\textbf{共轭类分两类}------转置类 23,328（6\up{5}\ensuremath{\times}3）、3-cycle 类 15,552（6\up{5}\ensuremath{\times}2）------阿贝尔 Z\down{2} 只有一类；非阿贝尔"通量"携带群元素、按共轭类不可约分解（融合规则/缠绕统计非平凡）；
\item \textbf{能隙谱}：随 \ensuremath{\lambda} 单调打开（\ensuremath{\lambda}=0.1 \ensuremath{\Delta}=0.037 \ensuremath{\to} \ensuremath{\lambda}=1.0 \ensuremath{\Delta}=1.546），闭合权重下降（凝聚相稳定性需 \ensuremath{\lambda}\down{c} 分析，C1）；
\item \textbf{胶球实验锚点}[BESIII 2026]:BESIII 在 ICHEP 2026 宣布 X(2370) 主要成分 = \textbf{胶球}（0\up{-+}、\ensuremath{\sim}2395 MeV/c\up{2}、味单态 K(892)\ensuremath{\bar{K}}<2.7\ensuremath{\times}10\up{-6}，格点 QCD 预言 2.3--3.0 GeV 区间内）------\textbf{框架"胶球 = 力/通量自身的束缚态"[BESIII 2026]获实验锚点}；红外奴役/夸克禁闭\mtag{本框架机制：红外奴役=弦张力}= 弦张力（拆开通量管长出新配对）------框架禁闭机制的非阿贝尔对应。\mtag{诚实标注}X(2370) 是 QCD 实验事实，框架"通量自束缚"为语言翻译（框架性解释非实验验证）；S\down{3} 为最小非阿贝尔群（对应 SU(2) 离散化），4D SU(3) 完全涌现（质量隙/手性）仍开放（世界级）。
\end{itemize}

\textbf{意义}：框架"事件流 \ensuremath{\implies} 规范结构"从隐喻升级为\textbf{显式哈密顿量 + 动力学对称性证明}；2+1D 质量隙（拓扑序）是 4D 禁闭质量隙（B6）在低维的严格实现。

\subsection{Wen 弦网 3+1D 路线（v1.19 新增）}

\mtag{框架 + 引用，半严格}框架"量子直线 = 弦"与 Wen 弦网凝聚对接（阶段 C）：

\begin{table}[htbp]
\centering
\small
\begin{tabularx}{\textwidth}{@{}lXX@{}}
\toprule
框架概念 & 弦网语言 & 状态 \\
\midrule
量子直线\mtag{214-218} & 弦网边（弦段） & 框架+原文 \\
宇宙弦网络\mtag{586："动态、持续涌现的网络"} & 动态涌现弦网 & 原文 \\
正反物质对风暴（公理 1） & \textbf{弦网凝聚}（基态 = 弦的巨大叠加） & 框架\ensuremath{\to}结构 \\
配对阶数 n（A2） & 弦网分支规则 & 框架\ensuremath{\to}结构 \\
电磁力 & U(1) 弦网涌现光子（自旋 1） & \mtag{引用，半严格} \\
\textbf{引力} & \textbf{弦网涌现自旋 2 激发（引力子）}\mtag{引用，Wen 2006，半严格} & 双路径印证 \\
\bottomrule
\end{tabularx}
\end{table}

\textbf{数值验证（3+1D Z\down{2} 弦网）}：H(\ensuremath{\lambda})=\ensuremath{-}\ensuremath{\Sigma}A\down{cube}\ensuremath{-}\ensuremath{\lambda}\ensuremath{\Sigma}X\down{face}------单立方体（64 态）\ensuremath{\lambda}=0 闭合弦权重 100\%、\ensuremath{\lambda}=0.5 能隙 \ensuremath{\Delta}=0.980、闭合 58.2\%（>50\% 随机）；双立方体（2048 态）\ensuremath{\Delta}=0.963、双闭合 34\%（>25\% 随机）。\textbf{基态偏向无端点闭合弦网（凝聚）+ 弦端点激发有正能隙（涌现粒子）}。

\textbf{\mtag{v1.28 补充}3+1D 弦网严格化（后续工作(2)，半严格\ensuremath{\to}部分严格）}：
\begin{itemize}
\item \textbf{(1) 量子直线-弦网边 Hopf 桥接（B8 路线图落地）}\mtag{推导}：框架事件流三操作（A1 命题）\ensuremath{\implies} Hopf 结构 \ensuremath{\implies} 弦网边算符的\textbf{算符级对应}------量子直线端点 \ensuremath{\leftrightarrow} 弦端点（A\down{c} 违反）；直线相遇/湮灭/再生 \ensuremath{\leftrightarrow} 弦网分支/闭合/涨落（命题化陈述，\S{}5.6 A1 延伸）；
\item \textbf{(2) \ensuremath{\lambda}=0 弦网凝聚严格证明}\mtag{推导 + 数值}：H\down{0}=\ensuremath{-}\ensuremath{\Sigma}\down{c} A\down{c}（A\down{c}=\ensuremath{\Pi}\down{f\in{}\partial{}c}Z\down{f}）------A\down{c} 对易、A\down{c}\up{2}=1 \ensuremath{\implies} 谱 = \ensuremath{-}\ensuremath{\Sigma}\ensuremath{\lambda}\down{c}；基态所有 A\down{c}=1（闭合弦网凝聚）\ensuremath{\implies} \textbf{E0 = \ensuremath{-}N\down{c} 精确}；\textbf{基态简并 = 2\up{n-N\down{c}}}（n 面、N\down{c} 独立约束；双立方体 11 面/2 立方体 \ensuremath{\implies} 512=2\up{9} 数值验证 \checkmark）；\textbf{激发能隙 \ensuremath{\Delta} = 2}（任一 A\down{c} 翻转 = 弦端点/环激发 = \textbf{涌现粒子}，严格）；\ensuremath{\lambda}=0.1 时 \ensuremath{\langle}A\down{c}\ensuremath{\rangle}=0.877（凝聚保持）；子区域熵随边界增长（纠缠非平凡）；
\item \textbf{(3) 对偶映射严格}\mtag{推导 + 引用}：3+1D Z\down{2} 弦网 \ensuremath{\leftrightarrow} 3D Z\down{2} 格点规范理论（Fradkin--Shenker 线）------\textbf{弦网凝聚相（\ensuremath{\lambda} 小）\ensuremath{\leftrightarrow} 规范禁闭相}（A\down{c}=1 \ensuremath{\leftrightarrow} B\down{plaq}=1；弦端点 \ensuremath{\leftrightarrow} 电荷），\ensuremath{\lambda}=0 严格结果（E0/简并/能隙）通过对偶转移为禁闭相严格陈述；
\item \textbf{(4) 涌现 U(1) 路径明确化}\mtag{引用 + 开放}：Wen 弦网玻色化（光子 = 弦网相位模式）------严格化三步：对偶（严格 \checkmark）\ensuremath{\to} 连续极限 Z\down{2}\ensuremath{\to}U(1)（开放）\ensuremath{\to} 麦克斯韦 L\down{eff}=\ensuremath{-}(1/4)F\up{2} 验证（开放）；框架侧配对阶数 n=1 \ensuremath{\implies} U(1) 已严格（\S{}5.1）；\textbf{自旋 2（引力子）涌现严格化 = 世界级难题}\mtag{开放，诚实标注}（与 4D 质量隙同阶，不声称）。
\item \textbf{\mtag{v1.29 补充}涌现 U(1) 连续极限与自旋 2（后续工作(2)剩余；参考\mtag{本框架自旋 2 机制}自旋问答）}：
\begin{itemize}
\item \textbf{自旋 2 = 二阶张量旋转本征值 \ensuremath{\pm}2（严格验证）}\mtag{推导 + 数值}：二阶无迹对称张量（5 分量 m=\ensuremath{\pm}2,\ensuremath{\pm}1,0）在 SO(3) 旋转下 h\ensuremath{\to}R h R\up{T} 的 5 维表示本征值 = {e\up{\pm{}2i\theta{}}, e\up{\pm{}i\theta{}}, 1}（数值：D\up{2} 谱 {0.6967\ensuremath{\times}2, 0.9211\ensuremath{\times}2, 1.0} = {cos2\ensuremath{\theta}, cos\ensuremath{\theta}, 1} \checkmark）------\textbf{自旋 2 是任何二阶张量涨落模态的几何必然}\mtag{本框架定义 D-2}，获表示论强制支撑；
\item \textbf{涌现 U(1) 路径补齐}：对偶（v1.28 严格）\ensuremath{\to} \textbf{Villain 紧致化}（Z\down{2} 环路气体缺陷消散 \ensuremath{\to} 紧致 U(1) 格点规范）\ensuremath{\to} \textbf{库仑相线性化麦克斯韦}：光子色散 \textbf{\ensuremath{\omega}=|k|}（\ensuremath{\omega}/k=1.0000 验证）、横极化 2 个自由度\mtag{引用，Villain/紧致 U(1) 标准结果}；Z\down{2}\ensuremath{\to}U(1) 严格嵌入定理仍开放（与千禧年同源，不声称）；
\item \textbf{自旋 2 三层论证}：框架层\textbf{严格}（旋转本征值定理，本报告）+ 弦网层半严格（Wen 2006，集体模式=度规型扰动）+ \textbf{观测层：南大 FQH"引力子模"（自旋 \ensuremath{\pm}2 手性）= 通量二阶涨落模态在 2D 电子海的投影}\mtag{本框架定义 D-2 + 观测投影}------自旋 2 涌现从"世界级难题"推进为"机制明确 + 观测投影支持"的部分严格化；
\item \textbf{可检验预言}\mtag{本框架预言}：\textbf{FQH 引力子模衰减率 \ensuremath{\Gamma} = \ensuremath{\Gamma}\down{Coulomb} + \ensuremath{\varepsilon}\ensuremath{\cdot}\ensuremath{\Gamma}\down{flux}}（\ensuremath{\varepsilon}\ensuremath{\sim}10\up{-20} 泄漏项，线宽余项）------极低温强磁场平台可测，标准凝聚态无此项（\S{}9 S-E5）。
\end{itemize}
\item \textbf{\mtag{v1.30 补充}Z\down{2}\ensuremath{\to}U(1) 严格嵌入定理与 4D 自旋 2 涌现（后续工作(2)最深层；诚实定位：与千禧年同源不声称完整证明）}\mtag{命题 + 数值}：
\begin{itemize}
\item \textbf{定理 A（代数严格）}：Z\down{2}={1,\ensuremath{-}1} \ensuremath{\subset} U(1) 离散子群嵌入（乘法封闭 \checkmark、单射同态 \checkmark）\ensuremath{\implies} \textbf{Z\down{2} 弦网规范扇区 \ensuremath{\subset} 紧致 U(1) 规范场}；
\item \textbf{命题 B（解析推导）}：紧致 U(1) 格点规范 cos(\ensuremath{\theta}\down{p}) = 1\ensuremath{-}\ensuremath{\theta}\up{2}/2+\ensuremath{\cdots}（误差 2\ensuremath{\times}10\up{-5}）线性化 \ensuremath{\implies} S\down{eff} = \ensuremath{\int}(1/4g\up{2})F\up{2} = \textbf{麦克斯韦}；
\item \textbf{数值验证 C}：库仑相光子格点色散 \ensuremath{\omega}(k)=2\ensuremath{\sqrt{\Sigma{}sin\up{2}(k/2)}}，长波 \textbf{\ensuremath{\omega}/|k|=0.9999}（光速 c=1 \checkmark）；
\item \textbf{自旋 2 判据链}：(1) Fierz--Pauli 唯一性（引用：无质量自旋 2 \ensuremath{\implies} 线性化 GR）(2) 旋转本征值 \ensuremath{\pm}2（严格 v1.29）(3) Weinberg 软引力子定理（引用判据，软极限 \ensuremath{\propto}1/q 可测）(4) 弦网低能谱（半严格，Wen）------\textbf{唯一缺口 = Z\down{2}\ensuremath{\to}U(1) 连续极限重整化群 / 4D 无质量自旋 2 场论涌现（世界级，诚实不声称）}。
\end{itemize}
\end{itemize}

\begin{itemize}
\item \textbf{\mtag{v1.34 补充}弦网低能谱与自旋 2 判据实证（行动清单 P1-1，判据链(4)数值化）}\mtag{数值}：用\textbf{链轴四重旋转对称性（C4）}把 Z\down{2} 弦网低能谱分解为三个螺旋度扇区（[R,H]=0、R\up{4}=I 精确验证 \checkmark）：标量（C4=+1）/ \textbf{螺旋度 \ensuremath{\pm}2（C4=\ensuremath{-}1，自旋 2 判据签名）} / 矢量（C4=\ensuremath{\pm}i，光子候选）。系统：双立方体（11 面 2,048 态密集全谱）+ 三立方体链（16 面 65,536 态 Lanczos）。(i) \textbf{\ensuremath{\lambda}=0 精确能级全验证}：2 立方体 512/1,024/512、3 链 8,192/24,576/24,576/8,192------与解析 2\up{n-N\down{c}}\ensuremath{\cdot}C(N\down{c},k) 一致 \checkmark；\textbf{有隙谱}（间隙随 \ensuremath{\lambda} 单调 0.04\ensuremath{\to}1.92，无零能集体模式）；(ii) \textbf{自旋 2 间隙}：\ensuremath{\lambda}=0.3 时螺旋\ensuremath{\pm}2 扇区 \textbf{\ensuremath{\Delta}\ensuremath{\pm}2=0.4919（双立方体）/0.4864（三链）}------尺寸无关（仅变 1\%），张量/矢量模式比标量重 \ensuremath{\sim}19\%（\ensuremath{\Delta}标量=0.41）；(iii) \textbf{载体实证}：二阶张量算符空间 Z\down{fZ}\down{g} 承载 C4=\ensuremath{-}1（螺旋度\ensuremath{\pm}2）本征态 \textbf{30/120 维（3 链）/14/55（双立方体）}，远超一阶空间（3/16）------\textbf{自旋 2 = 二阶张量扇区}（与\mtag{本框架定义 D-2}"自旋 2 = 二阶通量张量涨落"算符层对齐）。\mtag{诚实结论}Z\down{2} 级无无质量自旋 2（阿贝尔有隙，预期物理------引力子需 U(1) 4D 连续极限）；本实证给出自旋 2 模式在弦网谱中的\textbf{谱学定位}（有隙、张量扇区、尺寸无关），非涌现完成。
\end{itemize}

\begin{itemize}
\item \textbf{\mtag{v1.36 补充}3+1D \ensuremath{\lambda}\down{c} 相变 + Wilson RG 数值路径（行动清单 P1-3/P1-4，凝聚相稳定性边界 + 连续极限 RG 层）}\mtag{数值}：\textbf{（a）\ensuremath{\lambda}\down{c} 相变}------通用面构造 build\down{grid}(nx,ny,nz)（链 2/3/4\ensuremath{\times}1\ensuremath{\times}1 + 板 2\ensuremath{\times}2\ensuremath{\times}1，2,048--2.1\ensuremath{\times}10\up{6} 态 LinearOperator Lanczos）+ 序参量 \ensuremath{\langle}A\down{c}\ensuremath{\rangle} + \ensuremath{\chi}=\ensuremath{-}d\ensuremath{\langle}A\ensuremath{\rangle}/d\ensuremath{\lambda}：粗扫（步长 0.2）\ensuremath{\chi} 峰全在 \ensuremath{\lambda}=0（网格伪像），\textbf{细扫（步长 0.02）四尺寸 \ensuremath{\chi} 峰全部锁定 \ensuremath{\lambda}\down{\chi{}}=0.140}（\ensuremath{\chi}\down{max} 2.24--2.28 几乎不变）\ensuremath{\implies} \textbf{3+1D Z\down{2} 弦网凝聚\ensuremath{\to}极化相变点 \ensuremath{\lambda}\down{c}\ensuremath{\approx}0.14}（伪临界值，凝聚相 \ensuremath{\lambda}<0.14 稳定）；\textbf{维数-临界耦合递减谱系}：3+1D 弦网 \ensuremath{\lambda}\down{c}\ensuremath{\approx}0.14 vs 2+1D toric code \ensuremath{\kappa}\down{c}\ensuremath{\approx}0.25--0.32（P0-3b）vs Ising \ensuremath{\beta}\down{c} 2D/3D=0.44/0.22------维数升高临界耦合变小（涨落更强更易退凝聚），\ensuremath{\lambda}\down{c}\ensuremath{\approx}0.14 与对偶 3D Ising \ensuremath{\beta}\down{c}=0.2217 同量级（v1.28 对偶自洽）；\mtag{本框架机制：相变 = 通量场湍流态\ensuremath{\leftrightarrow}各向同性稳态冷凝点}。\textbf{（b）Wilson RG 路径}（Z\down{2}\ensuremath{\to}U(1) 连续极限动力学）------(1) U(1) 自由场（高斯）配分函数 RG 不变性（det 谱分解精确）：缩放律 \textbf{\ensuremath{\beta}'=\ensuremath{\beta}\ensuremath{\cdot}b\up{2-d}}（d=2 临界尺度不变、d=3 g*=0 高斯固定点 IR 自由、d=1 流向强耦合）------\textbf{麦克斯韦 = 高斯固定点 g*=0}；(2) \textbf{1D Ising decimation 精确}：\ensuremath{\beta}'=atanh(tanh\up{2}\ensuremath{\beta})，唯一固定点 \ensuremath{\beta}*=0（流向高斯/连续极限 \checkmark）；(3) \textbf{2D Ising 块 RG（Metropolis + majority rule，16\ensuremath{\to}8\ensuremath{\to}4）------Z\down{2}\ensuremath{\to}U(1) 路径核心演示}：\ensuremath{\beta}=0.3（<\ensuremath{\beta}\down{c}=0.44）\ensuremath{\implies} 0.193\ensuremath{\to}0.162\ensuremath{\to}\textbf{0.063（流向 0 = U(1) 连续极限）}、\ensuremath{\beta}=0.5（>\ensuremath{\beta}\down{c}）\ensuremath{\implies} 0.694\ensuremath{\to}0.677\ensuremath{\to}\textbf{1.900（流向 \ensuremath{\infty} = 凝聚相）}------Z\down{2} 弦网离散化在 \ensuremath{\beta}<\ensuremath{\beta}\down{c} 时 RG 流向 U(1) 高斯固定点。\mtag{诚实标注}格点块平均产生次近邻修正（配分函数不变式漂移 0.35--0.61 非零、最近邻关联被块内稀释------real-space RG 近似，\ensuremath{\beta}'=\ensuremath{\beta}b\up{2-d} 为连续极限主导项）；高斯模型为自由场（RG 平凡，真实 U(1) 需 HOTRG 非微扰）；\ensuremath{\lambda}\down{c} 为小系统伪临界（\ensuremath{\lambda}\down{c}(\ensuremath{\infty}) 需 3D 盒子 = MPS/DMRG）。与 v1.30 四层证据闭合：定理 A（群嵌入）\ensuremath{\to} 命题 B（麦克斯韦线性化）\ensuremath{\to} 数值验证 C（光子 \ensuremath{\omega}=|k|）\ensuremath{\to} RG 流（本工作）------"Z\down{2}\ensuremath{\to}U(1) 连续极限"从单一开放项升级为四层证据；与 v1.34 衔接：3+1D 凝聚相（\ensuremath{\lambda}<\ensuremath{\lambda}\down{c}）即"流向高斯固定点的吸引域"的 RG 语义。
\end{itemize}

\textbf{双路径印证}：引力涌现的统计力学路径（\S{}4，完整 \checkmark）与弦网路径（涌现自旋 2，半严格 \textbf{!}；v1.28 已严格化凝聚层/对偶层）\textbf{两条独立路径指向"引力 = 涌现"}------若等价，弦网集体模式 = 通量场统计涨落（图像性对应，非严格证明）。



\section{动力学（最小作用量）}

\subsection{作用量泛函}

\mtag{推导}最小拉格朗日量密度（多表示版）：

\[\mathcal{L} = \sqrt{-g}\Big[\frac{1}{16\pi G}R - \tfrac14 F_{\mu\nu}F^{\mu\nu} - \tfrac12(\partial\Phi)^2 - V(\Phi) + i\bar\psi\gamma^\mu\nabla_\mu\psi - m\bar\psi\psi - g\Phi\bar\psi\psi - eA_\mu\bar\psi\gamma^\mu\psi\Big]\]

\begin{itemize}
\item 钉扎项 g\ensuremath{\Phi}\ensuremath{\bar{\psi}}\ensuremath{\psi}：环境依赖有效质量 m\down{eff}(x) = m + g\ensuremath{\Phi}(x)------"粒子被风暴钉住"；
\item \ensuremath{\Phi} 为湮灭-再生噪声的凝聚代表（辅助场）。
\end{itemize}

\subsection{场方程（变分导出）}

\begin{table}[htbp]
\centering
\small
\begin{tabularx}{\textwidth}{@{}lXX@{}}
\toprule
通道 & 场方程 & 对应 \\
\midrule
标量 & \ensuremath{\square}\ensuremath{\Phi} = V'(\ensuremath{\Phi}) + g\ensuremath{\bar{\psi}}\ensuremath{\psi} & \textbf{D[\ensuremath{\Phi}]=S 的具体实现} \\
矢量 & \ensuremath{\partial}\down{\mu{}F}\up{\mu{}\nu{}} = eJ\up{\nu{}} & Maxwell \\
张量 & G\down{\mu{}\nu{}} = 8\ensuremath{\pi}G T\down{\mu{}\nu{}} & Einstein \\
\bottomrule
\end{tabularx}
\end{table}

\textbf{退化检验}：g=0、V=\ensuremath{\tfrac12}m\up{2}\ensuremath{\Phi}\up{2} \ensuremath{\to} 自由 Dirac+标量+GR；\ensuremath{\psi}=0、\ensuremath{\Phi}=const \ensuremath{\to} GR+宇宙学常数；弱场低速 \ensuremath{\to} g\down{00} \ensuremath{\approx} \ensuremath{-}(1+2\ensuremath{\Phi}\down{N}/c\up{2})。

\subsection{Noether 守恒律}

\mtag{推导}时间平移\ensuremath{\to}能量、空间平移\ensuremath{\to}动量、转动\ensuremath{\to}角动量、内部 U(1)\ensuremath{\to}电荷（4 条守恒律）；局部 U(1) 缺位（单标量无 A\down{\mu{}}，多表示版补齐）。

\subsection{对称性自发破缺（\ensuremath{\Phi}\up{4} 热相变）}

\[V(\Phi) = -\tfrac12\mu^2\Phi^2 + \tfrac{\lambda}{4}\Phi^4,\qquad V_T = V(\Phi) + \tfrac12 cT^2\Phi^2\]

\begin{itemize}
\item T > T\down{c} = \ensuremath{\mu}/\ensuremath{\sqrt{c}}：\ensuremath{\langle}\ensuremath{\Phi}\ensuremath{\rangle}=0（统一相）；T < T\down{c}：墨西哥帽（破缺）；
\item 复标量升级 \ensuremath{\to} Anderson--Higgs \ensuremath{\to} W/Z 质量 = "穿堂风管道集体模"。
\end{itemize}

\textbf{物理映射}："Higgs=通量场相变" \ensuremath{\iff} 序参量 \ensuremath{\langle}\ensuremath{\Phi}\ensuremath{\rangle} 从 0 涨至 v；"破缺=模态边界浮现" \ensuremath{\iff} 温度（通量密度）降过 T\down{c}。

\subsection{测地线方程}

\mtag{引用}点粒子作用量 S\down{pp} = \ensuremath{-}m\ensuremath{\int}d\ensuremath{\tau} 变分导出 \ensuremath{\Gamma} 方程；牛顿极限 \ensuremath{\vec{a}} = \ensuremath{-}\ensuremath{\nabla}\ensuremath{\Phi}\down{N}。"物体每瞬被调整使压强差最小" \ensuremath{\iff} \ensuremath{\Phi}\down{N} = 通量负压势。



\section{量子力学基础}

\subsection{问题：经典事件流推不出贝尔违背}

\mtag{引用 + 推导}经典隐变量流场只能给出 CHSH \ensuremath{\leq} 2，与实验 CHSH = 2\ensuremath{\sqrt{2}} 矛盾------\textbf{框架必须非交换化}。原文已自知（"你这套要想活，必须能推出 2\ensuremath{\sqrt{2}}\ldots{}\ldots{}源事件的结构性产物"）。

\subsection{非交换事件流}

\mtag{推导}\textbf{命题（非交换事件流 \ensuremath{\Rightarrow} 量子关联）}：把事件流状态从经典概率分布（交换代数 L\up{\infty{}）升级为算子代数（C}* 代数），"相遇---湮灭"事件对应非交换投影算子（Clifford 型 {\ensuremath{\Pi}\down{\mu{}}, \ensuremath{\Pi}\down{\nu{}}}\ensuremath{\neq}0），则：
(i) 纠缠对联合测量给出 CHSH = 2\ensuremath{\sqrt{2}}（Tsirelson 界）；
(ii) 互补性 V\up{2} + I\up{2} \ensuremath{\leq} 1 自动成立；
(iii) 不确定性 \ensuremath{\Delta}A\ensuremath{\cdot}\ensuremath{\Delta}B \ensuremath{\geq} \ensuremath{\tfrac12}|\ensuremath{\langle}[A,B]\ensuremath{\rangle}| 自动成立。

\textbf{框架已有非交换元素}：Clifford 代数（微变不可交换）、路径积分 (\ensuremath{-}1)\up{N\down{exchange}（交换符号）、叠加态几何起源（正反风暴中心全局等价）、量子直线（同源事件）。}

\textbf{数值验证}：Bell 态 |\ensuremath{\psi}\up{-}\ensuremath{\rangle} + 最优测量基 \ensuremath{\implies} S = \ensuremath{-}2.8284 = \ensuremath{-}2\ensuremath{\sqrt{2}} \checkmark；路径态内积 a 时 V\up{2}+I\up{2} = 1 \checkmark。

\subsection{量子游走：非交换事件流的动力学}

\mtag{推导}离散时间量子游走 U = S(C\ensuremath{\otimes}I)：硬币翻转 C（Hadamard，非交换）+ 条件位移 S。数值：量子 \ensuremath{\sigma} \ensuremath{\sim} t（弹道）vs 经典 \ensuremath{\sqrt{t}}（扩散）------非交换化保留相干性。

\subsection{定位澄清}

\mtag{推导}一旦非交换化，框架的量子部分与标准 QM \textbf{数学上重合}。框架的独特贡献 = \textbf{为"为什么是量子力学"提供本体论动机}（Gel'fand--Naimark 精神：交换 C* \ensuremath{\iff} 经典概率；非交换 C* \ensuremath{\iff} 量子态）。定位是 \textbf{'t Hooft 元胞自动机式}：底层确定事件流，可观测代数非交换。



\section{宇宙学}

\subsection{泄漏机制的能标依赖（\ensuremath{\varepsilon}\down{A}）}

\mtag{推导}\ensuremath{\varepsilon}\down{A}(E) = \ensuremath{\varepsilon}\down{0}(M\down{Pl}/E)\up{\beta{}，两边界条件：}
\begin{itemize}
\item E \ensuremath{\sim} 1 TeV \ensuremath{\to} \ensuremath{\varepsilon} \ensuremath{\sim} 6\ensuremath{\times}10\up{-10}（重子不对称 \ensuremath{\eta} \ensuremath{\approx} \ensuremath{\varepsilon}）；
\item E \ensuremath{\sim} 1 meV \ensuremath{\to} \ensuremath{\varepsilon} \ensuremath{\sim} 1（暗能量 \ensuremath{\rho}\down{\Lambda{}} \ensuremath{\sim} (meV)\up{4}）；
\item \ensuremath{\beta} = log(6\ensuremath{\times}10\up{-10})/log(10\up{-15}) \ensuremath{\approx} \textbf{0.615}，\ensuremath{\varepsilon}\down{0} \ensuremath{\approx} 7.6\ensuremath{\times}10\up{-20}。
\end{itemize}

\textbf{截断能标}\mtag{推导}：概率归一化 \ensuremath{\varepsilon}\down{A}(E\down{c}) = 1 \ensuremath{\implies} E\down{c} = M\down{Pl}\ensuremath{\cdot}\ensuremath{\varepsilon}\down{0}\up{1/\beta{}} = \textbf{1.01 meV}------与暗能量特征能标 2.3 meV 同量级（(E\down{c})\up{4} 差 27 倍，O(1) 系数待定）。

\subsection{统计衰减因子（\ensuremath{\varepsilon}\down{B}）与暗物质}

\mtag{推导}普朗克涨落原始密度 \ensuremath{\rho}\down{raw} = 2m\down{P}/l\down{P}\up{3} \ensuremath{\approx} 10\up{97} kg/m\up{3}；统计衰减 \ensuremath{\varepsilon}\down{B} = (l\down{P}/R\down{U})\up{2} \ensuremath{\approx} 10\up{-120}（几何来源，非 (1/3)\up{n------n} 实际=251）；暗物质 \ensuremath{\rho}\down{DM} = \ensuremath{\varepsilon}\down{B}\ensuremath{\cdot}\ensuremath{\rho}\down{raw} \ensuremath{\approx} 1.6\ensuremath{\times}10\up{-25} kg/m\up{3}（与观测 2.3\ensuremath{\times}10\up{-27} 差 \ensuremath{\sim}70 倍------\textbf{\mtag{待验证}}需 f\down{space} \ensuremath{\approx} 0.014 或机制修正）。

\textbf{\ensuremath{\varepsilon}\down{B} 的拓扑来源（H1--H4）\mtag{v1.7 补充}}："\ensuremath{\varepsilon} 是拓扑强制的'最小涡心密度'"（本框架机制）定量化\mtag{命题链 H1--H4，基于 Poincar\'{e}--Hopf}：
\begin{itemize}
\item \textbf{H1（S\up{2} 严格）}：\ensuremath{\chi}(S\up{2}) = 2 \ensuremath{\implies} 切向量场零点指数和 = 2 \ensuremath{\implies} 方向空间至少 2 个涡心\mtag{拓扑论证}；
\item \textbf{H2（涡心占比）}：\ensuremath{\varepsilon}\down{topo}(R) = (l\down{P}/R)\up{2}/(2\ensuremath{\pi})------事件方向落在涡心单元的概率；
\item \textbf{H3（分量结构）}：\ensuremath{\varepsilon}\down{total} = \ensuremath{\varepsilon}\down{topo} + \ensuremath{\varepsilon}\down{A}------普朗克尺度拓扑主导（0.159），TeV 以上动力学主导（\ensuremath{\varepsilon}\down{A}）；\ensuremath{\varepsilon}\down{0} = 7.67\ensuremath{\times}10\up{-20} 是动力学分量锚定值；
\item \textbf{H4（重大发现）}：\ensuremath{\varepsilon}\down{B}(宇宙) = 1.35\ensuremath{\times}10\up{-123} \ensuremath{\approx} \ensuremath{\varepsilon}\down{topo}(宇宙)\ensuremath{\times}2\ensuremath{\pi}------\textbf{暗物质衰减因子恰好处于拓扑强制最小值量级}，\ensuremath{\varepsilon}\down{B} 的几何来源获得拓扑背书；
\item \textbf{S\up{3} 澄清\mtag{修正}}：\ensuremath{\chi}(S\up{3})=0 且 S\up{3} 可平行化（切丛平凡），毛球定理不适用于三维空间切片（早期版本修正）------空间切片涡心需挂谷+泄漏补充。
\end{itemize}

\subsection{暗物质叠加模型}

\mtag{推导}\textbf{暗物质 = 本支普朗克涨落残余（\ensuremath{\varepsilon}\down{B} 机制）+ 镜像支引力投影（跨支机制）}：

\[\rho_{DM} = \varepsilon_B\,\rho_{raw}\,f_{space} + \kappa_{mirror}\,\rho_{mirror}\]

\begin{itemize}
\item 镜像支投影：\ensuremath{\kappa}\down{extra} \ensuremath{\sim} \ensuremath{\Omega}\down{mirror}/\ensuremath{\Omega}\down{total} \ensuremath{\approx} 0.5（引力跨支共享，电磁/强/弱被支界过滤）\ensuremath{\implies} \ensuremath{\rho} \ensuremath{\sim} 4\ensuremath{\times}10\up{-27} kg/m\up{3} \textbf{与观测同量级}------\textbf{暗物质主体可能来自镜像支引力投影}\mtag{待验证}；
\item 可检验预言：\textbf{void-CPT-\ensuremath{\kappa} 异常}（空洞区弱透镜偏强 + 质量-光脱耦，10--30 Mpc/h 尺度最显；候选对象 CDG-2、J0613+52）。
\end{itemize}

\subsection{暗能量与 w(z)}

\mtag{推导}暗能量 = \ensuremath{\varepsilon}\down{A} 饱和区贡献 \ensuremath{\rho}\down{\Lambda{}} \ensuremath{\sim} (E\down{c})\up{4}；饱和能标随宇宙膨胀缓慢漂移 E\down{c}(z) = E\down{c},0(1+z)\up{\varepsilon{}} \ensuremath{\implies}

\[w = -1 + \frac{4\varepsilon}{3} = -1 + 2\nu\]

观测约束（DESI 2024+CMB：w = \ensuremath{-}1.028\ensuremath{\pm}0.031）：|\ensuremath{\varepsilon}| < 0.023、|\ensuremath{\nu}| < 0.0155。\textbf{框架预测 w = const（wa=0）}------DESI DR2/Euclid 可直接检验。\textbf{\mtag{v2.7 补充：DESI DR2（2025-03）判别状态（响应审稿 M5）}}：DESI DR2+CMB 排除 \ensuremath{\Lambda}CDM \textbf{3.1\ensuremath{\sigma}}（+Pantheon+ 2.8\ensuremath{\sigma} / +Union3 3.8\ensuremath{\sigma} / +DESY5 4.2\ensuremath{\sigma}），CPL 最佳拟合 \textbf{w\down{0}\ensuremath{\approx}\ensuremath{-}0.76\ensuremath{\pm}0.06、wa\ensuremath{\approx}\ensuremath{-}0.77\ensuremath{\pm}0.23}------对框架 \textbf{wa=0 预言构成 \ensuremath{\sim}3.3\ensuremath{\sigma} 压力}（框架暗能量预言正被\textbf{已有数据}判别：若 DDE 信号确认 \ensuremath{\implies} 该预言被否证；若为系统效应 \ensuremath{\implies} 框架存活）。\textbf{裁决时间线}：DESI 最终分析 2026-04 + Euclid 首宇宙学数据 2026-10------\textbf{0-2 年内最终裁决}。诚实标注：DDE 与系统效应之争未决（ACT DR6/SPT-3G CMB 支持 \ensuremath{\Lambda}CDM）；w\down{0} 判别为 wCDM/CPL 跨模型对照（非直接可比，以 wa=0 检验为准）。

\subsection{原始错误修正记录}

\mtag{修正}(1) "meV 能标 = 哈勃尺度"对应错误（差 \ensuremath{\sim}30 个数量级；meV 波长 1.24 mm）；(2) \ensuremath{\rho}\down{raw} 量纲笔误（kg/(m\up{3}\ensuremath{\cdot}s) 应为 kg/m\up{3}）；(3) "120"来源 = (l\down{P}/R\down{U})\up{2} 而非 (1/3)\up{120。}



\section{可检验预言}

\begin{table}[htbp]
\centering
\small
\begin{tabularx}{\textwidth}{@{}lXXX@{}}
\toprule
\# & 预言 & 观测窗口 & 状态 \\
\midrule
1 & \textbf{B1--B5 黑洞强场预言包}：极端 Kerr 霍金辐射角分布偏离轴对称（B1）；原初黑洞熵体积修正 \ensuremath{\Delta}S/S\ensuremath{\sim}l\down{P}/R\ensuremath{\sim}10\up{-8}（B2）；霍金辐射-落入物质纠缠（B3，信息翻转判决签名）；ER 桥 packing 上限\ensuremath{\to}BBH 角模畸变（B4）；非热修正 \ensuremath{\propto}1/S\down{BH}（B5） & EHT/LIGO、原初黑洞 \ensuremath{\gamma} 射线、LISA 时代、ET/CE & 理论可检验（见 \S{}4.10） \\
2 & \textbf{void-CPT-\ensuremath{\kappa} 异常}（空洞弱透镜偏强、质量-光脱耦） & DES/Y3、DESI、Euclid、Rubin LSST & 可检验（分析方案已给出） \\
3 & \textbf{w = \ensuremath{-}1+4\ensuremath{\varepsilon}/3 常数偏离（wa=0）} & \textbf{DESI DR2（已有，2025-03）}、Euclid 2026-10 & \textbf{判别中：wa=0 遭 3.3\ensuremath{\sigma} 压力（DDE vs 系统效应未决），2026 最终裁决} \\
4 & \textbf{SU(3) \ensuremath{\beta} 函数阶梯跳跃}（第四代阈值） & FCC-hh、LHAASO/CTA & 理论可算，数值待推导 \\
5 & 第四代费米子（>1 TeV，z>1000 衰变）痕迹 & LHC、宇宙线、CMB-S4 & 与 LHC 未发现存在张力 \\
6 & \textbf{霍金辐射非热修正} \ensuremath{\sim}1/S（原初黑洞 \ensuremath{\sim}10\up{-20}） & 原初黑洞蒸发 & 接近窗口边缘 \\
6' & \textbf{火墙厚度非热修正} \ensuremath{\sim}\ensuremath{\sqrt{l\down{P}/R}}（原初黑洞 \ensuremath{\sim}10\up{-9}\ensuremath{\sim}10\up{-4}，F4 发射层预言） & 原初黑洞 \ensuremath{\gamma} 射线 & 理论可检验 \\
7 & PeV 中微子味比率异常 & IceCube-Gen2、GRAND & 定性预言 \\
8 & 纠缠熵面积律系数（框架含 ln 2 因子） & 量子模拟/多体系统 & 理论可检验 \\
9 & \textbf{拓扑序激发（量子双/弦网）}：环面基态简并 \ & G\ & \up{2}、任意子分数统计、弦端点激发能隙 \ensuremath{\Delta}\ensuremath{\sim}1（\S{}5.6--5.8） \\
10 & \textbf{禁闭弦张力数值实现 + 相变点定位（v1.20--v1.23）}：电荷对 \ensuremath{\sigma}\down{E}\ensuremath{\approx}2\ensuremath{\kappa}（X 场，热力学极限收敛 <0.02\%）与磁通对 \ensuremath{\sigma}\down{m}\ensuremath{\approx}2\ensuremath{\lambda}（Z 场，约束扇区：0.21/1.09/3.04 @ \ensuremath{\lambda}=0.5/1.0/2.0）；\ensuremath{\lambda}=0/\ensuremath{\kappa}=0 严格 deconfined；\textbf{相变定位（P0-3/P0-3b）}：开边界 FSS 0.23--0.27 \ensuremath{\to} 环面半极化点 \ensuremath{\kappa}*\ensuremath{\approx}0.20\ensuremath{\pm}0.01（尺寸无关）\ensuremath{\to} 真相变 \textbf{\ensuremath{\kappa}\down{c}\ensuremath{\approx}0.25--0.32}（与文献 X 场版本 h\down{xc}\ensuremath{\approx}0.32 一致），弦张力数值处于禁闭相 \checkmark & 量子模拟（Z\down{2} 格点规范实现）、超导电路 & 数值可复现（脚本 \down{scale}\down{extrap3}/\down{p01}\down{sector}/\down{p03}\down{fss}/\down{p03}\down{torus}*） \\
\textbf{S1} & \textbf{\ensuremath{\alpha} = 1.00\ensuremath{\pm}0.06}（风暴过载阴影兼容 GR，v1.15） & EHT M87*/Sgr A* & \checkmark 已兼容 \\
\textbf{S2} & \textbf{QNM 频率 f = f\down{GR}/\ensuremath{\alpha}}（\ensuremath{\alpha}\ensuremath{\neq}1 偏离 \ensuremath{\pm}5--11\%） & ET/CE 未来 & 可检验（\ensuremath{\pm}0.01） \\
\textbf{S3} & \textbf{火墙弥散 \ensuremath{\delta}/R \ensuremath{\sim} 10\up{-9}（原初黑洞，框架独有）} & 原初黑洞蒸发/引力波 & 框架独有签名 \\
\textbf{S4} & \textbf{过载区无奇点 + 声速消失 + 光学透明}（钳制物态 P=P\down{P}\ensuremath{\cdot}tanh(\ensuremath{\rho}/\ensuremath{\rho}\down{P})，v1.24 微观推导：P\ensuremath{\leq}\ensuremath{\rho}c\up{2} 处处满足、\ensuremath{\partial}P/\ensuremath{\partial}\ensuremath{\rho}\ensuremath{\geq}0 无坍缩；c\down{s}\up{2}=c\up{2}\ensuremath{\cdot}sech\up{2}(\ensuremath{\rho}/\ensuremath{\rho}\down{P})，\ensuremath{\rho}\ensuremath{\geq}\ensuremath{\rho}\down{P} 无弹性波；v1.25 均匀浆糊 \ensuremath{\bar{\rho}}=3c\up{6}/(32\ensuremath{\pi}G\up{3}M\up{2}) + 光学深度 \ensuremath{\tau}\ensuremath{\sim}10\up{-40}（宏观）/0.06（普朗克窗口）；\textbf{v1.31 吸积盘辐射传输：标准薄盘主体与 GR 一致 + 过载区全透明 + 火墙非热修正（原初窗口）}------框架独有） & 吸积流/过载区辐射、普朗克质量黑洞蒸发、原初黑洞 \ensuremath{\gamma} 射线 & \checkmark 流体模型+微观推导+径向结构+吸积盘传输（\S{}4.6） \\
\textbf{S5} & 阴影边缘普朗克弥散（宏观 \ensuremath{\sim}10\up{-19}\ensuremath{\sim}10\up{-24} 不可测） & 超长基线干涉 & 理论预言 \\
\textbf{S-E1} & \textbf{EMRI 近日点进动率偏高 \ensuremath{\delta}=0.30\ensuremath{\cdot}x\down{p}}（1--3\%，2PN g\down{rr} 33\% 分叉签名；v1.26 提案 + v1.27 Fisher + v1.32 Kerr 贝叶斯 + \textbf{v1.35 a-\ensuremath{\delta} 联合}） & LISA EMRI（M\ensuremath{\sim}10\up{5}--10\up{7} M\ensuremath{\odot}） & \textbf{v1.35 修正后：\ensuremath{\Delta}lnL 判据 27\ensuremath{\sigma}（2yr）/28\ensuremath{\sigma}（4yr），a-\ensuremath{\varphi}0 全自由稳健（\ensuremath{\sigma}\down{\delta{}}=1.6\ensuremath{\times}10\up{-6}，v1.32 的 9.7\ensuremath{\times}10\up{-4} 欠采样虚高 600\ensuremath{\times}）------单源判决} \\
\textbf{S-E2} & \textbf{EMRI 引力波相位偏离 GR}（进动调制相位差 \ensuremath{\sim}2\ensuremath{\times}\ensuremath{\Delta}\ensuremath{\Phi}\down{prec}） & LISA & 可检验 \\
\textbf{S-E3} & \textbf{EMRI 轨道衰减/啁啾相位差}（3PN g\down{00} 11\% 分叉） & LISA & 可检验（\ensuremath{\Delta}\ensuremath{\Phi}\down{chirp}\ensuremath{\sim}2\ensuremath{\times}10\up{2} rad） \\
\textbf{S-E4} & \textbf{进动差-半径标度律} \ensuremath{\delta}\ensuremath{\propto}x\down{p}（0.6\% @ p=50M \ensuremath{\to} 3.0\% @ p=10M）------框架独有，多源联合检验 & LISA & 可检验 \\
\textbf{S-E5} & \textbf{FQH 引力子模衰减率含 \ensuremath{\varepsilon} 项}（\ensuremath{\Gamma}=\ensuremath{\Gamma}\down{Coulomb}+\ensuremath{\varepsilon}\ensuremath{\cdot}\ensuremath{\Gamma}\down{flux}，\ensuremath{\varepsilon}\ensuremath{\sim}10\up{-20} 泄漏；本框架预言，v1.29） & 极低温强磁场 FQH 平台（线宽余项） & 可检验（框架独有） \\
\textbf{S-E6} & \textbf{胶球 = 非阿贝尔通量自束缚态}（框架翻译[BESIII 2026]；BESIII X(2370) 已确认：0\up{-+}、2395 MeV、味单态，v1.33 锚点） & BESIII/BEPCII、J/\ensuremath{\psi} 衰变、未来强子对撞机 & 框架性解释（QCD 事实 + 框架翻译） \\
\bottomrule
\end{tabularx}
\end{table}



\section{结论与开放问题}

\subsection{框架地位评估}

本框架从四条公理出发，为四种基本力的统一提供了\textbf{涌现式统计力学论证}，完成度 \ensuremath{\sim}90\%：

\begin{itemize}
\item \textbf{引力侧}（最完整）：\ensuremath{\gamma}=1、\ensuremath{\beta}=1 全通过（PPN 双参数判决）；微变权命题；玻尔兹曼分布（热平衡命题）；Jacobson 涌现路径；S=A/4 微态计数（k=2\up{1/4}）；\textbf{packing 命题包}（k\ensuremath{\geq}4 下界新证 + Tammes 正四面体 + 1/4 系数几何来源）；\textbf{命题 6}（内空间编码优先性，框架内演绎）；\textbf{观测者模型 O1--O4}（互信息框架 + 均匀先验严格化）；\textbf{现在平面选择 N1--N4}（涡动力学 + 对称性破缺）；\textbf{初始微扰来源 P1--P4}（"必然的偶然"，与宇宙学侧同源）------\textbf{"引力 = 通量场的统计力学涌现"获得从公理到方程级、从微观到观测者的完整论证链}；
\item \textbf{规范侧}（v1.19 大幅强化 + v1.20 数值闭环）："2 vs 3"配对阶数（引理 A--D 表示论层严格）；\textbf{量子双代数命题包 A1--A6}（规范结构 = 离散代数涌现，\S{}5.6）；\textbf{框架哈密顿量 B1--B5}（涌现规范对称性 [H\down{F},A]=[H\down{F},B]=0 严格 + 2+1D 质量隙 + 环面简并 |G|\up{2}，\S{}5.7）；\textbf{禁闭弦张力数值实现}（X 场电荷 \ensuremath{\sigma}\down{E}\ensuremath{\approx}2\ensuremath{\kappa} 热力学极限收敛 + Z 场磁通 \ensuremath{\sigma}\down{m}\ensuremath{\approx}2\ensuremath{\lambda} 约束扇区 + deconfined 退化严格验证，\S{}5.7）\textbf{+ 相变点定位}（v1.23：环面半极化 \ensuremath{\kappa}*\ensuremath{\approx}0.20\ensuremath{\pm}0.01、真相变 \ensuremath{\kappa}\down{c}\ensuremath{\approx}0.25--0.32 与文献 h\down{xc}\ensuremath{\approx}0.32 一致，P0-3b，\S{}5.7）；\textbf{弦网 3+1D 路线}（量子直线=弦边、四力=弦网序涌现、引力双路径印证，\S{}5.8）------\textbf{规范侧从"结构启发式"升级为"命题包 + 显式哈密顿量 + 涌现对称性命题 + 禁闭数值实现"}；
\item \textbf{量子侧}：非交换化闭环（CHSH=2\ensuremath{\sqrt{2}}、互补性）；'t Hooft 式定位；
\item \textbf{宇宙学侧}：\ensuremath{\varepsilon} 分工清晰、w(z) 参数化；暗物质叠加模型（定量未完全闭合）；初始微扰与重子不对称/暗能量同源（F26）。
\end{itemize}

\subsection{数值可靠性评估（v2.3，响应第三方审稿 M3）}

\textbf{判据校准（\ensuremath{\kappa}\down{c}）}：细网格（52 点）4 尺寸环面数据分析------内部判据（m=0.5 / \ensuremath{\chi} 峰 / \ensuremath{\Delta} 最小）1/L 外推给出 0.11--0.17，对文献精确值 h\down{c}\ensuremath{\approx}0.3285（iPEPS）\textbf{系统性偏低 47--67\%}（toy-model 尺度落差被定量证实）；\textbf{决定性诊断：m(h\down{c}) 随 L 外推 \ensuremath{\to} 0.93 \ensuremath{\neq} 0 \ensuremath{\implies} \ensuremath{\langle}X\ensuremath{\rangle}（单点磁化）非相变序参量}（v1.23"m=0.5 为半极化点"判据修正获得机制级确认）；\ensuremath{\chi}\down{max} \ensuremath{\propto} L\up{1.12}（2D Ising 类 \ensuremath{\gamma}/\ensuremath{\nu}=1.75）支持相变存在但被小格点宽化。\textbf{校准结论}：真相变点取文献值 \textbf{\ensuremath{\kappa}\down{c}\ensuremath{\approx}0.25--0.32}（框架对称约定与文献等价，v1.23），判据基础升级为"文献校准 + 弦序参量 \ensuremath{\langle}W\ensuremath{\rangle} + Binder 交叉 + 大格点 MPS/DMRG"。

\textbf{\ensuremath{\lambda}\down{c} 有限尺寸边界}：4 尺寸 \ensuremath{\chi}\down{max} \ensuremath{\propto} L\up{-0.04}（不增长）\ensuremath{\implies} \ensuremath{\chi} 峰为宽化伪峰；\ensuremath{\lambda}\down{\chi{}} 无漂移（全 0.140）\ensuremath{\implies} \textbf{\ensuremath{\lambda}\down{c}(\ensuremath{\infty}) 保守区间 [0.140, 0.140]，精确定位需 3D 盒子 MPS/DMRG}（开放数值问题，诚实保留）。

\textbf{EMRI 系统误差上界}：quadrupole 波形截断系统相位误差 \ensuremath{\delta}\ensuremath{\Phi}\down{sys}\ensuremath{\approx}1 rad（文献量级，Huerta \& Gair 2011 / Babak 2017）\ensuremath{\implies} \textbf{\ensuremath{\sigma}\down{\delta{}},tot\ensuremath{\approx}2.8\ensuremath{\times}10\up{-6}（系统误差与统计误差同量级，\ensuremath{\sigma}\down{sys}/\ensuremath{\sigma}\down{stat}\ensuremath{\approx}1.4）}------\ensuremath{\sigma}\down{\delta{}} 的参数精度受 quadrupole 限制（Teukolsky 可降至 \ensuremath{\sim}2\ensuremath{\times}10\up{-8}）；但 \ensuremath{\delta}=2\ensuremath{\times}10\up{-2} \ensuremath{\gg} \ensuremath{\sigma}\down{\delta{}},tot（7108\ensuremath{\sigma}）\ensuremath{\implies} \textbf{27\ensuremath{\sigma}/28\ensuremath{\sigma} 框架判别不受系统误差威胁}。

\textbf{\ensuremath{\to} toy-model 适用边界明确化}：相变判据（\ensuremath{\kappa}\down{c}）依赖文献校准与序参量升级；\ensuremath{\lambda}\down{c}(\ensuremath{\infty}) 为开放数值问题；EMRI 判别稳健但精度受波形近似限制（详见《M3补强报告》）。

\subsection{开放问题（按重要性排序）}

\begin{enumerate}
\item \sout{内空间编码优先性的动力学证明} \textbf{\checkmark 已解决（v1.2）}：命题 6 三引理链闭合（各向同性抹除外空间信息 + 事件点状性排除方向自由度 + 旋量幸存承载内空间信息）------三前提均为框架公理直引，未引入新假设；\textbf{"方向离散化"假设 \checkmark 已解决（v1.21，B4）}：普朗克管径\mtag{本框架设定：方向离散化}\ensuremath{\implies} 方向簇 k=4（王虹/Tammes）\ensuremath{\implies} Bekenstein 上界恢复（\S{}4.9、\S{}4.10）；
\item \textbf{非阿贝尔分量的完全涌现}：SU(2)\ensuremath{\times}SU(3) 从通量场"色-方向"耦合的动力学（含质量隙、手性）------\checkmark 已推进（B1：f\up{abc} 拓扑 + Jacobi=管网络 + m\down{g}\ensuremath{\approx}\ensuremath{\sqrt{\sigma{}}}）；\textbf{v1.33 推进（P0-2）}：S\down{3} 非阿贝尔 3+1D 面模型数值（46,656 态）------\ensuremath{\lambda}=0 凝聚 E0=\ensuremath{-}N\down{c}、简并 6\up{n-1}=7776、\textbf{激发按共轭类分两类}（非阿贝尔指纹）+ 胶球实验锚点（BESIII X(2370)，\S{}5.7）；4D SU(3) 完全涌现仍开放；
\item \textbf{强场度规形式}：框架强场是"过载区/翻转面"（机制完整，七环链）------\textbf{v1.21 更新}：各向同性指数度规（\S{}4.5b）的阴影 5.436GM/c\up{2} 与 EHT 相容（修正 v1.14 的"无度规承诺"为"有与观测相容的几何预言"）；\textbf{v1.24 更新}：过载区钳制物态已获微观推导（P=P\down{P}\ensuremath{\cdot}tanh(\ensuremath{\rho}/\ensuremath{\rho}\down{P}) + 声速剖面 c\down{s}\up{2}=c\up{2}\ensuremath{\cdot}sech\up{2}(\ensuremath{\rho}/\ensuremath{\rho}\down{P})，\S{}4.6）；饱和模态确定视界几何的完整推导仍开放（\S{}4.6）；
\item \textbf{统计\ensuremath{\to}几何的严格化}：遍历性从框架动力学推出、混沌传播适用性、截断定量------\checkmark 已推进（B2：遍历性从公理 1、截断 r\down{c}=GM/c\up{2}=视界、B2-2 碰撞核层次）；\textbf{v2.6 推进（响应审稿 M1）}：玻尔兹曼涌现数值证明（DSMC 温度 0.00\%/分布 7.0\%/H 定理 \ensuremath{-}2.27/遍历性 0.00\%；硬球对照 H \ensuremath{-}2.38）------热平衡=长时极限涌现非假设（循环风险解除）；H\down{event} 三层候选（经典碰撞/谱作用/量子 Lindblad）已构造；LHC 能标相容论证（预言无 LHC 新物理）；湮灭型 BBGKY 完整证明仍开放（诚实保留）；
\item \sout{观测限制\ensuremath{\to}信息容量的严格化} \textbf{\checkmark 已解决（v1.3）}：O2 给出互信息框架 S\down{access}(R) = I(R;\ensuremath{\Sigma})（\S{}4.9）；\textbf{v1.21 强化}：互信息面积律 I=H(\ensuremath{\Sigma}) 从体=表面投影推出（B5，\S{}4.9）；
\item \textbf{\ensuremath{\beta}\down{ppn} 之外的更高阶 PPN 参数}：通量遮挡的非线性修正------\checkmark 已解决（v1.21，B3）：完整 PPN 参数集（\ensuremath{\gamma}=\ensuremath{\beta}=1 其余 0）+ 2PN/3PN 分叉量化（\S{}4.5b）；
\item \sout{两维时间卷曲的动力学} \textbf{\checkmark 已解决（v1.21，B4'）}："恰好两维被卷曲" = 3\ensuremath{-}1=2 定义性闭合（三平面等价 \ensuremath{\to} 锁定恰好 1 个（N1--N4）\ensuremath{\to} 剩余 2 个相对卷曲，无自由参数）；剩余严格化：N2 的 4\ensuremath{\pi} 相位锁定相容性方程（深度开放）；
\item \sout{毛球拓扑 \ensuremath{\implies} \ensuremath{\varepsilon} 下界定量形式化} \textbf{\checkmark 已解决（v1.7）}：命题链 H1--H4------\ensuremath{\varepsilon}\down{topo}(R) = (l\down{P}/R)\up{2}/(2\ensuremath{\pi})（Poincar\'{e}--Hopf 支撑）；重大发现 \ensuremath{\varepsilon}\down{B}(宇宙) \ensuremath{\approx} \ensuremath{\varepsilon}\down{topo}(宇宙)\ensuremath{\times}2\ensuremath{\pi}（暗物质衰减因子 = 拓扑分量宇宙尺度值）；S\up{3} 可平行化澄清（毛球不适用，早期版本修正）（\S{}8.2）；
\item \sout{"涨落不可避免"的动力学证明} \textbf{\checkmark 已解决（v1.21，B5'）}：三层结构------统计层（密度涨落 \ensuremath{\sigma}/\ensuremath{\rho}\ensuremath{\sim}1/\ensuremath{\sqrt{N}} 泊松）+ 拓扑层（毛球 Poincar\'{e}--Hopf \ensuremath{\chi}=2）+ 泄漏层（\ensuremath{\varepsilon} 泊松随机 \ensuremath{\implies} 涨落方向各向同性）------"绝对均匀"有限 N 下不静止（\S{}4.9，P2 遗留闭合）；
\item \textbf{为什么 n=3（规范侧 A6）}：Clifford 内空间 4 分量 / 数轴分类 / 费马可解性的规范侧统一（\S{}5.6 A6------无 Tammes 式唯一性定理可用，比引力侧更难）------\checkmark 已推进（B6：编号错位统一 + 五层论证链）；\textbf{v2.5 推进（响应审稿 M6）}：配对阶数\ensuremath{\to}规范群的三准则结构筛选（S1 维度匹配 d\down{min}=n / S2 同号配对阶数 n\ensuremath{^\prime} / S3 三阶反对称 \ensuremath{\varepsilon}\down{ijk}）------候选群池 10 个逐一排除唯一幸存（n=1\ensuremath{\to}U(1)、n=2\ensuremath{\to}SU(2)、n=3\ensuremath{\to}SU(3)）；\textbf{SO(3) 因 3 自对偶（n\ensuremath{^\prime}=2）被硬排除}（3\ensuremath{\otimes}3 含平凡 \ensuremath{\to} 无禁闭强制），"恰好 SU(3) 而非 SU(4)"获结构筛选支撑（n=4 由空间三维挂谷维数排除）；剩余偶然性 = 涡卷绕模数 n=1,2,3 的微观来源（\S{}5.2 诚实保留）；
\item \textbf{J、K 耦合强度的来源}：框架哈密顿量 H\down{F}=\ensuremath{-}J\ensuremath{\Sigma}A\down{v}\ensuremath{-}K\ensuremath{\Sigma}B\down{p} 的耦合强度需从事件流微观动力学推出（\S{}5.7 B1 的"形状"已有，强度未推导）------\textbf{v1.21 推进（B7）}：同源论证（J=配对、K=闭合同为风暴通量交互）\ensuremath{\implies} 量级 J/K\ensuremath{\sim}O(1)；toric code 相变 \textbf{h\down{xc}\ensuremath{\approx}0.32}（X 场/\ensuremath{\sigma}\up{x} 版本；0.328 为 \ensuremath{\sigma}\up{z} 场版本值）\ensuremath{\implies} J/K\ensuremath{\sim}3 参考值（框架侧环面得半极化 \ensuremath{\kappa}*\ensuremath{\approx}0.20、真相变 \ensuremath{\kappa}\down{c}\ensuremath{\approx}0.25--0.32，P0-3b，交叉验证 \checkmark）；具体值无第一性推导（仍开放）；
\item \textbf{\ensuremath{\varepsilon} 泄漏项的哈密顿量形式}：\ensuremath{-}\ensuremath{\lambda}\ensuremath{\Sigma}(\ensuremath{-}1)\up{g\neq{}e} 质量项驱动的拓扑序\ensuremath{\to}禁闭转变------\textbf{\checkmark 已落地（v1.20）}：\ensuremath{\varepsilon} 项哈密顿量已构造并数值研究（\S{}5.7），弦张力 \ensuremath{\sigma}\down{m}/\ensuremath{\sigma}\down{E} 双版本实现 + deconfined 退化严格验证；\textbf{\checkmark 相变点定位（v1.22 P0-3 开边界 FSS；v1.23 P0-3b 环面精确化）}：有限尺寸无相变点（\ensuremath{\lambda}>0 简并立即打破、弦密度连续下降）是有限格点表现；环面（周期边界）给出半极化点 \textbf{\ensuremath{\kappa}*\ensuremath{\approx}0.20\ensuremath{\pm}0.01}（尺寸无关）；\textbf{判据修正}------m=0.5 交叉为半极化点非相变点（环面实测 m(h=0.32)\ensuremath{\approx}0.84），真相变 \textbf{\ensuremath{\kappa}\down{c}\ensuremath{\approx}0.25--0.32}，与文献 X 场（\ensuremath{\sigma}\up{x）版本} \textbf{h\down{xc}\ensuremath{\approx}0.32} 一致（0.328 为 \ensuremath{\sigma}\up{z} 场版本值，v1.23 已修正；精确值待 MPS/DMRG，\S{}5.7）；
\item \textbf{3+1D 弦网严格化}：涌现麦克斯韦/GR 的一般定理（\S{}5.8 半严格\ensuremath{\to}严格，阶段 C 遗留）+ \textbf{4D YM 质量隙}（=千禧年问题，阶段 D，不声称）------\textbf{v1.21 确认（B8）}：候选路线图定位（量子直线=弦网边映射严格化），4D 不声称；\textbf{v1.28 推进（后续工作(2)）}：(1) 量子直线-弦网边 Hopf 桥接落地；(2) \textbf{\ensuremath{\lambda}=0 弦网凝聚严格证明}（E0=\ensuremath{-}N\down{c}、简并 2\up{n-N\down{c}}、\ensuremath{\Delta}=2）；(3) 对偶映射严格；(4) 涌现 U(1) 路径明确化（对偶严格，连续极限开放）；\textbf{v1.29 推进}：涌现 U(1) 路径补齐（Villain 紧致化 + 库仑相麦克斯韦 \ensuremath{\omega}=|k|，文献引用）；\textbf{自旋 2 = 二阶张量旋转本征值 \ensuremath{\pm}2 严格验证}（本框架自旋 2 机制）+ 三层论证（框架严格/弦网半严格/南大 FQH 投影证据）+ \textbf{FQH \ensuremath{\varepsilon} 项预言}（\S{}9 S-E5）；\textbf{v1.30 推进}：\textbf{定理 A（严格）} Z\down{2} \ensuremath{\subset} U(1) 群嵌入 + \textbf{命题 B（解析）} cos 线性化 = 麦克斯韦 + \textbf{数值验证 C} 光子 \ensuremath{\omega}=|k|（长波 0.9999）+ \textbf{自旋 2 判据链}（Fierz--Pauli + 旋转本征值 \ensuremath{\pm}2 + Weinberg 软定理 + 弦网低能谱）；\textbf{v1.34 推进（P1-1）}：\textbf{判据链(4)数值化}------C4 螺旋度分解（标量/螺旋\ensuremath{\pm}2/矢量三扇区）、自旋 2 模式有隙存在（\ensuremath{\Delta}\ensuremath{\pm}2\ensuremath{\approx}0.49 尺寸无关）、二阶张量扇区承载 30 个螺旋\ensuremath{\pm}2 本征态（对齐本框架定义 D-2）；\textbf{v1.36 推进（P1-3/P1-4）}：\textbf{3+1D \ensuremath{\lambda}\down{c}\ensuremath{\approx}0.14}（四尺寸 \ensuremath{\chi} 峰锁定，与 2D toric code 构成维数-临界耦合递减谱系）+ \textbf{Wilson RG 路径}（U(1) 高斯固定点缩放律 \ensuremath{\beta}'=\ensuremath{\beta}b\up{2-d} + 2D Ising 块 RG \ensuremath{\beta}<\ensuremath{\beta}\down{c} 流向连续极限------定理 A + 命题 B + 数值验证 C 补第(4)层 RG 流证据，四层闭合）；\textbf{剩余开放}：Z\down{2}\ensuremath{\to}U(1) 连续极限重整化群严格化（世界级，RG 数值路径已演示）、4D 无质量自旋 2 场论涌现（世界级）、\ensuremath{\lambda}\down{c}(\ensuremath{\infty}) 精确值（需 3D 盒子 MPS/DMRG）、麦克斯韦谱严格验证。
\end{enumerate}

\subsection{与主流物理的关系}

\begin{itemize}
\item \textbf{不推翻 GR/QM}：框架定位为二者的"统计力学地基"（GR 是几何描述，框架是统计流体描述；QM 的非交换结构获得本体论动机）；
\item \textbf{与候选路径的关系}：与 Sakharov 诱导引力、Verlinde 熵力、't Hooft 元胞自动机、因果集理论、圈量子引力（面积谱/ln 2）、弦论（KK/全息）均有结构性对应------框架提供统一这些线索的"事件流"语言；
\item \textbf{诚实声明}：所有\mtag{待验证}项不得作为推导依据；框架的数学化已完成主体结构，但若干微观推导仍待严格化。
\end{itemize}



\section*{附录 A　符号与约定}

\begin{itemize}
\item 度规号差 (\ensuremath{-},+,+,+)；自然单位 \ensuremath{\hbar}=c=k\down{B}=1（除非注明）；
\item M\down{Pl} = 1.22\ensuremath{\times}10\up{19} GeV；l\down{P} = 1.62\ensuremath{\times}10\up{-35} m；t\down{P} = 5.4\ensuremath{\times}10\up{-44} s；
\item \ensuremath{\varepsilon}\down{A} = 能标依赖泄漏概率；\ensuremath{\varepsilon}\down{B} = 统计衰减因子；
\item \mtag{公理}= 基本假设；\mtag{假设}= 本框架引入的物理假设；\mtag{推导}= 框架内推导；\mtag{引用}= 已确立结果；\mtag{待验证}= 需验证猜想。
\end{itemize}

\section*{附录 B　来源与 AI 协助声明}

本框架的物理构思源自作者本人的研究探索；构思与推演过程中使用了大型语言模型作为计算与表述辅助工具（2026 年）。\textbf{所有科学主张、推导、近似与数值结果均由作者提出并验证}；语言模型仅作为推演与写作助手，不构成科学权威来源。数值脚本（P0--P3 各报告附件）可复现文中全部数值结果。







\section*{附录 C　版本说明（开发记录）}

\begin{quote}
本附录为 v1.1--v1.36 开发过程的版本变更与沿革完整记录。发表版正文（\S{}1--\S{}10）为当前科学内容的完整表述，不依赖本附录。
\end{quote}

\subsection*{C.1 变更说明（v1.21--v1.36）}

\begin{quote}
\textbf{v1.21 变更}：B3（指数度规 + 完整 PPN + 阴影 EHT 相容）+ B4（方向离散化 + Bekenstein 上界双闭合）+ B5（观测者互信息面积律 + 16 倍压缩）+ B 分支收官（两维卷曲/涨落三层结构/J-K/弦网）并入引力侧论证
\textbf{v1.22 变更}：P0-3（\ensuremath{\kappa}\down{c} 有限尺寸标度）------框架 X 场版本拓扑序\ensuremath{\to}禁闭真相变定位 \textbf{\ensuremath{\kappa}\down{c} \ensuremath{\approx} 0.25\ensuremath{\pm}0.05}（\S{}5.7，补相图完整性），"有限尺寸无相变点"图像获热力学极限定位；蓝图 P0 阶段全部收官
\textbf{v1.23 变更}：P0-3b（环面周期边界精确化）------(1) 环面（周期边界）m=0.5 半极化点 \textbf{\ensuremath{\kappa}\*\ensuremath{\approx}0.20\ensuremath{\pm}0.01}（尺寸无关，消除开边界边界污染）；(2) \textbf{判据修正}：m=0.5 交叉为半极化点非相变点（环面实测 m(h=0.32)\ensuremath{\approx}0.84），真相变 \textbf{\ensuremath{\kappa}\down{c}\ensuremath{\approx}0.25--0.32}；(3) \textbf{文献对照修正}：X 场版本（\ensuremath{\sigma}\up{x} 场）对标 h\down{xc}\ensuremath{\approx}0.32（0.328 为 \ensuremath{\sigma}\up{z} 场版本值，P0-3 曾混用，已修正）（\S{}5.7）
\textbf{v1.24 变更}：钳制物态微观推导（候选 D）------过载区物态方程 P(\ensuremath{\rho})=P\down{P}\ensuremath{\cdot}tanh(\ensuremath{\rho}/\ensuremath{\rho}\down{P}) 从\textbf{碎片-风暴碰撞动力学完整涌现}（\S{}4.6 v1.18 段升级）：(1) \ensuremath{\rho}/\ensuremath{\rho}\down{P}=普朗克填充率（m\down{P}/l\down{P}\up{3}）；(2) tanh=双向净动量流（反射/放行）罗吉斯蒂饱和；(3) P\down{P}=E\down{P}/l\down{P}\up{3} 单元最大压强（公理 1 + \mtag{公理 1 推论：事件流最大频率}+\mtag{假设 H-7：风暴压强有限}）；\textbf{新预言：声速剖面 c\down{s}\up{2}=c\up{2}\ensuremath{\cdot}sech\up{2}(\ensuremath{\rho}/\ensuremath{\rho}\down{P})}（\ensuremath{\rho}\ensuremath{\gg}\ensuremath{\rho}\down{P} 声速消失 \ensuremath{\implies} 浆糊无结构）、能量条件 P \ensuremath{\leq} \ensuremath{\rho}c\up{2} 处处满足（无坍缩）
\textbf{v1.25 变更}：过载区径向结构与辐射传输（后续工作(3)，遗留(2)(3)关闭）------\textbf{等温球（线性区静水平衡）在 r\down{h} 处 M=4M 恒定超重 \ensuremath{\implies} 内部无净引力梯度（视界通量=0 压制）\ensuremath{\implies} 均匀浆糊 \ensuremath{\bar{\rho}}=3c\up{6}/(32\ensuremath{\pi}G\up{3}M\up{2})}（与 \ensuremath{\rho}\down{avg} 精确一致）；辐射传输：宏观黑洞光学深度 \ensuremath{\tau}\ensuremath{\sim}10\up{-40}（\textbf{光学透明}，过载区不吸收/不发射辐射的定量原因）、普朗克质量黑洞 \ensuremath{\tau}\ensuremath{\approx}0.06（临界窗口）（\S{}4.6）
\textbf{v1.26 变更}：LISA EMRI 3PN 检验提案（候选 E）------框架 2PN g\down{rr} 33\% 分叉 \ensuremath{\implies} 进动差 \ensuremath{\delta}=0.30\ensuremath{\cdot}x\down{p} 在 EMRI 中\textbf{相位累积放大}：\ensuremath{\Delta}\ensuremath{\Phi}=N\ensuremath{\cdot}\ensuremath{\Delta}\ensuremath{\omega}\down{GR}\ensuremath{\cdot}\ensuremath{\delta} 达 10\up{2}--10\up{3} rad（6 源/全参数空间可检验，裕度 1128--513240）；3PN g\down{00} 11\% 分叉 \ensuremath{\implies} 啁啾相位差 \ensuremath{\sim}2\ensuremath{\times}10\up{2} rad------EMRI 是框架 2PN/3PN 分叉的\textbf{第一道判决窗口}（与 PTA 互补）（\S{}4.5b + \S{}9 S-E1--S-E4）
\textbf{v1.27 变更}：EMRI 完整波形模板 + Fisher 精度（后续工作(4)）------Peters--Mathews 绝热演化 + quadrupole 波形（进动差 \ensuremath{\delta} 显式注入）+ LISA 噪声 + \textbf{7 参数 Fisher 矩阵}：代表源进动差 \ensuremath{\delta}=2\% 可测 \textbf{24\ensuremath{\sigma}（2yr）/167\ensuremath{\sigma}（4yr）}（\ensuremath{\sigma}\down{\delta{}}=8.3\ensuremath{\times}10\up{-4}/1.2\ensuremath{\times}10\up{-4}），GR 模型被排除 \textbf{\ensuremath{\Delta}lnL=288/13930------单源判决（模板内）}（\S{}4.5b + \S{}9 S-E1）
\textbf{v1.28 变更}：3+1D 弦网严格化（后续工作(2)，\S{}5.8 半严格\ensuremath{\to}部分严格）------(1) \textbf{量子直线-弦网边 Hopf 桥接}（B8 落地，算符级对应）；(2) \textbf{\ensuremath{\lambda}=0 弦网凝聚严格证明}：E0=\ensuremath{-}N\down{c}、基态简并 2\up{n-N\down{c}}（双立方体 512 \checkmark）、激发能隙 \textbf{\ensuremath{\Delta}=2}（涌现粒子）；(3) \textbf{对偶映射严格}（凝聚相\ensuremath{\leftrightarrow}3D Z\down{2} 格点规范禁闭相，Fradkin--Shenker）；(4) \textbf{涌现 U(1) 路径明确化}（对偶严格+连续极限候选）；自旋 2 仍开放（世界级难题）（\S{}5.8 + \S{}10.2 \#13）
\textbf{v1.29 变更}：涌现 U(1) 连续极限与自旋 2（后续工作(2)剩余；参考\textbf{本框架自旋 2 机制}）------(1) \textbf{自旋 2 = 二阶无迹对称张量的旋转本征值 \ensuremath{\pm}2（严格验证}，D\up{2} 谱 {cos2\ensuremath{\theta},cos\ensuremath{\theta},1}）------“自旋 2 是通量涨落通用语言”获表示论支撑；(2) 涌现 U(1) 路径补齐（对偶严格 + Villain 紧致化 + 库仑相麦克斯韦 \ensuremath{\omega}=|k|，文献引用）；(3) 自旋 2 三层论证（框架严格/弦网半严格/\textbf{南大 FQH 投影观测证据}）；(4) \textbf{可检验预言：FQH 引力子模衰减率含 \ensuremath{\varepsilon} 项}（本框架预言）（\S{}5.8 + \S{}9 S-E5 + \S{}10.2 \#13）
\textbf{v1.30 变更}：Z\down{2}\ensuremath{\to}U(1) 严格嵌入定理与 4D 自旋 2 涌现（后续工作(2)最深层；诚实定位：与千禧年同源不声称）------\textbf{定理 A（严格）} Z\down{2} \ensuremath{\subset} U(1) 群嵌入；\textbf{命题 B（解析推导）} cos(\ensuremath{\theta}\down{p}) 线性化 = 麦克斯韦；\textbf{数值验证 C} 库仑相光子 \ensuremath{\omega}=|k|（长波 0.9999）；\textbf{自旋 2 判据链}（Fierz--Pauli 唯一性 + 旋转本征值 \ensuremath{\pm}2 + Weinberg 软定理 + 弦网低能谱）------唯一缺口 = 连续极限重整化群/4D 谱（世界级）（\S{}5.8 + \S{}10.2 \#13）
\textbf{v1.31 变更}：吸积盘辐射传输（后续工作(3)，遗留(3)关闭）------标准薄盘（Shakura--Sunyaev 多温度黑体，T\ensuremath{\sim}1.2\ensuremath{\times}10\up{5} K @ 10\up{6} M\ensuremath{\odot} 0.01 \ensuremath{\dot{M}}\down{Edd}）为盘主体\textbf{与 GR 一致}（外区 1PN 不可区分）；\textbf{过载区全透明}（\ensuremath{\tau}\ensuremath{\sim}10\up{-40}，落入辐射无次级吸收）；\textbf{框架独有差异在视界表层}------火墙非热修正 \ensuremath{\sim}\ensuremath{\sqrt{l\down{P}/R}}（原初黑洞 10\up{-9}--10\up{-4} 窗口）+ 霍金-落入物质纠缠（\S{}4.6 + \S{}9 S4）
\textbf{v1.32 变更}：Kerr EMRI 波形与贝叶斯分析（后续工作(4)）------静态模板推广到 \textbf{Kerr（a=0.9）} + Fisher\ensuremath{\to}贝叶斯（剖面似然+MCMC）：\textbf{\ensuremath{\delta}\down{hat}=0.0200 精确恢复、\ensuremath{\sigma}\down{\delta{}}=9.7\ensuremath{\times}10\up{-4}、\ensuremath{\Delta}lnL=338 \ensuremath{\implies} GR 排除 \ensuremath{\sim}26\ensuremath{\sigma}（Kerr 下稳健）}；GR 对照无假阳性；\ensuremath{\sigma}\down{\delta{}} 与静态 Fisher（8.3\ensuremath{\times}10\up{-4}）同量级（\S{}4.5b + \S{}9 S-E1）
\textbf{v1.33 变更}：非阿贝尔 3+1D 数值实现（行动清单 P0-2；参考\textbf{禁闭机制 + BESIII 锚点}）------\textbf{S\down{3} 面模型（非阿贝尔弦网，6 面 46,656 态）}：\ensuremath{\lambda}=0 弦网凝聚 E0=\ensuremath{-}N\down{c}、简并 \textbf{6\up{n-1}=7776}（群阶替换推广 Z\down{2} 的 2\up{n-1}）；\textbf{非阿贝尔指纹：激发按共轭类分两类}（转置 23,328 / 3-cycle 15,552，阿贝尔无此结构）；能隙随 \ensuremath{\lambda} 单调打开；\textbf{胶球实验锚点[BESIII 2026]：BESIII X(2370)（0\up{-+}、2395 MeV、味单态）确认胶球 = 框架"通量自束缚"的可观测实例}（格点 QCD 2.3--3.0 GeV 对照；红外奴役=弦张力\mtag{本框架机制：禁闭}）（\S{}5.7 + \S{}9 S-E6 + \S{}10.2 \#2）
\end{quote}

\begin{quote}
\textbf{v1.34 变更}：弦网低能谱与自旋 2 判据实证（行动清单 P1-1）------\textbf{C4 链轴四重旋转螺旋度分解}（[R,H]=0、R\up{4}=I 精确）：标量/螺旋\ensuremath{\pm}2/矢量三扇区投影谱（2,048--65,536 态）；\ensuremath{\lambda}=0 精确能级全验证（2\up{n-N\down{c}}\ensuremath{\cdot}C(N\down{c},k) \checkmark）；\textbf{自旋 2 模式有隙存在}（\ensuremath{\Delta}\ensuremath{\pm}2\ensuremath{\approx}0.49，双立方体 0.4919 / 三链 0.4864，尺寸无关）；\textbf{载体实证：二阶张量算符空间 Z\down{fZ}\down{g} 承载 30 个螺旋\ensuremath{\pm}2 本征态}（\ensuremath{\gg}一阶 3）------自旋 2 = 二阶张量扇区（对齐\mtag{本框架定义 D-2}）；诚实结论：Z\down{2} 级无无质量自旋 2（阿贝尔有隙，预期物理）（\S{}5.8 + \S{}10.2 \#13 判据链(4)数值化）
\end{quote}

\begin{quote}
\textbf{v1.35 变更}：EMRI a-\ensuremath{\delta} 联合 MCMC（行动清单 P1-2，v1.32 标注"a 固定"落地；\textbf{发现并修正 v1.32 严重数值缺陷}）------(1) \textbf{欠采样}（v1.32 dt/载波周期=34.6 \ensuremath{\to} \ensuremath{\delta} 似然混叠梳状）：解析旋近 + 奈奎斯特采样；(2) \textbf{轨道相位梳}（a 进 2\ensuremath{\lambda}，2yr 累积 4.3\ensuremath{\times}10\up{5} rad）：\ensuremath{\varphi}0 剖面似然消除；(3) \textbf{Lense--Thirring 进动加入}（LT\ensuremath{\propto}x\up{3/2} vs 1PN\ensuremath{\propto}x，r-标度差破除 \ensuremath{\delta}-a 简并）------\textbf{\ensuremath{\sigma}\down{\delta{}} 修正 600\ensuremath{\times}：9.7\ensuremath{\times}10\up{-4} \ensuremath{\to} 1.6\ensuremath{\times}10\up{-6}（2yr，物理 \ensuremath{\sigma}\down{\delta{}}\ensuremath{\approx}1/(2wGR\down{T}\ensuremath{\cdot}SNR)）}；a-\ensuremath{\varphi}0 全自由下 \textbf{\ensuremath{\Delta}lnL=357/396 \ensuremath{\implies} 27\ensuremath{\sigma}/28\ensuremath{\sigma}}（v1.32 的 4yr"167\ensuremath{\sigma}"=混叠虚高）；\textbf{主简并 \ensuremath{\delta}-\ensuremath{\varphi}0（\ensuremath{\rho}\ensuremath{\approx}\ensuremath{-}0.85）而非 \ensuremath{\delta}-a}；GR 对照无假阳性（\S{}4.5b + \S{}9 S-E1）
\end{quote}

\begin{quote}
\textbf{v1.36 变更}：3+1D \ensuremath{\lambda}\down{c} 相变 + Wilson RG 数值路径（行动清单 P1-3/P1-4）------\textbf{\ensuremath{\lambda}\down{c}\ensuremath{\approx}0.14}（四尺寸 2\ensuremath{\times}1\ensuremath{\times}1/3\ensuremath{\times}1\ensuremath{\times}1/2\ensuremath{\times}2\ensuremath{\times}1/4\ensuremath{\times}1\ensuremath{\times}1 \ensuremath{\chi} 峰全部锁定 \ensuremath{\lambda}\down{\chi{}}=0.140，2,048--2.1\ensuremath{\times}10\up{6} 态；与 2D toric code \ensuremath{\kappa}\down{c}\ensuremath{\approx}0.25--0.32 构成\textbf{维数-临界耦合递减谱系}，对偶 3D Ising \ensuremath{\beta}\down{c}=0.2217 同量级；\mtag{本框架机制：相变=冷凝点}）；\textbf{Wilson RG 路径}：U(1) 高斯固定点缩放律 \ensuremath{\beta}'=\ensuremath{\beta}b\up{2-d}（d=2 临界/d=3 IR 自由）+ 1D Ising decimation 精确（\ensuremath{\beta}*=0）+ \textbf{2D Ising 块 RG：\ensuremath{\beta}<\ensuremath{\beta}\down{c} 流向 0（=U(1) 连续极限）、\ensuremath{\beta}>\ensuremath{\beta}\down{c} 流向 \ensuremath{\infty}（凝聚相）}------Z\down{2}\ensuremath{\to}U(1) 四层证据闭合（定理 A + 命题 B + 数值验证 C + RG 流）；诚实标注：格点块平均次近邻修正（\S{}5.8 + \S{}10.2 \#13）
\end{quote}

\subsection*{C.2 版本沿革}

论文初稿 v1.36 完。版本沿革：v1.1 撤回 \S{}4.6 强场外推；新增 \S{}4.10 packing 定理包。v1.2 定理 6。v1.3 O1--O4。v1.4 N1--N4。v1.5 P1--P4。v1.6 整合版。v1.7 毛球\ensuremath{\implies}\ensuremath{\varepsilon} 下界（H1--H4）。v1.8 火墙厚度判定（F1--F4）。v1.9 通量扩散方程严格解（F1'）。v1.10 3D 球对称精确解（F1''）。v1.11 有限 R 严格解（F1'''）。v1.12 1/r 稳态与弱场 g\down{00} 系数对齐（因子 2 解答）。v1.13 强场分叉裁决（PPN \ensuremath{\beta}：玻尔兹曼 \ensuremath{\beta}=1 通过）。v1.14 通量饱和强场度规数学排除。v1.15 风暴过载观测效应（单参数族 \ensuremath{\alpha} + S1--S5）。v1.16 \ensuremath{\alpha} 微观来源（熵一致性强制 \ensuremath{\alpha}=1）。v1.17 压强临界面积锁定微观机制（M\up{2}/A 普朗克临界）。v1.18 过载区流体模型（普朗克钳制物态）。v1.19 规范侧涌现论证整合（\S{}5.6 量子双/\S{}5.7 哈密顿量/\S{}5.8 弦网）。v1.20 禁闭弦张力数值实现（X 场 \ensuremath{\sigma}\down{E}\ensuremath{\approx}2\ensuremath{\kappa} + Z 场 \ensuremath{\sigma}\down{m}\ensuremath{\approx}2\ensuremath{\lambda}，P0-1 复核）。\textbf{v1.21 引力侧完整化（B 分支整合）}------(1) \S{}4.5b 新增：各向同性指数度规识别（三条推导收敛 + 文献 exponential metric 对接）+ 完整 PPN 参数集（\ensuremath{\gamma}=\ensuremath{\beta}=1 其余 0）+ 2PN/3PN 分叉量化（进动差 \ensuremath{\sim}30x\down{p} 数值）；(2) \S{}4.6 新增：\textbf{各向同性约定下阴影 5.436GM/c\up{2} 与 EHT 相容（+4.55\%），修正 v1.14"强场无度规承诺"为"有与观测相容的几何预言"}；(3) \S{}4.9 新增：方向离散化（普朗克管径 \ensuremath{\implies} k=4 \ensuremath{\implies} Bekenstein 上界双闭合，B4）+ 观测者互信息面积律（体=表面投影 \ensuremath{\implies} I=H(\ensuremath{\Sigma})，16 倍压缩，B5）；(4) \S{}10.2 更新：\#1 方向离散化 \checkmark、\#6 高阶 PPN \checkmark、\#7 两维卷曲 \checkmark（3\ensuremath{-}1=2 定义性）、\#9 涨落 \checkmark（三层结构）、\#11 J/K 推进（量级 O(1)）、\#13 弦网确认（候选路线图）。\textbf{v1.22 P0-3（\ensuremath{\kappa}\down{c} 开边界 FSS 0.23--0.27，\S{}5.7）。v1.23 P0-3b（环面精确化：半极化 \ensuremath{\kappa}*\ensuremath{\approx}0.20\ensuremath{\pm}0.01、真相变 \ensuremath{\kappa}\down{c}\ensuremath{\approx}0.25--0.32、文献对照修正 h\down{xc}\ensuremath{\approx}0.32，\S{}5.7）。v1.24 钳制物态微观推导（P=P\down{P}\ensuremath{\cdot}tanh(\ensuremath{\rho}/\ensuremath{\rho}\down{P}) 从碎片-风暴碰撞涌现 + 声速剖面 c\down{s}\up{2}=c\up{2}\ensuremath{\cdot}sech\up{2}(\ensuremath{\rho}/\ensuremath{\rho}\down{P}) + 能量条件 P\ensuremath{\leq}\ensuremath{\rho}c\up{2}，\S{}4.6）。v1.25 过载区径向结构与辐射传输（等温球超重 4M \ensuremath{\implies} 均匀浆糊 \ensuremath{\bar{\rho}}=3c\up{6}/(32\ensuremath{\pi}G\up{3}M\up{2}) + 光学透明 \ensuremath{\tau}\ensuremath{\sim}10\up{-40} / 普朗克窗口 0.06，\S{}4.6）。v1.26 LISA EMRI 3PN 检验提案（2PN 进动差 \ensuremath{\delta}=0.30\ensuremath{\cdot}x\down{p} 相位累积 10\up{2}--10\up{3} rad、6 源全可检验、S-E1--S-E4，\S{}4.5b/\S{}9）。v1.27 EMRI 完整波形模板 + Fisher 精度（进动差 \ensuremath{\delta} 可测 24\ensuremath{\sigma}/167\ensuremath{\sigma}、GR 排除 \ensuremath{\Delta}lnL 288--13930、单源判决，\S{}4.5b/\S{}9 S-E1）。v1.28 3+1D 弦网严格化（\ensuremath{\lambda}=0 凝聚严格证明 E0=\ensuremath{-}N\down{c}/简并 2\up{n-N\down{c}}/\ensuremath{\Delta}=2 + 对偶严格 + 涌现 U(1) 路径明确化，\S{}5.8/\S{}10.2 \#13）。v1.29 涌现 U(1) 连续极限与自旋 2（自旋 2 = 二阶张量旋转本征值 \ensuremath{\pm}2 严格验证 + Villain 紧致化 + 库仑相麦克斯韦 \ensuremath{\omega}=|k| + FQH \ensuremath{\varepsilon} 项预言 S-E5，\S{}5.8/\S{}9/\S{}10.2 \#13；参考原文 449--454）。v1.30 Z\down{2}\ensuremath{\to}U(1) 严格嵌入定理与 4D 自旋 2（定理 A/B/C：群嵌入严格 + 麦克斯韦线性化严格 + 光子 \ensuremath{\omega}=|k| 数值 + 自旋 2 判据链，\S{}5.8/\S{}10.2 \#13）。v1.31 吸积盘辐射传输（标准薄盘主体与 GR 一致 + 过载区全透明 + 火墙非热修正原初窗口，\S{}4.6/\S{}9 S4）。v1.32 Kerr EMRI 波形与贝叶斯（\ensuremath{\delta}\down{hat}=0.0200 精确恢复、\ensuremath{\sigma}\down{\delta{}}=9.7\ensuremath{\times}10\up{-4}、GR 排除 26\ensuremath{\sigma} Kerr 下稳健、GR 对照无假阳性，\S{}4.5b/\S{}9 S-E1）。v1.33 非阿贝尔 3+1D 数值（S\down{3} 面模型 46,656 态：简并 6\up{n-1}=7776 + 共轭类两类激发指纹 + BESIII 胶球 X(2370) 实验锚点，\S{}5.7/\S{}9 S-E6/\S{}10.2 \#2；参考原文 1250/1595--1596）。v1.34 弦网低能谱与自旋 2 判据实证（C4 螺旋度分解、\ensuremath{\Delta}\ensuremath{\pm}2\ensuremath{\approx}0.49 尺寸无关、二阶张量扇区承载，判据链(4)数值化，\S{}5.8/\S{}10.2 \#13）。v1.35 EMRI a-\ensuremath{\delta} 联合 MCMC（奈奎斯特采样修复欠采样、\ensuremath{\varphi}0 剖面消相位梳、LT 进动；\ensuremath{\sigma}\down{\delta{}} 修正 600\ensuremath{\times} 9.7e-4\ensuremath{\to}1.6e-6、27\ensuremath{\sigma}/28\ensuremath{\sigma}、\ensuremath{\delta}-\ensuremath{\varphi}0 主简并、GR 对照无假阳性，\S{}4.5b/\S{}9 S-E1）。v1.36 3+1D \ensuremath{\lambda}\down{c} 相变 + Wilson RG（\ensuremath{\lambda}\down{c}\ensuremath{\approx}0.14 四尺寸锁定 + 维数谱系；U(1) 高斯固定点 \ensuremath{\beta}'=\ensuremath{\beta}b\up{2-d} + 2D Ising 块 RG 流向连续极限------定理 A/B/C 四层闭合，\S{}5.8/\S{}10.2 \#13）。后续工作：(1) 英文翻译与 arXiv 投稿准备；(2) Teukolsky 波形 / 真实参数估计；(3) Kerr 吸积盘与相对论辐射传输；(4) 实验数据定量对照（EHT/LIGO/LISA/DESI/胶球谱）；(5) \ensuremath{\lambda}\down{c}(\ensuremath{\infty}) 精确值（3D 盒子 MPS/DMRG）与 HOTRG 非微扰。}



\section*{参考文献（初稿，发表版补全至 50+）}
\begin{thebibliography}{30}

\bibitem{ref1} Papapetrou A. (1954) Z. Phys. 139, 518; 141, 293. 
\bibitem{ref2} Yilmaz H. (1958) Phys. Rev. 111, 1417; (1971) Phys. Rev. Lett. 27, 1399; (1973) Found. Phys. 3, 195. 
\bibitem{ref3} Rosen N. (1962) Ann. Phys. 17, 269. 
\bibitem{ref4} Makukov M., Mychelkin E. (2018) Phys. Rev. D 98, 064016. 
\bibitem{ref5} Boonserm P., Ngampitipan T., Simpson A., Visser M. (2021) Phys. Rev. D 103, 024019. 
\bibitem{ref6} Ibison M. (2006) "The Yilmaz cosmology".

\bibitem{ref7} Sakharov A.D. (1968) Dokl. Akad. Nauk 177, 70. 
\bibitem{ref8} Jacobson T. (1995) Phys. Rev. Lett. 75, 1260. 
\bibitem{ref9} Padmanabhan T. (2010) arXiv:1003.5665; (2012) arXiv:1206.4916. 
\bibitem{ref10} Verlinde E. (2011) JHEP 04, 029. 
\bibitem{ref11} Bekenstein J.D. (1973) Phys. Rev. D 7, 2333. 
\bibitem{ref12} Hawking S.W. (1975) Commun. Math. Phys. 43, 199. 
\bibitem{ref13} Unruh W.G. (1976) Phys. Rev. D 14, 870. 
\bibitem{ref14} 't Hooft G. (1993) arXiv:gr-qc/9310026; (1999) arXiv:gr-qc/9903084. 
\bibitem{ref15} Susskind L. (1995) J. Math. Phys. 36, 6377.

\bibitem{ref16} Bombelli L. et al. (1987) Phys. Rev. Lett. 59, 521. 
\bibitem{ref17} Malament D. (1977) J. Math. Phys. 18, 1399.

\bibitem{ref18} Kitaev A.Yu. (2003) Ann. Phys. 303, 2. 
\bibitem{ref19} Wen X.-G. (2003) Phys. Rev. D 68, 065003. 
\bibitem{ref20} Levin M., Wen X.-G. (2005) Rev. Mod. Phys. 77, 871. 
\bibitem{ref21} Fradkin E., Shenker S.H. (1979) Phys. Rev. D 19, 3682. 
\bibitem{ref22} Tupitsyn I.S. et al. (2010) Phys. Rev. B 82, 085114. 
\bibitem{ref23} Tammes P.M.L. (1930) Proc. Sect. Sci. K. Akad. Wet. Amsterdam 33, 31. 
\bibitem{ref24} Wilson K.G. (1974) Phys. Rev. D 10, 2445.

\bibitem{ref25} Babak S. et al. (2017) Phys. Rev. D 95, 103012. 
\bibitem{ref26} Huerta E.A., Gair J.R. (2011) Phys. Rev. D 84, 064023. 
\bibitem{ref27} Barack L., Cutler C. (2004) Phys. Rev. D 69, 082005. 
\bibitem{ref28} Peters P.C., Mathews J. (1963) Phys. Rev. 131, 435. 
\bibitem{ref29} Will C.M. (2014) Living Rev. Relativ. 17, 4. 
\bibitem{ref30} Poisson E., Will C.M. (2014) Gravity (Cambridge).



\end{thebibliography}

\end{document}
